\documentclass[11pt]{amsart}
\pagestyle{plain}

\usepackage{mathrsfs}
\usepackage{amssymb}
\usepackage{titletoc}
\usepackage{amscd}
\usepackage{amsmath, amssymb}
%\usepackage{amsthm}
\usepackage{amsfonts}
\usepackage[colorlinks,linkcolor=blue,citecolor=blue, pdfstartview=FitH]{hyperref}
\usepackage[all,cmtip]{xy}

\usepackage{stmaryrd}
\usepackage{upgreek}
\usepackage{tikz-cd}
\usepackage{bm}


\setlength{\textwidth}{5.6in} \setlength{\oddsidemargin}{0.3in}
\setlength{\evensidemargin}{0.3in}\setlength{\footskip}{0.3in}
\setlength{\headsep}{0.3in}
%%%%%%%%%%%%%%%%%%%%%%%%%%%%%%%%%%%%%%%%%%%%%%%%%%%%%%%%%%%%%%%%%%%%%%%%%%%
\numberwithin{equation}{section}
\renewcommand{\theequation}{\thesection.\arabic{equation}}


\theoremstyle{plain}
\newtheorem{thm}{Theorem}[section]
\newtheorem{theorem}[thm]{Theorem}
\newtheorem{dfn}[thm]{Definition}
\newtheorem{facts}[thm]{Facts}
\newtheorem{lem}[thm]{Lemma}
\newtheorem{corollary}[thm]{Corollary}
\newtheorem{prop}[thm]{Proposition}
\newtheorem{addendum}[thm]{Addendum}
\newtheorem{variant}[thm]{Variant}
%%%%%%%%%%%%%%%%%%%% Text roman %%%%%%%%%%%%%%%%%%%%%%%%%%%%%
\theoremstyle{definition}
\newtheorem{construction}[thm]{Construction}
\newtheorem{notations}[thm]{Notations}
\newtheorem{question}[thm]{Question}
\newtheorem{problem}[thm]{Problem}
\newtheorem{rmk}[thm]{Remark}
\newtheorem{remarks}[thm]{Remarks}
\newtheorem{definition}[thm]{Definition}
\newtheorem{claim}[thm]{Claim}
\newtheorem{assumption}[thm]{Assumption}
\newtheorem{assumptions}[thm]{Assumptions}
\newtheorem{properties}[thm]{Properties}
\newtheorem{exm}[thm]{Example}
\newtheorem{exercise}[thm]{Exercise}
\newtheorem{comments}[thm]{Comments}
\newtheorem{blank}[thm]{}
\newtheorem{observation}[thm]{Observation}
\newtheorem{defn-thm}[thm]{Definition-Theorem}
\newtheorem{defn-pro}[thm]{Definition-Proposition}
\newtheorem{conjecture}[thm]{Conjecture}



%%%%%%%%%%%%%%%%%%%%%%%%%%%%%%%%%%%%%%%%%%%%%%%%%%%%%%%%%%%%%%%%%%%%%%%%%%%
\usepackage{fancyhdr}
%\pagestyle{fancy}
% \pagestyle{myheadings} %\pagestyle{headings}
\renewcommand{\leftmark}{}
\renewcommand{\rightmark}{}

%%%%%%%%%%%%%%%%%%%%%%%%%%%%%%%%%%%%%%%%%%%%%%%%%%%%%%%%%%%%%%%%%%%%%%%%%%%


%Letters 
\newcommand{\A}{{\mathbb A}}
\newcommand{\sA}{{\mathcal A}}
\newcommand{\bA}{{\mathbf C}}
\newcommand{\B}{{\mathbb B}}
\newcommand{\sB}{{\mathcal B}}
\newcommand{\C}{{\mathbb C}}
\newcommand{\sC}{{\mathcal C}}
\newcommand{\D}{{\mathsf D}}
\newcommand{\sD}{{\mathcal D}}
\newcommand{\E}{{\mathbb E}}
\newcommand{\sE}{{\mathcal E}}
\newcommand{\F}{{\mathbb F}}
\newcommand{\bF}{{\textbf F}}
\newcommand{\sF}{{\mathcal F}}
\newcommand{\G}{{\mathbb G}}
\newcommand{\sG}{{\mathcal G}}
\renewcommand{\H}{{\mathbb H}}
\newcommand{\sH}{{\mathcal H}}
\newcommand{\I}{{\mathrm I}}
\newcommand{\sI}{{\mathcal I}}
\newcommand{\J}{{\mathbb J}}
\newcommand{\sJ}{{\mathcal J}}
\renewcommand{\k}{\Bbbk} 
\newcommand{\K}{{\mathbb K}} 
\newcommand{\sK}{{\mathcal K}}
\renewcommand{\L}{{\mathbb L}}
\newcommand{\sL}{{\mathcal L}}
\newcommand{\m}{{\mathfrak m}}
\newcommand{\M}{{\mathbb M}}
\newcommand{\sM}{{\mathcal M}}
\newcommand{\N}{{\mathbb N}}
\newcommand{\sN}{{\mathcal N}}
\renewcommand{\O}{{\mathbb O}}
\newcommand{\sO}{{\mathcal O}}
\renewcommand{\P}{{\mathsf P}}
\newcommand{\rP}{{\mathrm P}}
\newcommand{\sP}{{\mathcal P}}
\newcommand{\Q}{{\mathbb Q}} 
\newcommand{\rQ}{{\mathrm Q}}
\newcommand{\sQ}{{\mathcal Q}}
\newcommand{\fQ}{{\mathfrak Q}}
\newcommand{\R}{{\mathbb R}}
\newcommand{\sR}{{\mathcal R}}
%\newcommand{\rS}{{\mathrm S}}
\newcommand{\sS}{{\mathcal S}}
\newcommand{\fS}{{\mathfrak{S}}}
\newcommand{\T}{{\mathbb T}}
\newcommand{\sT}{{\mathcal T}}
\newcommand{\sfT}{{\mathsf T}}
\newcommand{\U}{{\mathbb U}}
\newcommand{\sU}{{\mathcal U}}
\newcommand{\fU}{{\mathfrak U}}
\newcommand{\V}{{\mathbb V}}
\newcommand{\sV}{{\mathcal V}}
\newcommand{\W}{{\mathbb W}}
\newcommand{\sW}{{\mathcal W}}
\newcommand{\X}{{\mathbb X}}
\newcommand{\Xd}{\check{\X}}
\newcommand{\sX}{{\mathfrak X}}
\newcommand{\Y}{{\mathbb Y}}
\newcommand{\sY}{{\mathcal Y}}
\newcommand{\Z}{{\mathbb Z}}
\newcommand{\sZ}{{\mathcal Z}} 

%algebra
\newcommand{\id}{\mathrm{id}} 
\renewcommand{\d}{{\mathrm{d}}} 
\newcommand{\im}{\mathrm{im}} 
\newcommand{\pr}{\mathrm{pr}} 
\newcommand{\tr}{{\mathrm{tr}}} 
\newcommand{\Tr}{{\mathrm{Tr}}} 
\newcommand{\ord}{{\text{ord}}} 
\newcommand{\xs}{{\ \xrightarrow{\sim}\ }} 
\newcommand{\diag}{\mathrm{diag}}
\newcommand{\rad}{{\mathrm{rad}}} 

\newcommand{\Hom}{\mathrm{Hom}}
\newcommand{\End}{\mathrm{End}}
\newcommand{\Aut}{\mathrm{Aut}}
\newcommand{\Ext}{\mathrm{Ext}}
\newcommand{\gr}{{\mathrm{gr}}}
\newcommand{\Sym}{{\mathrm{Sym}}}
\newcommand{\op}{{\mathrm{op}}}
\newcommand{\cha}{{\mathrm{char}}} 
\newcommand{\Gal}{{\mathrm{Gal}}} 
\newcommand{\Frac}{{\mathrm{Frac}}} 


%cat
\newcommand{\Mod}{\text{-}\mathrm{Mod}}
%\renewcommand{\mod}{\text{-}\on{mod}}
\newcommand{\comod}{\text{-}\mathrm{comod}}
\newcommand{\gmod}{\text{-}\mathrm{gmod}}
\newcommand{\gdom}{\mathrm{gmod}\text{-}}
\newcommand{\DG}{\mathrm{DG}}
\newcommand{\Rep}{\mathrm{Rep}}
\newcommand{\rep}{\mathrm{rep}}
\newcommand{\Vect}{\mathsf{Vect}}
\newcommand{\Coh}{\mathrm{Coh}}
\newcommand{\QCoh}{\mathrm{QCoh}}
\newcommand{\Perv}{\mathrm{Perv}} 


%AG 
\newcommand{\Spec}{\mathrm{Spec}} 
\newcommand{\Proj}{\mathrm{Proj}} 
\newcommand{\supp}{\mathrm{supp}} 
\newcommand{\sing}{\mathrm{sing}} 
\newcommand{\codim}{\mathrm{codim}} 
\newcommand{\rk}{\mathrm{rk}} 
\newcommand{\act}{\mathrm{act}} 
\newcommand{\pt}{\mathrm{pt}} 
\newcommand{\SL}{\mathrm{SL}}

\newcommand{\Div}{\mathrm{Div}} 
\newcommand{\Pic}{\mathrm{Pic}} 
\newcommand{\CH}{\mathrm{CH}} 

\newcommand{\Gm}{{\mathbb{G}_{\mathrm{m}}}}
\newcommand{\Ga}{{\mathbb{G}_{\mathrm{a}}}}

\newcommand{\ad}{\mathrm{ad}}
\newcommand{\Ad}{\mathrm{Ad}}
\newcommand{\Lie}{\mathrm{Lie}} 

\newcommand{\qc}{\mathrm{qc}}
\newcommand{\fg}{\mathrm{fg}}
\newcommand{\pp}{\mathrm{prop}}

%Lie 
\newcommand{\g}{{\mathfrak{g}}}
\renewcommand{\b}{{\mathfrak{b}}}
\renewcommand{\t}{{\mathfrak{t}}}
\newcommand{\n}{{\mathfrak{n}}}
\newcommand{\p}{{\mathfrak{p}}}
\renewcommand{\sl}{{\mathfrak{s}}{\mathfrak{l}}}
\newcommand{\gl}{{\mathfrak{g}}{\mathfrak{l}}}

%sheaf 
\newcommand{\sEnd}{\mathcal{E}nd}
\newcommand{\sHom}{\mathcal{H}om}
\newcommand{\cinf}{\mathcal{C}^{\infty}_X} 
\newcommand{\IC}{\mathrm{IC}} 

%functor
\newcommand{\For}{\mathsf{For}}
\newcommand{\res}{\mathrm{res}}
\newcommand{\ind}{\mathrm{ind}}
\newcommand{\Res}{\mathrm{Res}}
\newcommand{\Ind}{\mathrm{Ind}}

%new
\newcommand{\Br}{{\mathfrak{Br}}}
\newcommand{\Spr}{{\widetilde{\sN}}}
\newcommand{\Gro}{{\widetilde{\g}}}
\newcommand{\Db}{{D^{\mathrm{b}}}}
\newcommand{\semi}{{{\frac{\infty}{2}}}}
\newcommand{\quar}{{{\frac{\infty}{4}}}}
\newcommand{\oblv}{\mathrm{oblv}}
\newcommand{\Av}{\mathrm{Av}}
\newcommand{\av}{\mathrm{av}}
\newcommand{\IW}{\mathcal{IW}}

\newcommand{\simto}{\xrightarrow{\sim}}
\newcommand{\Phis}{\Phi^{\mathrm{s}}}
\newcommand{\ex}{\mathrm{ext}}
\newcommand{\af}{\mathrm{aff}}
%\newcommand{\aff}{\mathrm{aff}}
\newcommand{\Shv}{\mathsf{Shv}}
\newcommand{\Shvc}{\mathsf{Shv}_{\mathrm{c}}}
\newcommand{\Shvfg}{\mathsf{Shv}_{\mathrm{fg}}}
\newcommand{\Ras}{\mathrm{Ras}}
\newcommand{\Fl}{\mathrm{Fl}}
\newcommand{\Gr}{\mathrm{Gr}}
\newcommand{\tFl}{\widetilde{\Fl}}
\newcommand{\bi}{\mathbf{i}}
\newcommand{\bj}{\mathbf{j}}
\newcommand{\bq}{\mathbf{q}}
\newcommand{\bs}{\mathbf{s}}
\newcommand{\Iu}{\mathrm{I}_{\mathrm{u}}}

\newcommand{\comp}{\mathrm{c}}
%\newcommand{\colim}{\mathrm{colim}}
\newcommand{\rS}{\mathfrak{S}}
\newcommand{\uFl}{\underline{\Fl}}
\newcommand{\fd}{\mathrm{fd}}
\newcommand{\ZCen}{\mathsf{Z}}
\newcommand{\Gait}{\mathrm{Ga}}
\newcommand{\Sat}{\mathsf{Sat}}

\DeclareMathOperator*{\colim}{colim}

\def\baselinestretch{1.1}

\begin{document}

\title{Parabolic semiinfinite sheaves on affine flag varieties}

\makeatletter
%\let\uppercasenonmath\@gobble% disables title uppercase
\let\MakeUppercase\relax% disables author uppercase
%\let\scshape\relax% disables section smallcaps
\makeatother

\author{Simon Riche}
\address{Universit\'e Clermont Auvergne, CNRS, LMBP, F-63000 Clermont-Ferrand, France, Institut Universitaire de France.}
\email{simon.riche@uca.fr}

\author{Quan Situ}
\address{Institute of Science and Technology Austria (ISTA), Am Campus 1, 3400 Klosterneuburg, Austria.}
\email{quan.situ@ist.ac.at}

\thanks{
This project has received
funding from the European Research Council (ERC) under the European Union's Horizon 2020
research and innovation programme (grant agreement No.~101002592).
}

\date{}

\begin{abstract}
We study sheaves on affine flag varieties of reductive algebraic groups in parabolic semiinfinite settings, generalizing some constructions and results from \cite{Gaits18,adr} (for the case of the Borel subgroup) and~\cite{DL23} (for characteristic-$0$ coefficients and a particular case of parabolic semiinfinite condition). In particular we define perverse t-structures and Raskin equivalences in this context, and study an analogue of the Gaitsgory sheaf. In future work, these $\infty$-categories will be related to representation-theoretic $\infty$-categories associated with the Langlands dual group.
\end{abstract}

\maketitle

{\small{\setcounter{tocdepth}{1} \tableofcontents}}

%%%%%%%%%%%%%%%%%%%%%%%%%
\section{Introduction} 
%%%%%%%%%%%%%%%%%%%%%%%%%

%--------------------------------------------
\subsection{Semiinfinite sheaves on affine flag varieties}
%--------------------------------------------

This paper is a continuation of a joint project with Pramod Achar and Gurbir Dhillon aiming at developing a framework of \emph{semiinfinite sheaves} on affine flag varieties of connected reductive algebraic groups and exploring some applications to combinatorial, geometric and representation-theoretic questions involving the reductive group under consideration or its Langlands dual.
This study has its roots in predictions in the 1990s (by Lusztig and Finkelberg--Mirkovi{\'c}) that constructible sheaves on \emph{semiinfinite affine flag varieties} should provide geome\-tric models for some categories of representations of the Langlands dual Lie algebra; see~\cite[\S 1.1]{adr} for precise references. 
The concrete realizations of these predictions have been hindered by the difficulty of making sense of such categories of sheaves, but a new construction has been proposed more recently (with different applications in mind) by Gaitsgory~\cite{Gaits18}, in Lurie's framework of (presentable stable) $\infty$-categories, and we expect that this construction will allow us to realize these geometric models and use them more easily.

The first step of this project has been accomplished in~\cite{adr}, where Pramod Achar, Gurbir Dhillon and the first author have explained how to adapt Gaitsgory's construction, initially written for sheaves on affine Grassmannians with coefficients in characteristic $0$, to sheaves on affine flag varieties with coefficients of arbitrary characteristic, and explored some basic properties of these $\infty$-categories. An application in geometry has been explained in~\cite{adr-Drinfeld}, and first applications to representation theory will be developed in~\cite{adr-Gr,adrs}.

%--------------------------------------------
\subsection{Parabolic versions}
%--------------------------------------------

The goal of the present paper is to explain how to construct some variants of the main players of~\cite{adr} associated with parabolic subgroups, and generalize many of the results of~\cite{adr} to that setting.
We also take this opportunity to work out in detail the version of these constructions on general \emph{partial} affine flag varieties. (In the setting of~\cite{adr}, the possibility of this extension was mentioned, but the concrete details were left to the reader except in the important special case of affine Grassmannians.) Parabolic semiinfinite sheaves (on affine Grassmannians, and for coefficients of characteristic $0$) have already been investigated by Gurbir Dhillon and Sergey Lysenko in~\cite{DL23}, which provided some suggestions of constructions and statements; but even in their setting our versions are more general, and many of our proofs are different.

Let us now explain what we mean more precisely. Let $G$ be a connected reductive algebraic group over an algebraically closed field $\mathbb{F}$, and let $T \subset B \subset G$ be a maximal torus and a Borel subgroup. The choice of $B$ determines an Iwahori subgroup $\I$ in the loop group $\sL G$ of $G$, and we can consider the affine flag variety $\Fl=\sL G / \I$.
Let $W_\ex$ be the extended affine Weyl group of $G$, and $S_\af \subset W_\ex$ be the subset of simple reflections determined by $B$. Then for any finitary subset $K \subset S_\af$ we have a parahoric subgroup $\I_K \subset \sL G$ containing $\I$, and the associated partial affine flag variety $\Fl_K=\sL G/\I_K$. Given now a subset $J$ of the set $S \subset S_\af$ of simple reflections in the \emph{finite} Weyl group $W$, we have a standard parabolic subgroup $P_J \subset G$, with a Levi decomposition $P_J = M_J \ltimes U_J$. Any further subset $L \subset J$ is a finitary subset of the set of affine simple reflections of $M_J$, so we have an associated parahoric subgroup $\I_{L \subset J} \subset \sL M_J$, and the semiinfinite parabolic subgroup of $\sL G$ defined by $L \subset J$ is
\[
\I^\semi_{L \subset J} = \I_{L \subset J} \ltimes \sL U_J.
\]

This construction recovers several well-known and well-studied objects:
\begin{itemize}
\item
if $L=J=\varnothing$ one obtains the semiinfinite Iwahori subgroup appearing in particular in~\cite{Gaits18,adr};
\item
if $J=S$ we have $\I_{L \subset S} = \I_L$;
\item
if $L=J$ one obtains the parabolic semiinfinite subgroups appearing in~\cite{DL23};
\item
if $L=\varnothing$ we obtain the ``quarterinfinite'' parahoric subgroups that appear in the discussion in~\cite[\S 8.11]{bez-losev}.
\end{itemize}

%--------------------------------------------
\subsection{Contents}
%--------------------------------------------

Our main objects of study in this paper are the $\infty$-categories
\begin{equation}
\label{eqn:intro-cat}
\Shv \left( \I^\semi_{L \subset J} \backslash \Fl_K \right)
\end{equation}
of \'etale $\I^\semi_{L \subset J}$-equivariant sheaves on $\Fl_K$, with coefficients in a field $\k$ (of characteristic either $0$ or a prime invertible in $\F$), for the various possible choices of $L$, $J$ and $K$, and the natural functors that relate them. In particular:
\begin{enumerate}
\item
We study the relevant combinatorics of $W_\ex$, defining in particular a quarterinfinite length and a quarterinfinite order associated with $J$; see Section~\ref{sec:Wext}.
\item
We define the standard and costandard objects in the $\infty$-categories~\eqref{eqn:intro-cat}, and study their behaviour under various natural functors; see Section~\ref{sec:quar-semi-sheaves}.
\item
We define and study a ``Raskin equivalence'' relating~\eqref{eqn:intro-cat} to the $\infty$-category $\Shv \left( \I_{L} \backslash \Fl_K \right)$; see Section~\ref{sec:Raskin}. (This construction is a direct generalization of constructions in~\cite{Gaits18,adr,DL23}, which have their roots in earlier work of Raskin.)
\item
We define a ``perverse'' t-structure on~\eqref{eqn:intro-cat}, prove t-exactness of various natural functors, prove that standard and costandard objects are perverse in some cases, and describe the t-structure on $\Shv \left( \I \backslash \Fl_K \right)$ obtained by transporting this t-structure via the Raskin equivalence (for $L=\varnothing$) in terms similar to the description of the ``new t-structure'' of Bezrukavnikov--Lin; see Section~\ref{sec:perv}.
\item
We study a natural parabolic analogue of the ``Gaitsgory sheaf'' (or semiinfinite intersection cohomology sheaf) considered in~\cite{Gaits18,adr}; see Section~\ref{sec:Gaitsgory}.
\end{enumerate}
All of these results are natural extensions of statements appearing in~\cite{Gaits18,adr,DL23}, and their proofs are mostly based on similar arguments (with sometimes additional technical difficulties).

The only section containing a result which might be considered more original is Section~\ref{sec:t-exactness}. Here we prove t-exactness of a certain functor
\[
\Shv \left( \I^\semi_{J \subset J} \backslash \Fl_K \right) \to \Shv \left( \I^\semi_{\varnothing \subset \varnothing} \backslash \Fl_K \right)
\]
defined by combining two different choices of Raskin equivalences. The validity of this statement is suggested by expected relations with representation theory described in~\S\ref{ss:into-RT} below; but in fact we are planning to use it to prove some of these conjectural equivalences. This proof requires the consideration of analogues of the $\infty$-categories~\eqref{eqn:intro-cat} where the equivariance condition on the right is replaced by a Whittaker-type condition; for this reason we also study such $\infty$-categories systematically.

%--------------------------------------------
\subsection{Relation to representation theory}
\label{ss:into-RT}
%--------------------------------------------

Our main motivation for the systematic study of the $\infty$-categories $\Shv \left( \I^\semi_{L \subset J} \backslash \Fl_K \right)$ is that we expect them to provide useful ``geometric models'' for several interesting $\infty$-categories of representations of objects related to the Langlands dual reductive group over $\k$. Namely, recall that the \emph{geometric Satake equivalence}~\cite{mv} provides an equivalence of abelian monoidal categories
\begin{equation}
\label{eqn:intro-geom-Satake}
(\Shvfg(\sL^+ G \backslash \Gr)^\heartsuit, \star^{\sL^+ G}) \simeq (\Rep^{\heartsuit,\fd}(G^\vee_\k), \otimes)
\end{equation}
between (finitely generated) $\sL^+G$-equivariant perverse sheaves on the affine Grassmannian $\Gr=\Fl_S = \sL G / \sL^+ G$ (where $\I_S = \sL^+ G$ is the positive loop group) and finite-dimensional representations of the split reductive Langlands dual group $G^\vee_\k$ over $\k$. 

In the particular case $K=S$ and $L=J$ we will construct in~\cite{adrs} an equivalence of abelian monoidal categories
\begin{equation}
\label{eqn:intro-Satake-parabolic}
 \Shv_\omega \left( \I^\semi_{J \subset J} \backslash \Gr \right)^\heartsuit \simeq \Rep^{\heartsuit,\fd}(P_{J,\k}^\vee)
\end{equation}
where $P^\vee_{J,\k} \subset G^\vee_\k$ is the standard parabolic subgroup associated with $J$ (now seen as a subset of the simple reflections of the Weyl group of $G^\vee_\k$, which canonically identifies with $W$), and the symbol ``$\omega$'' indicates a full subcategory introduced formally in the body of the paper. Here the left-hand side is equipped with the fusion product. In the special case $J=S$ this equivalence will coincide with~\eqref{eqn:intro-geom-Satake}. It will generalize an equivalence conjectured in~\cite[\S 1.3.3]{DL23}\footnote{The parabolic subgroup appearing in~\cite[\S 1.3.3]{DL23} is the \emph{negative} parabolic subgroup attached to $J$. The precise relation between the two equivalences should be given by conjugation by the longest element $w_0$ in the Weyl group. For instance, for $J=\varnothing$, in our version the simple object $\overline{\Delta}^\semi_\lambda$ of~\cite[Proposition~6.20]{adr} will correspond to the $1$-dimensional representation of the (positive) Borel subgroup $P^\vee_{\varnothing,\k}$ corresponding to $\lambda$, while the version considered in~\cite{DL23} or~\cite[Proposition~1.5.7]{Gaits18} involves the action of $w_0$.} and the special case $J=\varnothing$ studied (under some assumptions on $\mathrm{char}(\k)$, and without the monoidal structure) in~\cite{adr-Gr}.

Now we turn to the $\infty$-category $\Shv \left( \I^\semi_{J \subset J} \backslash \tFl \right)$, where $\tFl$ is the quotient of $\sL G$ by the pro-unipotent radical $\Iu$ of $\I$. (Sheaves on $\tFl$ are also studied in the body of the paper, in parallel with the $\infty$-categories~\eqref{eqn:intro-cat}.) We assume that $\k$ has positive characteristic $\ell \geq h$ where $h$ is the Coxeter number of $G$, and consider the Lie algebra $\mathfrak{g}^\vee_\k$ of $G^\vee_\k$. 
Following~\cite{adr-cat} we consider the $\infty$-category
\[
(\mathfrak{g}^\vee_\k,P^\vee_{J,\k})\Mod
\]
of $(\mathfrak{g}^\vee_\k,P^\vee_{J,\k})$-modules, and its block $(\mathfrak{g}^\vee_\k,P^\vee_{J,\k})\Mod_0$ containing the trivial module. (Here $(\mathfrak{g}^\vee_\k,P^\vee_{J,\k})\Mod$ is a presentable stable $\infty$-category equipped with a t-structure whose heart is the abelian category of $P^\vee_{J,\k}$-equivariant $U(\mathfrak{g}^\vee_\k)$-modules such that the differential of the action of $P^\vee_{J,\k}$ coincides with the restriction of the action of $\mathfrak{g}^\vee_\k$.) Generalizing the picture sketched in~\cite[\S 1.4]{adr} (corresponding to $J=\varnothing$), we expect that for any $J$ there exists a t-exact equivalence of $\infty$-categories
\begin{equation}\label{equ:parabolic-cat-O}
\Shv \left( \I^\semi_{J \subset J} \backslash \tFl \right) \simeq (\mathfrak{g}^\vee_\k,P^\vee_{J,\k})\Mod_0,
\end{equation}
and that these equivalences satisfy some natural compatibility properties with each other and with some other constructions like~\eqref{eqn:intro-Satake-parabolic}. We expect to prove at least some particular cases of such equivalences in future work. 
An abelian version of the equivalence \eqref{equ:parabolic-cat-O} will be shown in an upcoming work of the second author with Peter Fiebig and Jonathan Gruber \cite{FGS}. 
%of the abelian categories 
%\[
%\Shv_\omega \left( \I^\semi_{J \subset J} \backslash \tFl \right)^\heartsuit 
%\simeq (\mathfrak{g}^\vee_\k,P^\vee_{J,\k})\Mod^{\heartsuit,\fg}_0
%\]

Finally, in~\cite[\S 8.11]{bez-losev}, Bezrukavnikov and Losev conjecture an equivalence relating a (characteristic-$0$ analogue of a) certain derived category of $U(\mathfrak{g}^\vee_\k)$-modules with a fixed nontrivial Frobenius character and a category of $\Iu$-equivariant sheaves on $\sL G / \I_{\varnothing \subset J}^\semi$, mentioning that the definition of the latter category is problematic since ``the space $\sL G / \I_{\varnothing \subset J}^\semi$ behaves pretty badly.'' We expect that a more suitable replacement of this category is the $\infty$-category $\Shv \left( \I^\semi_{\varnothing \subset J} \backslash \tFl \right)$, and are planning to come back to this question in future work.

%--------------------------------------------
\subsection{Acknowledgements}
%--------------------------------------------

As noted above, this paper is part of a joint project with Pramod Achar and Gurbir Dhillon, and although they were not directly involved in the study of the specific aspects considered here their insights on these questions were very helpful. 
We thank Ivan Losev for explaining to us the idea of \cite{bez-losev} during the early stages of this work. 
The second author thanks Xinyu Li for his interest and helpful discussions. 

We acknowledge the use of Copilot and Claude for editing purposes and to check correctness of our proofs. No mathematical content was generated or suggested by AI.

%%%%%%%%%%%%%%%%%%%%%%%%%
\section{Affine Weyl group and affine flag varieties} 
\label{sec:Wext}
%%%%%%%%%%%%%%%%%%%%%%%%%

%---------------------------------------------------------------
\subsection{Roots and weights} 
\label{ss:roots-weights}
%---------------------------------------------------------------

Let $\F$ be an algebraically closed field, whose characteristic will be denoted $p$, and let $G$ be a connected reductive algebraic group over $\F$. 
We fix a Borel subgroup $B\subset G$ and a maximal torus $T\subset B$, and denote by $U\subset B$ the unipotent radical of $B$. 

Let $\X=X^*(T)$ be the character lattice, and $\Xd=X_*(T)$ be the cocharacter lattice. 
Let $\Phi \subset \X$ be the root system of $G$ with respect to $T$, let $\Phi^+\subset \Phi$ be the subset of positive roots determined by $B$ (i.e.~the set of nonzero $T$-weights in $\mathrm{Lie}(B)$), and let $\Phis \subset \Phi^+$ be the associated system of simple roots.
Let $\Xd^{++}\subset \Xd^+ \subset \Xd$ be the subsets of strictly dominant and dominant coweights respectively, and $\check{\Phi} \subset \Xd$ be the set of coroots. 
For $\alpha \in \Phi$, we will denote as usual by $\check{\alpha} \in \check{\Phi}$ the associated coroot. 
We set $2\rho=\sum_{\alpha\in \Phi^+} \alpha \in \X$.
For $\alpha\in \Phi$, we denote by $U_\alpha$ the associated one-parameter subgroup in $G$, and fix an isomorphism $u_\alpha:\Ga\xs U_\alpha$ such that $t u_\alpha(a) t^{-1} = u_\alpha(\alpha(t) a)$ for any $a \in \Ga$ and $t \in T$. 

Let $W$ be the Weyl group of $(G,T)$, endowed with the Coxeter group structure with set of Coxeter generators $S=\{s_\alpha : \alpha\in \Phis\}$, where $s_\beta$ is the reflection associated with $\beta\in \Phi$. 
We will denote by $w_0\in W$ the longest element. 
For every $w \in W$, we fix a lift $\dot{w}\in \mathrm{N}_G(T)$ of $w$. 
%We will denote by $\preceq$ the ``dominance'' order on $\X$ determined by $\Phi^+$, i.e.~the order such that $\lambda \preceq \mu$ if and only if $\mu-\lambda$ is a sum of positive roots.

If $J$ is a subset of $S$, we set $J^\comp:=S\smallsetminus J$. 
%Let $W_J$ be the subgroup of $W$ generated by $J$. 
For such $J$, we will denote by $P_J$ the standard parabolic subgroup of $G$ associated with $J$, with Levi decomposition $P_J=M_J\ltimes U_{J}$ where $M_J$ is the Levi subgroup of $P_J$ containing $T$ and $U_{J}$ is the unipotent radical of $P_J$. 
Let
\[
\Phis_J=\{\alpha\in \Phis \mid s_\alpha\in J\}, \quad \Phi^+_J=\Phi^+ \cap \Z\Phis_J, \quad \Phi_J=\Phi \cap \Z\Phis_J
\]
be the subsets of simple roots, positive roots, and roots of the group $M_J$, with respect to the maximal torus $T \subset M_J$ and the Borel subgroup $B_J:=B \cap M_J$. 
We set $2\rho_J=\sum_{\alpha\in \Phi^+_J} \alpha$. 
We will consider the following subsets of $\Xd$ attached to $J$:
\begin{align*}
\Xd_J^0 &:=\{\lambda\in \Xd \mid \forall \alpha\in \Phis_J, \, \langle \lambda, \alpha\rangle = 0\}, \\
\Xd^+_J &:=\{\lambda\in \Xd \mid \forall \alpha\in \Phis_J, \, \langle \lambda, \alpha\rangle \geq 0\}, \\
\Xd^{++}_J &:=\{\lambda\in \Xd \mid \forall \alpha\in \Phis_J, \, \langle \lambda, \alpha\rangle > 0\}. 
\end{align*}
%We call any $\lambda\in \Xd^+_J$, resp. $\lambda\in \Xd^{++}_J$, a \textit{$J$-dominant} weight, resp. a \textit{strictly $J$-dominant} weight. 
%A weight $\lambda$ is called \textit{$J$-dominant enough}, if $\langle \lambda, \alpha\rangle>0$ is large enough for any $\alpha\in \Phis_J$. 
(To compare with the notation of~\cite{DL23}, our $\Xd_J^0$ corresponds to their $\Lambda_{M,ab}$ where $M=M_J$, and the subset $\Xd_J^0 \cap \Xd^+=\Xd_J^0 \cap \Xd_{J^\comp}^+$ that will appear frequently below corresponds to their $\Lambda_{M,ab}^+$.) We will also denote by $W_J \subset W$ the parabolic subgroup generated by $J$, i.e.~the Weyl group of $(M_J,T)$.

%---------------------------------------------------------------
\subsection{Affine Weyl group} 
%---------------------------------------------------------------

%===============================
\subsubsection{Affine roots and affine Weyl group} 
%===============================

For the statements in this subsection, we refer to~\cite[\S\S 1--2]{macdonald}.
We consider the affine real roots 
\[
\Phi_\af=\Phi+\Z\hbar= \{\alpha+n\hbar : \alpha\in \Phi,n\in \Z\} \subset \X \oplus \Z \hbar,
\] 
together with the collection of positive roots
\[
\Phi^+_\af=(\Phi^++\Z_{\geq 0}\hbar) \sqcup (-\Phi^++\Z_{> 0}\hbar).
\] 
Following usual conventions, for $\alpha \in \Phi_\af$, we will write $\alpha >0$ for the condition that $\alpha \in \Phi^+_\af$, and $\alpha<0$ for the condition that $\alpha \in -\Phi^+_\af$.

We define the extended affine Weyl group $W_\ex$ and the affine Weyl group $W_\af$ as
\[
W_\ex=W\ltimes \Xd \quad \text{and} \quad W_\af= W\ltimes \Z\check{\Phi}. 
\]
For $\lambda\in \Xd$ we will denote by $t_\lambda := (1 \ltimes \lambda) \in W_\ex$ the translation element corresponding to $\lambda$. 
For $\alpha \in \Phi$ and $n \in \Z$,
%$\alpha+n\hbar\in \Phi_\af$, 
we set $s_{\alpha+n\hbar} = t_{n\check{\alpha}} s_\alpha =s_\alpha t_{-n\check{\alpha}} \in W_\af$. 
The group $W_\af$ admits a structure of Coxeter group, with set of Coxeter generators $S_\af=S \cup \{s_{-\gamma+\hbar}\}_\gamma$ where $\gamma$ runs over the maximal roots in $\Phi^+$. 
Let $\ell: W_\af\rightarrow \Z_{\geq 0}$ be the length function associated with this Coxeter structure. 
It is well known that we have
\begin{equation}
\label{eq:length}
%\ell(t_\lambda w) = \sum_{\substack{\alpha \in \Phi^+ \\ w(\alpha) \in \Phi^+}} |\langle \lambda, w(\alpha) \rangle | + \sum_{\substack{\alpha \in \Phi^+ \\ w(\alpha) \in -\Phi^+}} |1- \langle \lambda, w(\alpha) \rangle |,
\ell(t_\lambda w) = \sum_{\substack{\alpha \in \Phi^+ \\ w^{-1}(\alpha) >0}} |\langle \lambda, \alpha \rangle | + \sum_{\substack{\alpha \in \Phi^+ \\ w^{-1}(\alpha) <0}} |1 + \langle \lambda,\alpha\rangle |
%\qquad\text{for $w \in W_\fin$, $\lambda \in \Z\fR^\vee$.}
\end{equation}
for $w \in W$ and $\lambda \in \Z\check{\Phi}$. 
This formula makes sense more generally when $\lambda \in \Xd$, which allows to extend the function $\ell$ to $W_\ex$. 

%We extend it to a function $\ell: W_\ex\rightarrow \Z_{\geq 0}$ using \eqref{equ length}. 
Setting $\Omega := \{w \in W_\ex \mid \ell(w)=0\}$, it is a standard fact that $\Omega$ is a finitely generated abelian group, that conjugation by $\Omega$ on $W_\ex$ preserves $S_\af$, that multiplication induces a group isomorphism
\[
\Omega \ltimes W_\af \simto W_\ex,
\]
and finally that $\ell(w\omega) = \ell(w) = \ell(\omega w)$ for any $\omega \in \Omega$ and $w \in W_\ex$.

\begin{rmk}
\label{rmk:Coxeter-structures}
The Coxeter structure on $W_\af$ considered above is the same as the one used in~\cite{adr}. It is \emph{not} the same as the one used in~\cite{AR18}; more precisely these two structures are twists of each other by the automorphism of $W_\af$ which is the identity on $W$ and sends each $t_\lambda$ to $t_{-\lambda}$ ($\lambda \in \Z\check{\Phi}$), or alternatively are twists of each other by the automorphism given by conjugation by $w_0$.
\end{rmk}

Let $\leq$ be the Bruhat order on $W_\ex$, which is generated by 
\[
x< xt \quad \text{for any reflection $t\in W_\af$ such that $\ell(x)<\ell(xt)$},
\]
where we call $t\in W_\af$ a reflection if it is contained in
\[
\bigcup_{y\in W_\ex} yS_\af y^{-1} = \{s_\beta : \beta \in \Phi_\af\}.
\] 

The group $W_\ex$ acts on $\Phi_\af$ by the formula
\begin{equation}
\label{eq:act-aff}
(t_\lambda w)(\alpha+n\hbar)=w(\alpha)+(n+\langle \lambda,w(\alpha)\rangle)\hbar 
\end{equation}
for $\lambda\in \Xd$, $w\in W$, $\alpha \in \Phi$ and $n \in \Z$.
In terms of this action, we have
\[
\ell(w) = \#\{\alpha \in \Phi_\af^+ \mid w^{-1}(\alpha)<0\}.
\]
For any $x\in W_\ex$ and $\alpha\in \Phi^+_\af$ we also have
\begin{equation}
\label{eq:Bruhat}
x < xs_\alpha \quad \text{if and only if}\quad x(\alpha)\in \Phi^+_\af. 
\end{equation}

%Moreover, $W$ is the parabolic subgroup in $W_\af$ generated by the subset $S := S_\af \cap W$ of $S_\af$ and, setting $\Omega := \{w \in W_\ex \mid \ell(w)=0\}$, $\Omega$ is an abelian group whose action by conjugation preserves $S_\af$, and multiplication induces a group isomorphism
%\[
%\Omega \ltimes W_\af \simto W_\ex.
%\]

%===============================
\subsubsection{Quarterinfinite length and order} 
\label{sss:quar-length}
%===============================

Let $J\subset S$.
%, and recall $J^c=S\smallsetminus J$. 
We define the \emph{quarterinfinite length}\footnote{The term ``quarterinfinite'' used here and below is taken from the discussion in~\cite[\S 8.11]{bez-losev}, which was one of the starting points for our investigations (see also~\S\ref{ss:into-RT}).} associated with $J$ on $W_\ex$ by 
\begin{multline*}
%\label{equ quar-length}
\ell^J(t_\lambda w) = \langle \lambda, 2\rho-2\rho_J \rangle 
+ \#\{ \alpha \in \Phi^+\smallsetminus \Phi^+_J \mid w^{-1}(\alpha) <0 \} \\ 
+ \sum_{\substack{\alpha \in \Phi_J^+ \\ w^{-1}(\alpha) >0}} | \langle \lambda, \alpha \rangle | 
+ \sum_{\substack{\alpha \in \Phi_J^+ \\ w^{-1}(\alpha) <0}} | 1+ \langle \lambda, \alpha \rangle |
\end{multline*}
for $\lambda \in \Xd$ and $w \in W$.
Note that for a fixed $x \in W_\ex$, for $N \gg 0$ we have
\[
\ell^J(x)=\ell(t_{\theta} x)-\ell(t_\theta)
\]
for any $\theta \in \Xd_J^0$ such that $\langle \theta,\alpha \rangle \geq N$ for all $\alpha \in \Phis_{J^\comp}$. 
%which does not depend on the choice of $\theta$. 
%More explicitly, for $w \in W$ and $\lambda \in \Xd$ we have 
%\begin{equation}\label{equ quar-length}
%\begin{aligned}
%\ell^J(t_\lambda w) = \langle \lambda, & 2\rho-2\rho_J \rangle 
%+ |w(\Phi^+) \cap -(\Phi^+\backslash \Phi^+_J)| \\ 
%&+ \sum_{\substack{\alpha \in \Phi^+ \\ w(\alpha) \in \Phi^+_J}} | \langle \lambda, w(\alpha) \rangle | 
%+ \sum_{\substack{\alpha \in \Phi^+ \\ w(\alpha) \in -\Phi^+_J}} | 1- \langle \lambda, w(\alpha) \rangle |.
%\end{aligned}
%\end{equation}
If $J=S$ then $\ell^S$ is the length function $\ell$ from~\eqref{eq:length}; if $J=\varnothing$ then $\ell^{\varnothing}$ is the semiinfinite length considered e.g.~in~\cite[Remark~3.5]{adr}. 

\begin{rmk}
\phantomsection
\label{rmk:lJ}
\begin{enumerate}
\item
\label{it:lJ-translation}
It is clear from the definition that for $\theta \in \Xd_J^0$ and $x \in W_\ex$ we have $\ell^J(t_\theta x) = \ell^J(x) + \langle \theta, 2\rho \rangle$.
\item
\label{it:lJ-WJext}
Consider the extended affine Weyl group $W_{J,\ex}$ of the reductive group $M_J$ (with maximal torus $T$ and Borel subgroup $B_J$), and denote its length function by $\ell_J$. Then $W_{J,\ex}$ identifies with a subgroup of $W_\ex$, and any $w \in W_\ex$ can be written uniquely as $w=yx$ with $y \in W_{J,\ex}$ and $x \in W$ the unique minimal-length representative of its coset $W_J x$. Writing $y=t_\lambda z$ with $\lambda \in \Xd$ and $z \in W_J$, using the facts that $z^{-1}$ permutes $\Phi^+ \smallsetminus \Phi_J^+$ and that $x^{-1}(\Phi_J^+) \subset \Phi^+$ (which is classical, and follows from~\eqref{eq:Bruhat}) we see that $\ell^J(w) = \langle \lambda, 2\rho-2\rho_J \rangle + \ell(x) + \ell_J(y)$.
\end{enumerate}
\end{rmk}

We also consider the \emph{quarterinfinite order} $\preceq^J$ on $W_\ex$ defined as follows. Consider the subset $W_\ex^{+,J} \subset W_\ex$ consisting of elements $x \in W_\ex$ which satisfy $\ell(t_\mu x)=\ell(t_\mu) + \ell(x)$ for any $\mu \in \Xd_J^0 \cap \Xd^+_{J^\comp}$. In particular, this subset contains all the elements of the form $t_\lambda w$ with $w \in W$ and $\lambda \in \Xd$ which satisfies $\langle \lambda, \beta \rangle \geq 0$ for all $\beta \in \Phi^+ \smallsetminus \Phi_J^+$, and is stable under left multiplication by any $t_\theta$ with $\theta \in \Xd_J^0 \cap \Xd^+_{J^\comp}$. As in~\cite[\S 3.1]{adr} one checks that for $x,y \in W_\ex$ the following conditions are equivalent:
\begin{enumerate}
\item
there exists $\mu \in \Xd^0_J$ such that $t_\mu x$ and $t_\mu y$ belong to $W_\ex^{+,J}$ and $t_\mu x \leq t_\mu y$;
\item
for any $\mu \in \Xd^0_J$ such that $t_\mu x$ and $t_\mu y$ belong to $W_\ex^{+,J}$, we have $t_\mu x \leq t_\mu y$.
\end{enumerate}
%subgroup $W_{J^\comp} \subset W$ generated by $J^\comp$, and the subset $W_\ex^{J^\comp}$ of elements $w$ which are of minimal length in $wW_{J^\comp}$. Note that for any fixed $x \in W_\ex$, for $N \gg 0$ we have $t_\theta x \in W_\ex^{J^\comp}$ for any $\theta \in \Xd_J$ such that $\langle \theta,\beta \rangle \geq N$ for all $\alpha \in \Phis_{J^\comp}$.
When these conditions are satisfied, we write $x \preceq^J y$.
%define the \textit{quarter-infinite order associated to $J$} on $W_\ex$ by 
%\[
%x \preceq^{J} y \quad \text{if and only if}\quad t_\theta x \leq t_\theta y,
%\]
%for $\theta\in \Xd_J$ that is $J^c$-dominant enough, which does not depend on the choice of $\theta$. 

This order is ``$\Xd^0_J$-periodic'' in the sense that for $x,y \in W_\ex$ and $\lambda \in \Xd^0_J$ we have $x \preceq^J y$ if and only if $t_\lambda x \preceq^J t_\lambda y$.
When $J=S$ we have $W_\ex^{+,S} = W_\ex$, and $\preceq^S$ is the Bruhat order. When $J=\varnothing$, $W_\ex^{+,\varnothing}$ is the subset denoted $W_\ex^+$ in~\cite[\S 3.1]{adr}, and $\preceq^{\varnothing}$ is the ``semiinfinite order,'' which is denoted $\leq^\semi$ in~\cite{adr}.
%see e.g. \cite{adr} (where it is denoted by $\leq^\semi$). 
The order $\preceq^{J}$ can also be characterized using the quarterinfinite length: it is the order generated by 
\[
x \preceq^{J} xt \quad \text{for any reflection $t\in W_\af$ such that}\ \ell^J(x)<\ell^J(xt). 
\]
In particular, $\ell^J$ is increasing with respect to $\preceq^J$ (and the standard order on $\Z$).

We set 
\[
{}^J\Phi^+_\af:= (\Phi_J^++\Z_{\geq 0}\hbar) \sqcup (-\Phi_J^++\Z_{> 0}\hbar) \sqcup ((\Phi^+\smallsetminus \Phi_J^+) +\Z \hbar). 
\]
Note that for $\alpha \in \Phi_\af$, we have $\alpha \in {}^J\Phi^+_\af$ if and only if there exists $\theta \in \Xd_J^0$ such that $t_{\theta+\theta'} (\alpha) \in \Phi_\af^+$ for any $\theta' \in \Xd_J^0 \cap \Xd_{J^\comp}^+$.
The following lemma is an analogue of~\eqref{eq:Bruhat}.
%, which can be easily deduced.
%using \eqref{equ quar-length}. 

\begin{lem}
\label{lem:quar-order}
For any $x\in W_\ex$ and $\alpha\in \Phi^+_\af$, we have $x\prec^{J} xs_\alpha$ if and only if $x(\alpha)\in {}^J\Phi_\af^+$. 
\end{lem}

\begin{proof}
Assume that $x\prec^{J} xs_\alpha$, and choose $\theta \in \Xd_J^0$ such that $t_\theta x$ and $t_\theta x s_\alpha$ belong to $W_\ex^{+,J}$. Then for any $\theta' \in \Xd_J^0 \cap \Xd_{J^\comp}^+$ we have  $t_{\theta+\theta'} x \leq t_{\theta + \theta'} x s_\alpha$, so $t_{\theta+\theta'} x(\alpha) \in \Phi_\af^+$ by~\eqref{eq:Bruhat}. This forces $x(\alpha)\in {}^J\Phi_\af^+$.

Conversely, assume that $x(\alpha)\in {}^J\Phi_\af^+$, and choose $\theta \in \Xd^0_J$ such that $t_\theta x$ and $t_\theta x s_\alpha$ belong to $W_\ex^{+,J}$ and $t_{\theta} x(\alpha) \in \Phi_\af^+$. (Such a $\theta$ indeed exists.) Then we have $t_{\theta} x < t_{\theta} x s_\alpha$ by~\eqref{eq:Bruhat}, hence $x\prec^{J} xs_\alpha$.
\end{proof}

%===============================
\subsubsection{Minimal coset representatives} 
%===============================

For any $J\subset S_\af$, we denote by $W_J$ the subgroup of $W_\ex$ generated by $J$. 
Recall that $J$ is called \emph{finitary} if $W_J$ is finite; in this case
%If $J\subsetneq S_\af$ such that $W_J$ is finite, 
we will denote by $w_J$ the longest element in $W_J$. (Note in particular that any subset of $S$ is finitary, and that we have $w_S=w_0$.)

For $J\subset S$ and $K\subset S_\af$ finitary we set 
\begin{gather*}
{}^JW_\ex:=\{x\in W_\ex \mid 
\forall w\in W_J, \, \ell(x)\leq \ell(wx)\}, 
\\
W_\ex^{(K,J)}:=\{x\in W_\ex \mid 
\forall w\in W_K, \, \ell^{J}(x)\leq \ell^{J}(xw)\}.
\end{gather*}
In particular, ${}^J W_\ex$ is the set of minimal-length representatives for the cosets in $W_J \backslash W_\ex$ (with respect to $\ell$), and for $x \in W_\ex$ we have $x \in {}^J W_\ex$ if and only if $sx>x$ for any $s \in J$. When $J=S$, we will write $W_\ex^K=W_\ex^{(K,S)}$; this set is the set of minimal-length representatives for the cosets in $W_\ex / W_K$. (We will use similar conventions for the variants of this notation introduced below.)

%For $J,K\subset S_\af$, we denote by ${}^JW_\ex^{K}$ the collection of minimal representatives in $(W_J,W_K)$-cosets of $W_\ex$, namely
%\[
%{}^JW_\ex^{K}=\{x\in W_\ex \mid \ell(x)\leq\ell(wxw'),\ \forall w\in W_J, w'\in W_K\}. 
%\]
%We abbreviate ${}^JW_\ex={}^JW_\ex^{\emptyset}$ and $W_\ex^{K}={}^{\emptyset}W_\ex^K$. 

\begin{lem}
\label{lem:min-length}
Let $J\subset S$ and $K\subset S_\af$ be a finitary subset. Let $x\in W_\ex$.
\begin{enumerate}
\item 
\label{it:min-length-0}
%Both ${}^JW_\ex$ and $W_\ex^{(K,J)}$ are stable under left multiplication by elements in $\Xd_J^0$.
If $\theta \in \Xd_J^0$, we have
\[
x \in {}^J W_\ex \Leftrightarrow t_\theta x \in {}^J W_\ex, \qquad
x \in W_\ex^{(K,J)} \Leftrightarrow t_\theta x \in W_\ex^{(K,J)}.
\]
\item 
\label{it:min-length-1}
We have $x\in {}^JW_\ex$ if and only if $x$ is minimal in $W_Jx$ for the quarterinfinite length $\ell^J$. 
In this case, for any $w\in W_J$ we have $\ell^{J}(wx)=\ell^{J}(x)+\ell(w)$. 
\item 
\label{it:min-length-2}
If $x\in W_\ex^{(K,J)}$, for any $w\in W_K$ we have $\ell^J(xw)=\ell^J(x)+\ell(w)$. 
\end{enumerate}
\end{lem}

\begin{proof}
\eqref{it:min-length-0}
Let $\theta \in \Xd_J^0$.
By~\eqref{eq:Bruhat}, we have $x\in {}^JW_\ex$ if and only if $x^{-1}(\alpha)\in \Phi^+_\af$ for any $\alpha\in \Phis_J$. 
Since $t_{-\theta}$ fixes $\Phis_J$ (under the action \eqref{eq:act-aff}), 
%multiplication on the left by $t_\theta$ stabilizes ${}^J W_\ex$. 
this is equivalent to the condition that $x^{-1}t_{-\theta}(\alpha)\in \Phi^+_\af$ for any $\alpha\in \Phis_J$, i.e.~$t_\theta x \in {}^J W_\ex$.

Now, choose $\theta' \in \Xd_J^0$ such that for any $w \in W_K$ we have
\[
\ell^J(xw) = \ell(t_{\theta+\theta'} xw) - \ell(t_{\theta+\theta'}), \quad \ell^J( t_{\theta}xw) = \ell(t_{\theta'} t_{\theta}xw) - \ell (t_{\theta'}).
\]
%Here $w$ commutes with $t_\theta$
We have $x \in W_\ex^{(K,J)}$ if and only if $\ell(t_{\theta+\theta'} xw) \geq \ell(t_{\theta+\theta'} x)$ for any $w \in W_K$, which is equivalent to $t_\theta x \in W_\ex^{(K,J)}$.

\eqref{it:min-length-1}
Let $x \in W_\ex$, and fix $\theta\in \Xd_J^0$ 
such that for any $w \in W_J$ we have
\[
\ell^J(wx) = \ell(t_\theta wx) - \ell(t_\theta).
\]
If $x \in {}^J W_\ex$ then $t_\theta x \in {}^J W_\ex$ by~\eqref{it:min-length-0}; since each $w$ as above commutes with $t_\theta$ we deduce that $\ell(t_\theta wx)=\ell(w) + \ell(t_\theta x)$, hence $\ell^J(wx)=\ell^J(x)+\ell(w)$. In particular, $x$ is minimal in $W_J x$ for $\ell^J$. Conversely, if $\ell^J(wx) \geq \ell^J(x)$ for any $w \in W_J$, then $t_\theta x \in {}^J W_\ex$, hence $x \in {}^J W_\ex$ by~\eqref{it:min-length-0}.

%Fix now $x \in W_\ex$, and $\mu \in \Xd^0_J$ such that each $t_\mu wx$ with $w \in W_J$ belongs to $W_\ex^{+,J}$. If $x \in {}^J W_\ex$ then $t_\mu x \in {}^J W_\ex$, so $t_\mu x \leq t_\mu wx$ for any $w \in W_J$ (because $w$ commutes with $t_\mu$), so $x \preceq^J wx$ for any $w \in W_J$, and finally $\ell^J(x) \leq \ell^J(wx)$ for any $w \in W_J$. Conversely, if $x$ is minimal in $W_J x$ for $\ell^J$, then for any $s \in J$ we have $x \preceq^J sx$, hence $t_\mu x \leq s t_\mu x$, so that $t_\mu x \in {}^J W_\ex$, and finally $x \in {}^J W_\ex$.

%If these conditions are satisfied, one can choose 
%%$N$ such that for any 
%$\theta\in \Xd_J^0$ 
%%which satisfies $\langle \theta, \beta \rangle \geq N$ for any $\beta \in \Phis_{J^\comp}$, 
%such that for any $w \in W_J$ we have
%\[
%\ell^J(wx) = \ell(t_\theta wx) - \ell(t_\theta).
%% = \ell(w t_\theta x) - \ell(t_\theta).
%\]
%%If $x \in {}^J W_\ex$ then $t_\theta x \in {}^J W_\ex$ and therefore $\ell^J(wx) = \ell(w) + \ell^J(x)$ for any $w \in W_J$. Conversely, if $x$ is minimal in $W_J x$ for $\ell^J$, then for $\theta$ as above we have $t_\theta x \in {}^J W_\ex$, and then $x=t_{-\theta} t_\theta x \in {}^J W_\ex$.
%Here we have $t_\theta wx = wt_\theta x$, and $t_\theta x \in {}^J W_\ex$, so $\ell(t_\theta wx)=\ell(w) + \ell(t_\theta x)$, and finally $\ell^{J}(wx)=\ell^{J}(x)+\ell(w)$.

%For $x\in {}^JW_\ex$, we let $\theta$ be $J^c$-dominant enough such that $\ell^J(t_\theta wx)=\ell(t_\theta wx)$, for any $w\in W_J$.  Then we have $\ell^J(wx)=\ell(t_\theta wx)-\ell(t_\theta)=\ell(wt_\theta x)-\ell(t_\theta)=\ell(w)+\ell(t_\theta x)-\ell(t_\theta)=\ell(w)+\ell^J(x)$. 

\eqref{it:min-length-2}
%For $N \gg 0$, for any 
Choose $\theta \in \Xd_J^0$ such that 
%$\langle \theta, \beta \rangle \geq N$  for any $\beta \in \Phis_{J^\comp}$ and 
for any $w \in W_K$ we have
\[
\ell^J(xw) = \ell(t_\theta xw)-\ell(t_\theta).
\]
We have $t_\theta x\in W_\ex^{K}$, and then $\ell(t_\theta xw) = \ell(t_\theta x)+\ell(w)$ for any $w \in W_K$, which implies the desired formula.
%Let $\theta\in \Xd_J$ be $J^c$-dominant enough such that the order $\preceq^{J}$ coincides with $\leq$ on the (finite) set $\{t_\theta xw\}_{w\in W_K}$. Then $t_\theta x\in W_\ex^{K}$, and we have $\ell^J(xw)=\ell(t_\theta xw)-\ell(t_\theta)=\ell(t_\theta x)+\ell(w)-\ell(t_\theta)=\ell^J(x)+\ell(w)$. 
\end{proof}

It follows in particular from Lemma~\ref{lem:min-length}\eqref{it:min-length-2} that $W_\ex^{(K,J)}$ is a subset of representatives for the cosets in $W_\ex/W_K$. For $L \subset S$ we set 
\[
{}^{L}W_\ex^{(K,J)}:={}^{L}W_\ex\cap W_\ex^{(K,{J})}.
\]
%\left\{x\in W_\ex \left|
%\begin{gathered}
%xw_K\in {}^JW_\ex\\ 
%\ell^J(x)\leq \ell^J(xw),\ \forall w\in W_K
%\end{gathered}
%\right. \right\}. 
%If $J=S$ (resp. $J'=\varnothing$) we will omit $J$ (resp. $J'$) in the notation; e.g. ${}^{J'} W_\ex^{K}={}^{J'}W_\ex^{(K,S)}$, 
%$W_\ex^{(K,J),\mathrm{w}}={}^\varnothing W_\ex^{(K,J),\mathrm{w}}$, and 
%$W_\ex^{K,\mathrm{w}}=W_\ex^{(K,S),\mathrm{w}}$. In particular, 
When $J=S$, ${}^{L}W_\ex^{K}$ is the set of minimal-length representatives for the cosets in $W_{L} \backslash W_\ex / W_K$.
% is the set of minimal (under Bruhat order) representatives in $(W_J,W_K)$-cosets of $W_\ex$, and $W_\ex^{(K,J),\mathrm{w}}$ is the set of minimal (under order $\preceq^J$) representatives in $W_J$
For any subset $D \subset S_\af$ and $x\in W_\ex$, we set ${}^x D=xDx^{-1}$, $D^x=x^{-1}Dx$. For $D_2 \subset D_1$ finitary, we also set $W_{D_1}^{D_2} := W_{D_1} \cap W_\ex^{D_2}$, ${}^{D_2} W_{D_1} := W_{D_1} \cap {}^{D_2} W_\ex$.

\begin{lem}
\label{lem:bi-min}
Let $L \subset J\subset S$ and $K\subset S_\af$ be finitary. 
\begin{enumerate}
\item 
\label{it:bi-min-1}
%The compositions 
%\[
%{}^JW_\ex \hookrightarrow W_\ex \twoheadrightarrow W_J\backslash W_\ex , \quad 
%W_\ex^{(K,J)} \hookrightarrow W_\ex \twoheadrightarrow W_\ex/W_K
%\]
%and 
%\[
%{}^JW_\ex^{(K,J)} \hookrightarrow W_\ex \twoheadrightarrow W_J\backslash W_\ex /W_K
%\]
%are bijections. 
%For any subset $J' \subset J$, 
The composition
\[
{}^{L} W_\ex^{(K,J)} \hookrightarrow W_\ex \twoheadrightarrow W_{L}\backslash W_\ex /W_K
\]
is a bijection. 
\item 
\label{it:bi-min-2}
Let $x\in {}^{L}W_\ex^{(K,J)}$. 
We have $W_{L}\cap xW_Kx^{-1} = W_{L \cap ({}^xK)}$. Moreover, for any $w \in W_{L}^{L\cap ({}^xK)}$ we have $wx\in W_\ex^{(K,J)}$, and all the representatives in $W_\ex^{(K,J)}$ of the cosets of the form $wxW_K$ with $w \in W_L$ are of this form.
%We have $W_{L}\cap xW_Kx^{-1} = W_{L \cap ({}^xK)}$, and 
%\[
%W_{L}^{L\cap ({}^xK)}:=W_{L} \cap W_\ex^{L\cap ({}^xK)}=
%\{w\in W_{L} \mid wx\in W_\ex^{(K,J)}\}. 
%\]
\item 
\label{it:bi-min-3}
Let $x\in {}^{L} W_\ex^{(K,J)}$. 
We have $x^{-1}W_{L} x\cap W_K = W_{(L^x)\cap K}$. Moreover, for any $w \in {}^{(L^x)\cap K}W_K$ we have $xw\in {}^{L}W_\ex$, and the minimal representatives in all cosets $W_L xw$ with $w \in W_K$ are of this form.
%We have $x^{-1}W_{L} x\cap W_K = W_{(L^x)\cap K}$, and 
%\[
%{}^{(L^x)\cap K}W_K:=W_K\cap ({}^{(L^x)\cap K}W_\ex)=
%\{w\in W_K\mid xw\in {}^{L}W_\ex\}. 
%\]
\end{enumerate}
\end{lem}

\begin{proof}
All the statements reduce to well-known facts about the subsets ${}^{L} W_\ex^K$, using the facts that for $\theta \in \Xd_J^0$, $t_\theta$ commutes with $W_{L}$ and 
%(as observed in the course of the proof of Lemma~\ref{lem:min-length}) 
left multiplication by $t_\theta$ stabilizes 
${}^{L}W_\ex$ and
$W_\ex^{(K,J)}$ (see Lemma~\ref{lem:min-length}\eqref{it:min-length-0}). More specifically, one proceeds as follows.

%We only show the third bijection; the other ones can be proved similarly. Let $\theta\in \Xd_J$. 
%As observed in the course of the proof of Lemma~\ref{lem:min-length}, left multiplication by $t_\theta$ stabilizes ${}^JW_\ex$. This operation also clearly stabilizes $W_\ex^{(K,J)}$, hence it stabilizes ${}^JW_\ex^{(K,J)}$. 
For~\eqref{it:bi-min-1}, we start with injectivity.
Let $x,y\in {}^{L}W_\ex^{(K,J)}$ be two elements which have the same image, i.e.~such that $W_{L} x W_K = W_{L} y W_K$. 
%For $N \gg 0$, if 
As in the proof of Lemma~\ref{lem:min-length}\eqref{it:min-length-2}, there exists $\theta \in \Xd_J^0$ 
%satisfies $\langle \theta, \beta \rangle \geq N$ for any $\beta \in \Phis_{J^\comp}$ then 
such that $t_\theta x$ and $t_\theta y$ belong to $W_\ex^{K}$, and then these elements belong to ${}^{L} W_\ex^{K}$. Since $W_{L} t_\theta x W_K = t_\theta W_{L} x W_K = t_\theta W_{L} y W_K = W_{L} t_\theta y W_K$, we have $t_\theta x=t_\theta y$, and then $x=y$.

To show surjectivity, we consider a double coset $X \in W_{L} \backslash W_\ex /W_K$.
%, and $N \gg 0$ such that for any element $w \in A$, if 
Since $X$ is finite, there exists $\theta \in \Xd_J^0$ such that for all $w \in X$ we have
%satisfies $\langle \theta, \beta \rangle \geq N$ for any $\beta \in \Phis_{J^\comp}$ we have
\[
\ell^J( w) = \ell(t_\theta w) - \ell(t_\theta).
\]
If $z$ is the unique element in $t_\theta X$ (which is a coset in $W_{L}\backslash W_\ex /W_K$) which belongs to ${}^{L} W_\ex^{K}$, then $t_{-\theta} z$ is an element of $X$ which belongs to ${}^{L} W_\ex^{(K,J)}$.

%we have $t_\theta x\in {}^JW_\ex^{K}$ when $\theta$ is $J^c$-dominant enough. 
%On the other hand, $t_\theta$ acts on $W_J\backslash W_\ex /W_K$ by multiplication from the left. 
%Therefore, our assertion follows from the bijection ${}^JW_\ex^{K}\leftrightarrow W_J\backslash W_\ex /W_K$. 

For~\eqref{it:bi-min-2} and \eqref{it:bi-min-3}, assume first that $J=S$. In this case,~\eqref{it:bi-min-3} follows from~\cite[Lemma~2.1.9 and Theorem~2.1.12]{geck-pfeiffer} and~\eqref{it:bi-min-2} is obtained by applying these statements to $x^{-1}$, with the roles of $K$ and $L$ switched. (These statements are proved for Coxeter groups in~\cite{geck-pfeiffer}, but they follow for an ``almost Coxeter group'' such as $W_\ex$.)
The general case follows,
%is similar, 
using multiplication by an appropriate $t_\theta$ with $\theta \in \Xd_J^0$.
%to reduce the claim to
%: As in (1), by replacing $x$ by $t_\theta x$ (where $\theta\in \Xd_J$ is $J^c$-dominant enough), we can assume that $x\in {}^JW_\ex^{K}$ and on $W_J xW_K$ the order $\preceq^J$ coincides with the Bruhat order. 
%Then the assertion follows from 
%the corresponding fact for the Bruhat order. (REFERENCE?)
\end{proof}

For $J$, $K$ as above, we now set
\[
W_\ex^{(K,J),\max} := W_\ex^{(K,J)} \cdot w_K,
\]
%Then 
so that by Lemma~\ref{lem:min-length}\eqref{it:min-length-2}
we have
\[
W_\ex^{(K,J),\max} = \{x \in W_\ex \mid \forall w\in W_K, \, \ell^{J}(x)\geq \ell^{J}(xw) \}.
\]
For $L\subset S$ we set 
\[
{}^{L}W_\ex^{(K,J),\max}:={}^{L}W_\ex\cap W_\ex^{(K,{J}),\max}. 
\]
Here, one can check (using e.g.~Lemma~\ref{lem:quar-min-max} below) that the natural map
\[
{}^{L}W_\ex^{(K,J),\max} \to W_{L} \backslash W_\ex / W_K
\]
is injective, but it is \emph{not} surjective in general.

\begin{lem}
\label{lem:quar-min-max}
Let $L \subset J\subset S$ and $K\subset S_\af$ be finitary.
For 
%any subset $J' \subset J$ and 
any $x\in W_\ex$, the following properties are equivalent:
\begin{enumerate}
\item 
$x\in {}^{L} W_\ex^{(K,J),\max}$; 
\item 
$x\in {}^{L} W_\ex$, and for any $w\in W_{L}$ we have $wx\in W_\ex^{(K,J),\max}$; 
\item 
$x\in W_\ex^{(K,J),\max}$, and for any $u\in W_K$ we have $xu\in {}^{L} W_\ex$; 
\item 
for any $w\in W_{L}$ and $u\in W_K$, we have $\ell^J(wxu)=\ell(w)+\ell^J(x)-\ell(u)$. 
\end{enumerate}
\end{lem}

\begin{proof}
%Let $\theta\in \Xd_J$. 
%By Lemma~\ref{lem:min-length}\eqref{it:min-length-2}, $t_\theta$ stabilizes $W_\ex^{(K,J),\mathrm{w}}$. 
%And for any $x\in W_\ex^{(K,J),\mathrm{w}}$, we have $t_\theta x\in W_\ex^{K,\mathrm{w}}$ when $\theta$ is $J^c$-dominant enough. 
%Hence the assertion can be deduced from the analogue one for the Bruhat order and length $\ell$. 
As for Lemma~\ref{lem:bi-min}, the assertion follows from the analogous statement for the usual length function, which follows (using right multiplication by $w_K$) from~\cite[Lemma~2.4]{AR23}.
\end{proof}

%\begin{rmk}
%\label{rmk:intersection-max-elt}
%If $x \in {}^L W_\ex^{(K,J),\max}$, then the element $xw_K \in {}^L W_\ex^{(K,J)}$ satisfies
%\[
%\ell^J(w(xw_K)u) = \ell(w) + \ell^J(xw_K) + \ell(u)
%\]
%for any $w \in W_L$ and $u \in W_K$. It follows that $L \cap ({}^{xw_K} K) = \varnothing$, i.e.~$L \cap ({}^x K) = \varnothing$.
%\end{rmk}

In the following lemma we consider the special case $K=S$ and $L=J$.

\begin{lem}
\label{lem:JSJw}
For any $J \subset S$ we have 
\[
{}^JW_\ex^{(S,J),\max}=\{w_Jt_\lambda w_0 : \lambda\in \Xd_J^{++}\}. 
\]
\end{lem}

\begin{proof}
This statement will be deduced from~\cite[Lemma~10.2]{AR18}. (This reference has some technical assumptions on the reductive group under consideration, but the general case can be easily reduced to that of~\cite{AR18}.) The latter statement describes the elements $x \in W_\ex$ which are minimal in $Wx$ and maximal in $xW_J$, with respect to the ``other'' Coxeter structure mentioned in Remark~\ref{rmk:Coxeter-structures}. Noting that the assignment $x \mapsto w_J x^{-1} w_0$ induces a bijection between this set and 
the set of elements $y$ which are maximal in $yW$ and minimal in $W_J y$, and taking into account the difference of Coxeter group structures, we obtain that with our present conventions the set of elements $y$ which are maximal in $yW$ and minimal in $W_J y$ consists of the elements $(w_J t_\nu v_\nu w_0 : \nu \in \Xd_J^{++})$ where $v_\nu \in W$ is the element of minimal length such that $v_\nu(\nu) \in \Xd^+$.

Let now $x \in {}^JW_\ex^{(S,J),\max}$, and choose $\theta \in \Xd_J^0$ such that
\[
\ell^J(w) = \ell(t_\theta w)-\ell(t_\theta)
\]
for any $w \in x W$. Writing $x=t_\mu y$ with $\mu \in \Xd$ and $y \in W$, we can (and shall) also assume that 
$\langle \theta+\mu, w_J(\alpha) \rangle \geq 0$ for any $\alpha \in \Phis_{J^\comp}$.
%$\theta+\mu \in \Xd_{J^\comp}^{+}$. 
Then $t_\theta x$ is minimal in $W_J t_\theta x$ by Lemma~\ref{lem:min-length}\eqref{it:min-length-0} and maximal in $t_\theta x W$ by our choice of $\theta$, hence $t_\theta x = w_J t_\nu v_\nu w_0$ for some $\nu \in \Xd_J^{++}$. We then have
\[
t_\theta x = t_{\theta+\mu} y = t_{w_J(\nu)} w_J v_\nu w_0.
\]
It follows that $\theta+\mu = w_J(\nu)$, hence $\langle \nu, \alpha \rangle \geq 0$ for $\alpha \in \Phis_{J^\comp}$, so that $\nu \in \Xd^+$ and therefore $v_\nu=1$. We deduce that
\[
x = w_J t_{\nu-\theta} w_0,
\]
hence $x$ has the required form with $\lambda=\nu-\theta \in \Xd_J^{++}$.

Conversely, let $\lambda \in \Xd_J^{++}$, and set $x = w_J t_\lambda w_0$. Choose $\theta \in \Xd_J^0$ such that 
\[
\ell^J(w) = \ell(t_\theta w)-\ell(t_\theta)
\]
for any $w \in x W$, and moreover
$\lambda+\theta \in \Xd^+$. Then the element $t_\theta x = w_J t_{\lambda + \theta} w_0$ is minimal in $W_J t_\theta x$ and maximal in $t_\theta x W$. It follows that $x \in W_\ex^{(S,J),\max}$ by our choice of $\theta$, and that $x \in {}^J W_\ex$ by Lemma~\ref{lem:min-length}\eqref{it:min-length-0}.
%, and therefore belongs to ${}^JW_\ex^{(S,J),\max}$ by our assumptions on $\theta$. 
Hence $x \in {}^JW_\ex^{(S,J),\max}$, as desired.
%It is direct to check that $w_Jt_\lambda w_0\in {}^JW_\ex^{(S,J),\max}$, provided $\lambda\in \Xd_J^{++}$. 
%Suppose $x=w't_\lambda w\in {}^JW_\ex^{(S,J),\max}$, where $w'\in W_J$, $w\in W$ and $\lambda\in \Xd$. Let also also 
%%$N>0$ be such that for 
%$\theta\in \Xd_J^0$ 
%%which satisfies $\langle \theta, \beta \rangle \geq N$ for any $\beta \in \Phis_{J^\comp}$, 
%be such that
%for any $y \in W_J x W$ we have $\ell^J(y)=\ell(t_\theta y)-\ell(t_\theta)$.
%%Let $\theta\in \Xd_J$ be $J^c$-dominant enough such that 
%Then $t_\theta x=w't_{\lambda+\theta}w\in {}^JW_\ex^{S,\max}$. 
%It forces that ${\lambda+\theta}\in \Xd_J^{++}$, see e.g. \cite[Lemma~10.2]{AR18}. 
%Therefore we have $\lambda\in \Xd_J^{++}$. 
%By Lemma~\ref{lem:quar-min-max}, each $(W_J,W)$-coset in $W_\ex$ can have at most one element in ${}^JW_\ex^{(S,J),\max}$. 
%The conclusion follows. 
\end{proof}

\begin{rmk}
\label{rmk:JSJw-length}
A direct computation shows that if $\lambda\in \Xd_J^{++}$ we have $\ell^J(w_Jt_\lambda w_0)= \langle \lambda, 2\rho \rangle + \ell(w_0) - \ell(w_J)$.
\end{rmk}

%---------------------------------------------------------------
\subsection{Affine flag varieties} 
%---------------------------------------------------------------

%===============================
\subsubsection{Loop groups, parahoric subgroups, and affine flag varieties} 
\label{sss:def-Fl}
%===============================

For an affine scheme $X$ over $\F$, we denote by 
\[
\sL{X},\quad \text{resp.}\ \sL^+X, \quad \text{resp.}\ \sL^-X,
\]
the ind-affine ind-scheme, resp. affine scheme, resp. ind-affine ind-scheme, represen\-ting the functor sending an $\F$-algebra $R$ to $X(R( \hspace{-1pt} (z) \hspace{-1pt} ))$, resp.~$X(R[ \hspace{-1pt} [z] \hspace{-1pt} ])$, resp.~$X(R[z^{-1}])$, where $z$ is an indeterminate. 
There is a natural morphism
\[
\mathrm{ev}^X_+:\sL^+X\rightarrow X, \quad \text{resp.} \quad \mathrm{ev}^X_-:\sL^-X\rightarrow X,
\]
induced by evaluating $z$, resp. $z^{-1}$, to zero. 
When $X$ is an affine group scheme, $\sL{X}$, resp. $\sL^+X$, resp. $\sL^-X$, admits a structure of group ind-scheme, resp. group scheme, resp. group ind-scheme. 
In this case we set
\[
\sL^{++}X=\ker(\mathrm{ev}^X_+), \quad \sL^{--}X=\ker(\mathrm{ev}^X_-).
\] 
We call $\sL X$ the loop group of $X$, and $\sL^+X$ the arc group of $X$. 

We will consider in particular these constructions for our groups $T$, $U$ and $G$.
For $\lambda\in \Xd$, we denote by $z^\lambda$ the $\F$-point of $\sL{G}$ given by the image of the $\F$-point $z\in \F(\!(z)\!)^\times = \sL\Gm(\F)$ under the composition $\sL\Gm\xrightarrow{\sL\lambda} \sL{T} \rightarrow \sL{G}$. 

We define the Iwahori subgroup $\I$ of $\sL^+ G$, resp.~its pro-unipotent radical $\Iu$, as
\[
\I=(\mathrm{ev}^G_+)^{-1}(B), \quad \text{resp.} \quad \Iu=(\mathrm{ev}_+^G)^{-1}(U).
\] 
For any finitary subset $J \subset S_\af$, we have an associated parahoric subgroup $\I_J \subset \sL G$ containing $\I$, see e.g.~\cite[\S 3.4]{AR23}. If we set $\Iu^J := \dot{w}_J \Iu (\dot{w}_J)^{-1}$, the intersection $\Iu \cap \Iu^J$ is the pro-unipotent radical of $\I_J$, and we have a canonical Levi factor $M_J \subset \I_J$ such that multiplication induces an isomorphism
\[
M_J \ltimes (\Iu \cap \Iu^J) \simto \I_J.
\]
%The case 
If $J \subset S$ 
%will play a special role below. In this case 
we have
\[
\I_{J}=(\mathrm{ev}^G_+)^{-1}(P_J), 
%\quad \subset \sL^+ G,
\]
and $M_J$ is the subgroup of $G$ defined in~\S\ref{ss:roots-weights} (now seen as a subgroup of $\sL G$).
%The intersection $B \cap M_J \subset M_J$ is a Borel subgroup of $M_J$, and 

%We will now introduce more subgroups of $\sL G$ associated with subsets of $S$.
If $L \subset J \subset S$, then $J$ is the set of simple reflections of the Weyl group $W_J$ of $(M_J,T)$, and we will denote by
%Given $J\subset S$, we will consider the similarly defined Iwahori subgroup 
$\I_{L \subset J}\subset \sL^+M_J$ the parahoric subgroup of $\sL^+ M_J$ associated with the subset $L \subset J$. (Here, the fixed Borel subgroup of $M_J$ used for this construction is $B_J$.)
%of $\sL{G}$ associated to $J$. 
We also set
\[
\I^\semi_{L \subset J} := \I_{L \subset J} \ltimes \sL U_J.
\]
The following special cases are particularly relevant.
\begin{itemize}
\item
When $L=J$, we will simplify the notation and set $\I_{J}^{\semi} := \I_{J \subset J}^{\semi}$. This group ind-scheme
%, called the semiinfinite parahoric subgroup associated with $J$, 
is considered in particular in~\cite{DL23}.
\item
When $L=\varnothing$, we will set $\I_{J}^{\quar}:= \I_{\varnothing \subset J}^\semi$, and call this group ind-scheme the quarterinfinite Iwahori subgroup associated with $J$. This object appears in the discussion in~\cite[\S 8.11]{bez-losev}.
\item
When $J=\varnothing$ (so that $L=\varnothing$) we have $\I_{\varnothing}^{\quar}=\I_{\varnothing}^{\semi}$, and this group coincides with the \textit{semiinfinite Iwahori subgroup} $\I^\semi$ considered in~\cite{adr} (and earlier, using different notation, in various references including~\cite{Gaits18}).
\item
When $J=S$, we have $\I_{L \subset S}^\semi = \I_L$; in particular $\I_{S}^\quar = \I$ and $\I_S^\semi = \sL^+ G$.
\end{itemize}

%We also define the \textit{quarter-infinite Iwahori subgroup} and the \textit{semi-infinite parahoric subgroup} of $\sL G$ associated with $J$ as 
%\[
%\I_{J}^{\quar}:= \I_{M_J}\ltimes \sL U_{J}, \quad
%%\]
%%where $\I_{M_J}\subset \sL^+M_J$ is the Iwahori subgroup of $M_J$ defined as for $\I$ (using the Borel subgroup $B \cap M_J$ of $M_J$). 
%%We define the \textit{semi-infinite parahoric subgroup associated to $J$} as 
%%\[
%\I_{J}^{\semi}:= \sL^+M_J\ltimes \sL U_{J}. 
%\]
%When $J=\varnothing$, we have
%\[
%\I_{\varnothing}^{\quar}=\I_{\varnothing}^{\semi},
%\]
%and this group coincides with the \textit{semi-infinite Iwahori subgroup} $\I^\semi$ considered in~\cite{adr}.
%When $J=S$, we have
%\[
%\I_{S}^{\quar}=\I \quad \text{and} \quad \I_{S}^{\semi}=\sL^+G.
%\]
%For general $J$, the group ind-scheme $\I_{J}^{\quar}$ appears in the discussion in~\cite[\S 8.11]{bez-losev}, while $\I_{J}^{\semi}$ is used in~\cite{DL23}.

We define the affine flag variety and the enhanced affine flag variety  by 
\[
\Fl=\sL G/\I, \quad \tFl=\sL G/\Iu.
\]
For a finitary subset $K\subset S_\af$, we also have the partial affine flag variety
\[
\Fl_K=\sL G/\I_K. 
\]
We will denote by $q:\tFl\rightarrow \Fl$ and $p_K: \Fl\rightarrow \Fl_K$ the natural projections. In case $K=S$, $\Fl_S$ is the familiar affine Grassmannian $\Gr = \sL G / \sL^+ G$.

%===============================
\subsubsection{Orbits} 
%===============================

For $x=t_\lambda w$ ($\lambda\in \Xd$, $w\in W$), we set $\dot{x}=z^\lambda \dot{w}\in \sL{G}(\F)$. 
Let $L \subset J \subset S$, and let $K \subset S_\af$ be a finitary subset.
%\footnote{SR: Parahoric subgroups can be defined more generally for $K \subset S_\af$. Below, can we consider more generally a finitary subset $K \subset S_\af$?} 
We have decompositions of $\Fl_K$ into $\I_{L \subset J}^\semi$-orbits
%$\I_J^{\quar}$, $\I_J^{\semi}$, and $\I_J$-orbits 
\begin{equation}
\label{eq:decomp}
\Fl_K= \bigsqcup_{x\in {}^L W_\ex^{(K,J)}} \rS^{L \subset J}_{K,x},
\end{equation}
where 
\[
\rS^{L \subset J}_{K,x}:= \I_{L \subset J}^{\semi} \dot{x} \I_K/\I_K.
\]
When $K=S$ and $L=J=\varnothing$, these orbits are the familiar ``semiinfinite orbits'' which play a crucial role in many constructions related to the \emph{geometric Satake equivalence}, see~\cite{mv}. When $K=\varnothing$ and $L=J=\varnothing$, these orbits play a role in some constructions in~\cite{ab} (see also~\cite[\S 4.5]{ARbook}).
When $J=S$ (so that $\I_{L \subset S}^\semi = \I_L$),
$\rS^{L \subset S}_{K,x}$ is a (smooth) scheme of finite type over $\F$, and we have
\begin{equation}
\label{eqn:dim-par-orbit}
\dim(\rS^{L \subset S}_{K,x}) = \ell(x) + \ell(w_L w_{L \cap ({}^xK)}) = \ell(x) + \ell(w_L) -\ell(w_{L \cap ({}^xK)}).
\end{equation}
(This formula follows e.g.~from standard facts about the Bruhat decomposition, together with Lemma~\ref{lem:bi-min}\eqref{it:bi-min-2}.)

We denote by $\bs^{L \subset J}_{K,x}$ the locally closed embedding of $\rS^{L \subset J}_{K,x}$ in $\Fl_K$. 
We will omit $L$, $J$ or $K$, in the notations above if it is the empty set. (In particular, $\rS^J_{K,x}$ will mean $\rS^{\varnothing \subset J}_{K,x}$, the $\I_J^\quar$-orbit of $\dot{x} \I_K/\I_K$.)
%(Note that the $\I^{\semi}$-orbits $\rQ_x^S=\rS_x^S$ in $\Fl^S = \Gr$ are the familiar ``semiinfinite orbits'' which play a crucial role in many constructions related to the \emph{geometric Satake equivalence}.)
We also have a decomposition 
\begin{equation*}
%\label{equ decomp-tFl}
\tFl = \bigsqcup_{x\in {}^L W_\ex} q^{-1}(\rS^{L \subset J}_x)
\end{equation*}
into $\I_{L \subset J}^{\semi}$-orbits.

For $\alpha\in \Phi$ and $n\in \Z$, we denote by $U_{\alpha+n\hbar}$ the one-parameter subgroup in $\sL G$ given by the image of $a\in\Ga \mapsto u_\alpha(az^n)$. The following lemma, which can be checked by standard computations, describes the $\I^\quar_J$-orbits on $\Fl$.
%The following lemma (which generalizes~\cite[Lemma~3.2]{adr}) is standard. Here and below, 
Here and below, a product $\prod_{n=0}^k$ with $k<0$ or a product $\prod_{n=1}^k$ with $k \leq 0$ should be interpreted as $\Spec(\F)$.

\begin{lem}
Let $x=t_\lambda w$ ($\lambda\in \Xd$, $w\in W$).
Multiplication on $\dot{x}\I/\I$ induces an isomorphism 
\begin{equation}
\label{eq:quar-orbit}
\begin{aligned}
&\left(\prod_{\substack{\alpha \in \Phi_J^+ \\ w^{-1}(\alpha)>0}} \prod_{n=0}^{\langle\lambda,\alpha\rangle-1}U_{\alpha+n\hbar} \right) \times 
\left(\prod_{\substack{\alpha \in \Phi_J^+ \\ w^{-1}(\alpha)<0}} \prod_{n=0}^{\langle\lambda,\alpha\rangle}U_{\alpha+n\hbar} \right) \times \\ 
&\qquad \left(\prod_{\substack{\alpha \in -\Phi_J^+ \\ w^{-1}(\alpha)>0}} \prod_{n=1}^{\langle\lambda,\alpha\rangle-1}U_{\alpha+n\hbar} \right) \times 
\left(\prod_{\substack{\alpha \in -\Phi_J^+ \\ w^{-1}(\alpha)<0}} \prod_{n=1}^{\langle\lambda,\alpha\rangle}U_{\alpha+n\hbar} \right) \times \\ 
&\qquad\qquad \left(\prod_{\substack{\alpha \in \Phi^+\smallsetminus \Phi_J^+ \\ w^{-1}(\alpha)>0}} \Ad_{z^{\lambda}}(\sL^{--} U_{\alpha}) \right) \times 
\left(\prod_{\substack{\alpha \in \Phi^+\smallsetminus \Phi_J^+ \\ w^{-1}(\alpha)<0}} \Ad_{z^{\lambda}}(\sL^{-} U_{\alpha}) \right) 
\xs \rS^{J}_{x}. 
\end{aligned}
\end{equation}
\end{lem}

Note that the affine root subgroups appearing in the left-hand side of~\eqref{eq:quar-orbit} are parametrized by
\[
\{ \alpha \in {}^J \Phi_\af^+ \mid x^{-1}(\alpha)<0\}.
\]
In the special case $J=S$, we obtain an isomorphism
\begin{equation*}
%\label{eq:Iwa-orbit}
\begin{aligned}
&\left(\prod_{\substack{\alpha \in \Phi^+ \\ w^{-1}(\alpha)>0}} \prod_{n=0}^{\langle\lambda,\alpha\rangle-1}U_{\alpha+n\hbar} \right) \times 
\left(\prod_{\substack{\alpha \in \Phi^+ \\ w^{-1}(\alpha)<0}} \prod_{n=0}^{\langle\lambda,\alpha\rangle}U_{\alpha+n\hbar} \right) \times \\
&\qquad\qquad \left(\prod_{\substack{\alpha \in -\Phi^+ \\ w^{-1}(\alpha)>0}} \prod_{n=1}^{\langle\lambda,\alpha\rangle-1}U_{\alpha+n\hbar} \right) \times 
\left(\prod_{\substack{\alpha \in -\Phi^+ \\ w^{-1}(\alpha)<0}} \prod_{n=1}^{\langle\lambda,\alpha\rangle}U_{\alpha+n\hbar} \right)
\xs \I \dot{x} \I/\I, 
\end{aligned}
\end{equation*}
which generalizes~\cite[Lemma~3.2]{adr}.

%The following lemmas will be needed in the sequel. 

Given locally closed sub-ind-schemes $X,Y \subset \Fl$, we will write $X \times^\I Y$ for the twisted product $\widetilde{X} \times^\I Y$ where $\widetilde{X}$ is the preimage of $X$ in $\sL G$.

\begin{lem}
\label{lem:min-push}
Let $J \subset S$, and let $K \subset S_\af$ be finitary.
Let $x\in W_\ex^{(K,J)}$.

\begin{enumerate}
\item
\label{it:min-push-1}
We have 
\[
(p_K)^{-1}(\rS^J_{K,x})=\bigsqcup_{{w\in W_K}} \rS^J_{xw},
%\quad \text{and}\quad 
%p_K:{}^J\rQ_{x}\xs {}^J\rQ^K_{x}. 
\]
and $p_K$ restricts to an isomorphism $\rS^J_{x}\xs \rS^J_{K,x}$.
\item
\label{it:min-push-2}
For any $w\in W_K$, multiplication induces an isomorphism
\[
%{}^J\rQ_{x}\times^{B}B \dot{w}B/B \xs 
\rS^J_{x}\times^{\I}\I \dot{w}\I/\I \xs 
\rS^J_{xw}. 
\]
\end{enumerate}
\end{lem}

\begin{proof}
This follows from~\eqref{eq:quar-orbit} 
%\eqref{eq:Iwa-orbit}, Lemma~\ref{lem:quar-order} 
and the techniques from Lemma~\ref{lem:min-length}. For instance, \eqref{it:min-push-2} follows from the equality
\begin{multline*}
\{ \alpha \in {}^J \Phi_\af^+ \mid (xw)^{-1}(\alpha)<0\} = \\
\{ \alpha \in {}^J \Phi_\af^+ \mid x^{-1}(\alpha)<0\} \sqcup x(\{\alpha \in \Phi_\af^+ \mid w^{-1}(\alpha)<0\}),
\end{multline*}
which can be deduced from the corresponding fact for the standard Bruhat order (corresponding to $J=S$) using the fact that if $\alpha \in {}^J \Phi_\af^+$, we have $t_\theta(\alpha) \in \Phi_\af^+$ if $\theta \in \Xd^0_J$ satisfies $\langle \theta, \beta \rangle \gg 0$ for any $\beta \in \Phis_{J^\comp}$.
\end{proof}

\begin{lem}
\label{lem:bi-min-av}
Let $L \subset J \subset S$, and let $K \subset S_\af$ be finitary.
Let $x\in {}^L W_\ex^{(K,J)}$. 
\begin{enumerate}
\item 
\label{it:bi-min-av-1}
For any $w\in W_L^{L\cap ({}^xK)}$, we have canonical isomorphisms 
\[
B\dot{w}B\times^{B} (\rS_{K,x}^{J}) \xs \I_J^{\quar}\dot{w}\I_J^{\quar}\times^{\I_J^{\quar}} (\rS_{K,x}^{J})
 \xs \rS_{K,wx}^J,
\]
where the second map is induced by multiplication. 
\item 
\label{it:bi-min-av-2}
The morphism
\[
\I^\semi_{L \subset J} \times^{\I^\quar_J} \rS_{K,x}^J  \rightarrow \rS_{K,x}^{L \subset J}
\]
induced by multiplication
is a locally trivial fibration with fiber isomorphic to the flag variety of $M_{L \cap ({}^xK)}$. 
\item 
\label{it:bi-min-av-3}
We have 
\[
\rS_{K,x}^{L \subset J}
=\bigsqcup_{w\in W_L^{L\cap ({}^xK)}} \rS^J_{K,wx}. 
\]
\end{enumerate}
\end{lem}

\begin{proof}
\eqref{it:bi-min-av-1}
By Lemma~\ref{lem:bi-min}\eqref{it:bi-min-2} we have $wx \in W_\ex^{(K,J)}$. Then by Lemma~\ref{lem:min-push}\eqref{it:min-push-1} we can assume that $K=\varnothing$. In this case, 
the proof is similar to that of Lemma~\ref{lem:min-push}\eqref{it:min-push-2}, the corresponding equality being
\begin{multline*}
\{\alpha \in {}^J \Phi_\af^+ \mid (wx)^{-1} (\alpha) <0\} = \\
\{\alpha \in \Phi_\af^+ \mid w^{-1}(\alpha)<0\} \sqcup w(\{\alpha \in {}^J \Phi_\af^+ \mid x^{-1}(\alpha)<0\}).
\end{multline*}
%the claim follows from~\eqref{eq:quar-orbit} together with Lemmas~\ref{lem:quar-order} and~\ref{lem:min-length}\eqref{it:min-length-1}.
%This follows from~\eqref{eq:Iwa-orbit}, \eqref{eq:quar-orbit}, and Lemmas~\ref{lem:quar-order}, \ref{lem:min-length}\ref{it:min-length-1} and \ref{lem:bi-min}\eqref{it:bi-min-2}. 
%(ADD MORE DETAILS HERE.)

\eqref{it:bi-min-av-2}
Since $\rS_{K,x}^{L \subset J}$ is an $\I^\semi_{L \subset J}$-orbit,
it suffices to analyze the fiber of this map over the image of $\dot{x}$.
% along $\I^\semi_J \times^{\I^\quar_J} {}^J\rQ_{x}^K  \rightarrow {}^J\rS_{x}^{K}$ 
This fiber is isomorphic to 
\[
\I^\semi_{L \subset J} \cap \Ad_{\dot{x}}(\I_K)/\I^\quar_J\cap \Ad_{\dot{x}}(\I_K)
\simeq \bigl( M_L \cap \Ad_{\dot{x}}(\I_K) \bigr) / \bigl( B_L \cap \Ad_{\dot{x}}(\I_K) \bigr). 
\]
Here, since $x \in {}^L W_\ex$, by~\eqref{eq:Bruhat} we have $x^{-1}(\Phi_L^+) \subset \Phi_\af^+$, i.e.~$\Phi_L^+ \subset x(\Phi_\af^+)$, hence $B_L \subset \Ad_{\dot{x}}(\I_K)$. On the other hand we have $W_L \cap x W_K x^{-1} = W_{L \cap ({}^xK)}$ by Lem\-ma~\ref{lem:bi-min}\eqref{it:bi-min-2}, hence $(M_L \cap \Ad_{\dot{x}}(\I_K)) / (B_L \cap \Ad_{\dot{x}}(\I_K))$ identifies with the flag variety of the Levi factor $M_{L\cap ({}^xK)}$ of $M_L$.

%The group $P_J\cap \Ad_{\dot{x}}(\I_K)$ admits a Levi factor given by $M=M_J\cap \Ad_{\dot{x}}(M_K)$, which is the reductive group associated with the root datum $(\X,\Phi_J\cap x(\Phi_K),\Xd,\check{\Phi}_J\cap x(\check{\Phi}_K))$ (where $\check{\Phi}_J$, $\check{\Phi}_K$ are the coroots corresponding to $\Phi_J$, $\Phi_K$). 
%Moreover $M\cap B$ a Borel subgroup of $M$.  
%By Lemma~\ref{lem:bi-min}\eqref{it:bi-min-2}, we have $M=M_{J\cap ({}^xK)}$. 
%Since $x\in {}^JW_\ex^{(K,J)}$, by \eqref{eq:Bruhat} and Lemma~\ref{lem:quar-order} we have $x^{-1}(\Phi^+_J)\subset \Phi^+_\af$ and $x(\Phi^+_K)\subset {}^J\Phi^+_\af$.  
%Hence by comparing the tangent spaces at the base point, the inclusion $M \subset P_J\cap \Ad_{\dot{x}}(\I_K)$ induces an isomorphism 
%\[
%M/M\cap B \xs P_J\cap \Ad_{\dot{x}}(\I_K)/B\cap \Ad_{\dot{x}}(\I_K) . 
%\]

\eqref{it:bi-min-av-3}
It is clear that $\rS_{K,x}^{L \subset J}$ is the union of the $\I^\quar_J$-orbits associated with the representatives in $W_\ex^{(K,J)}$ of the cosets $wxW_K$ for $w \in W_L$. The claim therefore follows from
Lemma~\ref{lem:bi-min}\eqref{it:bi-min-2}.
%, for any $w \in W_J^{J \cap ({}^x K)}$ we have $wx \in W_\ex^{(K,J)}$.
%
%We have 
%\[
%{}^J\rS_{x}^{K}
%=\bigsqcup_{\substack{w\in W_J, \\ wx\in W_\ex^{(K,J)}}} {}^J\rQ_{wx}^K,
%\] 
%therefore the assertion follows from Lemma~\ref{lem:bi-min}\eqref{it:bi-min-2}. 
\end{proof}

\begin{lem}
\label{lem:bi-min-push}
Let $L \subset J \subset S$, and let $K \subset S_\af$ be finitary.
Let $x\in {}^L W_\ex^{(K,J)}$. 
\begin{enumerate}
\item 
\label{it:bi-min-push-1}
For any $w\in {}^{(L^x)\cap K}W_K$, multiplication induces an isomorphism
%we have isomorphisms 
\[
%{}^J\rS_{x}\times^{B}B\dot{w}B/B \xs 
\rS^{L \subset J}_{x}\times^{\I}\I \dot{w}\I/\I \xs \rS^{L \subset J}_{xw}.
\]
%where the last map is induced by multiplication. 
\item
\label{it:bi-min-push-2}
The morphism $p_K: \rS^{L \subset J}_{x}\rightarrow \rS^{L \subset J}_{K,x}$ is a locally trivial fibration with fiber isomorphic to the flag variety of $M_{(L^x) \cap K}$. 
\item
\label{it:bi-min-push-3}
We have 
\[
(p_K)^{-1}(\rS^{L \subset J}_{K,x})=\bigsqcup_{w\in {}^{(L^x) \cap K}W_K} \rS^{L \subset J}_{xw}. 
\] 
\end{enumerate}
\end{lem}

\begin{proof}
%It is proved similarly as 
The proof is similar to that of
Lemma~\ref{lem:bi-min-av}, using Lemma~\ref{lem:bi-min}\eqref{it:bi-min-3}.
%, where (2) uses Lemma~\ref{lem:bi-min}(3). 
\end{proof}



%OLD VERSION:

%%===============================
%\subsubsection{Orbits} 
%%===============================
%
%For $x=t_\lambda w$ ($\lambda\in \Xd$, $w\in W$), we set $\dot{x}=z^\lambda \dot{w}\in \sL{G}(\F)$. 
%Let $J\subset S$, and let $K \subset S_\af$ be a finitary subset.
%%\footnote{SR: Parahoric subgroups can be defined more generally for $K \subset S_\af$. Below, can we consider more generally a finitary subset $K \subset S_\af$?} 
%We have decompositions of $\Fl_K$ into $\I_J^{\quar}$, $\I_J^{\semi}$, and $\I_J$-orbits 
%\begin{equation}\label{equ decomp}
%\Fl_K= \bigsqcup_{x\in W_\ex^{(K,J)}} {}^J\rQ^{K}_x = \bigsqcup_{x\in {}^JW_\ex^{(K,J)}} {}^J\rS^{K}_x = \bigsqcup_{x\in {}^JW_\ex^K} {}^J\Fl^{K}_x, 
%\end{equation}
%where 
%\[
%{}^J\rQ^K_x:= \I_J^{\quar} \dot{x} \I_K/\I_K, 
%\quad {}^J\rS^K_x:= \I_J^{\semi} \dot{x} \I_K/\I_K, 
%\quad \text{and}\quad {}^J\Fl^{K}_x=\I_J \dot{x} \I_K/\I_K. 
%\]
%Here ${}^J\Fl^{K}_x$ is a (smooth) scheme of finite type over $\F$, and we have
%\begin{equation}
%\label{eqn:dim-par-orbit}
%\dim({}^J\Fl^{K}_x) = \ell(x) + \ell(w_Jw_{J\cap ({}^xK)}) = \ell(x) + \ell(w_J) -\ell(w_{J\cap ({}^xK)}).
%\end{equation}
%(This formula follows from standard facts about the Bruhat decomposition, together with Lemma~\ref{lem:bi-min}\eqref{it:bi-min-2}.)
%
%We denote by ${}^J\bq^K_x$, ${}^J\bs^K_x$, resp. ${}^J\bi^K_x$, the locally closed embedding of ${}^J\rQ^K_x$, ${}^J\rS^K_x$, resp. ${}^J\Fl^{K}_x$, in $\Fl_K$. 
%We will omit $J$, resp. $K$, in the notations above if it is the empty set. (Note that the $\I^{\semi}$-orbits $\rQ_x^S=\rS_x^S$ in $\Fl^S = \Gr$ are the familiar ``semiinfinite orbits'' which play a crucial role in many constructions related to the \emph{geometric Satake equivalence}.)
%We have decompositions 
%\begin{equation*}
%%\label{equ decomp-tFl}
%\tFl= \bigsqcup_{x\in W_\ex} q^{-1}({}^J\rQ_x) = \bigsqcup_{x\in {}^JW_\ex} q^{-1}({}^J\rS_x)= \bigsqcup_{x\in {}^JW_\ex} q^{-1}({}^J\Fl_x). 
%\end{equation*}
%
%For $\alpha\in \Phi$ and $n\in \Z$, we denote by $U_{\alpha+n\hbar}$ the one-parameter subgroup in $\sL G$ given by the image of $a\in\Ga \mapsto u_\alpha(az^n)$. 
%The following lemma (which generalizes~\cite[Lemma~3.2]{adr}) is standard. Here and below, a product $\prod_{n=0}^k$ with $k<0$ or a product $\prod_{n=1}^k$ with $k \leq 0$ should be interpreted as $\Spec(\F)$.
%
%\begin{lem}
%Let $x=t_\lambda w$ ($\lambda\in \Xd$, $w\in W$). 
%Multiplication on $z^\lambda \dot{w}\I/\I$ induces an isomorphism 
%\begin{equation}
%\label{eq:Iwa-orbit}
%\begin{aligned}
%&\left(\prod_{\substack{\alpha \in \Phi^+ \\ w^{-1}(\alpha)>0}} \prod_{n=0}^{\langle\lambda,\alpha\rangle-1}U_{\alpha+n\hbar} \right) \times 
%\left(\prod_{\substack{\alpha \in \Phi^+ \\ w^{-1}(\alpha)<0}} \prod_{n=0}^{\langle\lambda,\alpha\rangle}U_{\alpha+n\hbar} \right) \times \\
%&\qquad\qquad \left(\prod_{\substack{\alpha \in -\Phi^+ \\ w^{-1}(\alpha)>0}} \prod_{n=1}^{\langle\lambda,\alpha\rangle-1}U_{\alpha+n\hbar} \right) \times 
%\left(\prod_{\substack{\alpha \in -\Phi^+ \\ w^{-1}(\alpha)<0}} \prod_{n=1}^{\langle\lambda,\alpha\rangle}U_{\alpha+n\hbar} \right)
%\xs \Fl_{x}. 
%\end{aligned}
%\end{equation}
%\end{lem}
%
%%Similarly we have the following, whose proof is direct. 
%The following statement can also be checked by direct computation.
%
%\begin{lem}
%Let $x=t_\lambda w$ ($\lambda\in \Xd$, $w\in W$).
%Multiplication on $z^\lambda \dot{w}\I/\I$ induces an isomorphism 
%\begin{equation}
%\label{eq:quar-orbit}
%\begin{aligned}
%&\left(\prod_{\substack{\alpha \in \Phi_J^+ \\ w^{-1}(\alpha)>0}} \prod_{n=0}^{\langle\lambda,\alpha\rangle-1}U_{\alpha+n\hbar} \right) \times 
%\left(\prod_{\substack{\alpha \in \Phi_J^+ \\ w^{-1}(\alpha)<0}} \prod_{n=0}^{\langle\lambda,\alpha\rangle}U_{\alpha+n\hbar} \right) \times \\ 
%&\qquad \left(\prod_{\substack{\alpha \in -\Phi_J^+ \\ w^{-1}(\alpha)>0}} \prod_{n=1}^{\langle\lambda,\alpha\rangle-1}U_{\alpha+n\hbar} \right) \times 
%\left(\prod_{\substack{\alpha \in -\Phi_J^+ \\ w^{-1}(\alpha)<0}} \prod_{n=1}^{\langle\lambda,\alpha\rangle}U_{\alpha+n\hbar} \right) \times \\ 
%&\qquad\qquad \left(\prod_{\substack{\alpha \in \Phi^+\smallsetminus \Phi_J^+ \\ w^{-1}(\alpha)>0}} \Ad_{z^{\lambda}}(\sL^{--} U_{\alpha}) \right) \times 
%\left(\prod_{\substack{\alpha \in \Phi^+\smallsetminus \Phi_J^+ \\ w^{-1}(\alpha)<0}} \Ad_{z^{\lambda}}(\sL^{-} U_{\alpha}) \right) 
%\xs {}^J\rQ_{x}. 
%\end{aligned}
%\end{equation}
%\end{lem}
%
%%The following lemmas will be needed in the sequel. 
%
%\begin{lem}
%\label{lem:min-push}
%Let $x\in W_\ex^{(K,J)}$.
%
%\begin{enumerate}
%\item
%\label{it:min-push-1}
%We have 
%\[
%(p_K)^{-1}({}^J\rQ^K_{x})=\bigsqcup_{{w\in W_K}} {}^J\rQ_{xw},
%%\quad \text{and}\quad 
%%p_K:{}^J\rQ_{x}\xs {}^J\rQ^K_{x}. 
%\]
%and $p_K$ restricts to an isomorphism ${}^J\rQ_{x}\xs {}^J\rQ^K_{x}$.
%\item
%For any $w\in W_K$, multiplication induces an isomorphism
%\[
%%{}^J\rQ_{x}\times^{B}B \dot{w}B/B \xs 
%{}^J\rQ_{x}\times^{\I}\I \dot{w}\I/\I \xs 
%{}^J\rQ_{xw}. 
%\]
%\end{enumerate}
%\end{lem}
%
%\begin{proof}
%This follows from~\eqref{eq:quar-orbit}, 
%%\eqref{eq:Iwa-orbit}, 
%Lemma~\ref{lem:quar-order} and Lemma~\ref{lem:min-length}. 
%\end{proof}
%
%\begin{lem}
%\label{lem:bi-min-av}
%Let $x\in {}^JW_\ex^{(K,J)}$. 
%\begin{enumerate}
%\item 
%\label{it:bi-min-av-1}
%For any $w\in W_J^{J\cap ({}^xK)}$, we have isomorphisms 
%\[
%B\dot{w}B\times^{B} ({}^J\rQ_x^{K})\xs \I_J^{\quar}\dot{w}\I_J^{\quar}\times^{\I_J^{\quar}} ({}^J\rQ_x^{K})
% \xs {}^J\rQ_{wx}^K,
%\]
%where the last map is induced by multiplication. 
%\item 
%\label{it:bi-min-av-2}
%The morphism
%\[
%\I^\semi_J \times^{\I^\quar_J} {}^J\rQ_{x}^K  \rightarrow {}^J\rS_{x}^{K}
%\]
%induced by multiplication
%is a locally trivial fibration with fiber isomorphic to the flag variety of $M_{J\cap ({}^xK)}$. 
%\item 
%\label{it:bi-min-av-3}
%We have 
%\[
%{}^J\rS_{x}^{K}
%=\bigsqcup_{w\in W_J^{J\cap ({}^xK)}} {}^J\rQ_{wx}^K. 
%\]
%\end{enumerate}
%\end{lem}
%
%\begin{proof}
%\eqref{it:bi-min-av-1}
%By \ref{lem:bi-min}\eqref{it:bi-min-2} we have $wx \in W_\ex^{(K,J)}$. Then by Lemma~\ref{lem:min-push}\eqref{it:min-push-1} we can assume that $K=\varnothing$. In this case, the claim follows from~\eqref{eq:quar-orbit} together with Lemmas~\ref{lem:quar-order} and~\ref{lem:min-length}\eqref{it:min-length-1}.
%%This follows from~\eqref{eq:Iwa-orbit}, \eqref{eq:quar-orbit}, and Lemmas~\ref{lem:quar-order}, \ref{lem:min-length}\ref{it:min-length-1} and \ref{lem:bi-min}\eqref{it:bi-min-2}. 
%%(ADD MORE DETAILS HERE.)
%
%\eqref{it:bi-min-av-2}
%Since ${}^J\rS_{x}^{K}$ is an $\I^\semi_J$-orbit,
%it suffices to analyze the fiber of this map over the image of $\dot{x}$.
%% along $\I^\semi_J \times^{\I^\quar_J} {}^J\rQ_{x}^K  \rightarrow {}^J\rS_{x}^{K}$ 
%This fiber is isomorphic to 
%\[
%\I^\semi_J\cap \Ad_{\dot{x}}(\I_K)/\I^\quar_J\cap \Ad_{\dot{x}}(\I_K)
%\simeq M_J\cap \Ad_{\dot{x}}(\I_K)/B_J \cap \Ad_{\dot{x}}(\I_K). 
%\]
%Here, since $x \in {}^J W_\ex$, by~\eqref{eq:Bruhat} we have $x^{-1}(\Phi_J^+) \subset \Phi_\af^+$, i.e.~$\Phi_J^+ \subset x(\Phi_\af^+)$, hence $B_J \cap \Ad_{\dot{x}}(\I_K) = B_J$. On the other hand we have $W_J \cap x W_K x^{-1} = W_{J\cap ({}^xK)}$ by Lemma~\ref{lem:bi-min}\eqref{it:bi-min-2}, hence $M_J\cap \Ad_{\dot{x}}(\I_K)/B_J \cap \Ad_{\dot{x}}(\I_K)$ identifies with the flag variety of the Levi factor $M_{J\cap ({}^xK)}$ of $M_J$.
%
%%The group $P_J\cap \Ad_{\dot{x}}(\I_K)$ admits a Levi factor given by $M=M_J\cap \Ad_{\dot{x}}(M_K)$, which is the reductive group associated with the root datum $(\X,\Phi_J\cap x(\Phi_K),\Xd,\check{\Phi}_J\cap x(\check{\Phi}_K))$ (where $\check{\Phi}_J$, $\check{\Phi}_K$ are the coroots corresponding to $\Phi_J$, $\Phi_K$). 
%%Moreover $M\cap B$ a Borel subgroup of $M$.  
%%By Lemma~\ref{lem:bi-min}\eqref{it:bi-min-2}, we have $M=M_{J\cap ({}^xK)}$. 
%%Since $x\in {}^JW_\ex^{(K,J)}$, by \eqref{eq:Bruhat} and Lemma~\ref{lem:quar-order} we have $x^{-1}(\Phi^+_J)\subset \Phi^+_\af$ and $x(\Phi^+_K)\subset {}^J\Phi^+_\af$.  
%%Hence by comparing the tangent spaces at the base point, the inclusion $M \subset P_J\cap \Ad_{\dot{x}}(\I_K)$ induces an isomorphism 
%%\[
%%M/M\cap B \xs P_J\cap \Ad_{\dot{x}}(\I_K)/B\cap \Ad_{\dot{x}}(\I_K) . 
%%\]
%
%\eqref{it:bi-min-av-3}
%This follows from~\eqref{it:bi-min-av-2} since, by Lemma~\ref{lem:bi-min}\eqref{it:bi-min-2}, for any $w \in W_J^{J \cap ({}^x K)}$ we have $wx \in W_\ex^{(K,J)}$.
%%
%%We have 
%%\[
%%{}^J\rS_{x}^{K}
%%=\bigsqcup_{\substack{w\in W_J, \\ wx\in W_\ex^{(K,J)}}} {}^J\rQ_{wx}^K,
%%\] 
%%therefore the assertion follows from Lemma~\ref{lem:bi-min}\eqref{it:bi-min-2}. 
%\end{proof}
%
%\begin{lem}
%\label{lem:bi-min-push}
%Let $x\in {}^JW_\ex^{(K,J)}$. 
%\begin{enumerate}
%\item 
%\label{it:bi-min-push-1}
%For any $w\in {}^{J^x\cap K}W_K$, multiplication induces an isomorphism
%%we have isomorphisms 
%\[
%%{}^J\rS_{x}\times^{B}B\dot{w}B/B \xs 
%{}^J\rS_{x}\times^{\I}\I \dot{w}\I/\I \xs {}^J\rS_{xw}.
%\]
%%where the last map is induced by multiplication. 
%\item
%\label{it:bi-min-push-2}
%The morphism $p_K:{}^J\rS_{x}\rightarrow {}^J\rS^K_{x}$ is a locally trivial fibration with fiber isomorphic to the flag variety of $M_{(J^x) \cap K}$. 
%\item
%\label{it:bi-min-push-3}
%We have 
%\[
%(p_K)^{-1}({}^J\rS^K_{x})=\bigsqcup_{w\in {}^{(J^x) \cap K}W_K} {}^J\rS_{xw}. 
%\] 
%\end{enumerate}
%\end{lem}
%
%\begin{proof}
%%It is proved similarly as 
%The proof is similar to that of
%Lemma~\ref{lem:bi-min-av}, using Lemma~\ref{lem:bi-min}\eqref{it:bi-min-3}.
%%, where (2) uses Lemma~\ref{lem:bi-min}(3). 
%\end{proof}

%%%%%%%%%%%%%%%%%%%%%%%%%
\section{Parabolic semiinfinite sheaves} 
\label{sec:quar-semi-sheaves}
%%%%%%%%%%%%%%%%%%%%%%%%%

%---------------------------------------------------------------
\subsection{Definition} 
\label{ss:qs-sheaves}
%---------------------------------------------------------------

We will work in the sheaf-theoretic context explained in \cite[\S\S 2.1--2.2]{adr}, in the case of \'etale sheaves.
%So let $\k$ be a field, which will play the role of coefficients for sheaves. 
%%From now on, w
%We assume that the pair $(\F,\k)$ is one of the following: 
%\begin{enumerate}
%\item
%\label{it:sheaves-etale}
%$\F$ is arbitrary, 
So let $\ell$ be a prime number invertible in $\F$, and let $\k$ be either a finite extension or an algebraic closure of $\F_\ell$ or $\Q_\ell$.
%and we consider \'{e}tale $\k$-sheaves; 
%\item 
%$\F=\C$, $\k$ is arbitrary, and we work with sheaves in the analytic topology. 
%\end{enumerate}
For a prestack $X$ over $\F$, we have a $\k$-linear presentable stable $\infty$-category $\Shv(X)$ of $\k$-sheaves on $X$. The functorialities that these $\infty$-categories satisfy are briefly discussed in~\cite[\S 2.2]{adr} (following in particular some constructions from~\cite{zhu}); in particular, for any map $f : X \to Y$ of $\F$-prestacks we have a $!$-pullback functor $f^! : \Shv(Y) \to \Shv(X)$, which sends the dualizing complex $\omega_Y$ of $Y$ to the dualizing complex $\omega_X$ of $X$. These functors, as well as all the functors between such $\infty$-categories that we will consider below, are continuous, i.e.~colimit-preserving.

\begin{rmk}
It is likely that all the results below have analogues for sheaves in the analytic topology. However, since a detailed treatment of the construction of such $\infty$-categories is not yet available in the literature, we will stick to \'etale sheaves. (For the statements involving Whittaker conditions, their counterparts should involve the ``Kirillov model.'')
\end{rmk}

%===============================
%\subsubsection{Definitions} 
%===============================

Fix subsets $L \subset J \subset S$, and
let $X$ be either $\Fl_K$ (for some finitary subset $K \subset S_\af$) or $\tFl$. 
Below we will study the $\infty$-category 
\[
\Shv \left( \I^\semi_{L \subset J} \backslash X \right).
%, \quad \Shv(\I_J^{\quar}\backslash X) \quad \text{and}\quad \Shv(\I_J^{\semi}\backslash X).
\]
These $\infty$-categories (for the various choices of $L$, $J$ and $X$) admit natural descriptions in terms of limits/colimits of $\infty$-cate\-gories of sheaves on algebraic stacks of finite type, which can be obtained by adapting (in the obvious way) the discussions in~\cite[\S\S 2.5--2.6]{adr} or~\cite[\S 1.2]{Gaits18}.

%For any $x\in W_\ex^K$, resp. $x\in {}^JW_\ex^K$, we have 
%\[
%\Shv(\I_J^{\quar}\backslash {}^J\rQ^K_x) \simeq \Vect_\k, \quad \text{resp.}\quad \Shv(\I_J^{\semi}\backslash {}^J\S^\semi_x) \simeq ?. 
%\]

If $L_1 \subset L_2 \subset J$,
since $\I_{L_1 \subset J}^{\semi}$ is a subgroup of $\I_{L_2 \subset J}^{\semi}$, we
have a natural forgetful functor
\[
\oblv^{\I_{L_2 \subset J}^{\semi}|\I_{L_1 \subset J}^{\semi}}: \Shv \left( \I_{L_2 \subset J}^{\semi} \backslash X \right) \rightarrow \Shv \left( \I_{L_1 \subset J}^{\semi} \backslash X \right)
\]
induced by $!$-pullback along the quotient map $\I_{L_1 \subset J}^{\semi} \backslash X \to \I_{L_2 \subset J}^{\semi}\backslash X$,
which admits left and right adjoints 
\[
\Av_!^{\I_{L_2 \subset J}^{\semi}|\I_{L_1 \subset J}^{\semi}}, \Av_*^{\I_{L_2 \subset J}^{\semi}|\I_{L_1 \subset J}^{\semi}} : \Shv \left( \I_{L_1 \subset J}^{\semi} \backslash X \right) \rightarrow \Shv \left( \I_{L_2 \subset J}^{\semi} \backslash X \right). 
\]
Here $\Av_!^{\I_{L_2 \subset J}^{\semi}|\I_{L_1 \subset J}^{\semi}}$ is given by $!$-pushforward along the map $\I_{L_1 \subset J}^{\semi} \backslash X \to \I_{L_2 \subset J}^{\semi}\backslash X$, and $\Av_*^{\I_{L_2 \subset J}^{\semi}|\I_{L_1 \subset J}^{\semi}}$ is given by $*$-pushforward along this map, followed by $[-2\ell(w_{L_2}w_{L_1})]$. In fact, since the projection map is proper we have
\[
\Av_*^{\I_{L_2 \subset J}^{\semi}|\I_{L_1 \subset J}^{\semi}} = \Av_!^{\I_{L_2 \subset J}^{\semi}|\I_{L_1 \subset J}^{\semi}}[-2\ell(w_{L_2}w_{L_1})].
\]

We also have pullback and pushforward functors 
\begin{gather*}
(p_K)^*, (p_K)^!: \Shv \left( \I^\semi_{L \subset J} \backslash \Fl_K \right) \rightarrow \Shv \left(\I^\semi_{L \subset J}  \backslash \Fl \right), \\ 
(p_K)_* = (p_K)_!: \Shv \left( \I^\semi_{L \subset J} \backslash \Fl \right) \rightarrow \Shv \left( \I^\semi_{L \subset J}  \backslash \Fl_K \right), \\
%\]
%\[
%\text{and}\quad 
q^*, q^!: \Shv \left( \I^\semi_{L \subset J} \backslash \Fl \right) \rightarrow \Shv \left( \I^\semi_{L \subset J} \backslash \tFl \right), \ 
q_*,q_!: \Shv \left( \I^\semi_{L \subset J}  \backslash \tFl \right) \rightarrow \Shv \left( \I^\semi_{L \subset J} \backslash \Fl \right). 
\end{gather*}
%Denote by $n_K$ the dimension of the flag variety of $M_K$. Then
We have $(p_K)^! \simeq (p_K)^*[2 \ell(w_K)]$, and we set
\[
(p_K)^\dag = (p_K)^*[\ell(w_K)].
\]
Similarly we have $q^! \simeq q^*[2\dim(T)]$.
%, and we set
%\[
%q^\dag = q^*[\dim(T)].
%\]

%---------------------------------------------------------------
\subsection{Standard and costandard sheaves} 
%---------------------------------------------------------------

For $L$, $J$, $K$ as above and $x \in {}^LW_\ex^{(K,J)}$,
we will also consider the $\infty$-category 
\begin{equation}
\label{eq:cat-orbt}
\Shv \left( \I^{\semi}_{L \subset J} \backslash \rS^{L \subset J}_{K,x} \right).
\end{equation}
As in~\cite[Proposition~1.3.5]{Gaits18} or \cite[Lemma~3.3.2]{DL23}, $!$-restriction to the point $\dot{x}\I_K/\I_K$ induces an equivalence 
\begin{equation}
\label{eq:semi-orbt}
\Shv \left( \I^{\semi}_{L \subset J} \backslash \rS^{L \subset J}_{K,x} \right) \xs \Shv \left( \mathrm{Stab}_{M_L}(\dot{x}\I_K/\I_K)\backslash \pt \right).
\end{equation}
To simplify notation we will write $\omega_{\rS^{L \subset J}_{K,x}}$ for $\omega_{\I^{\semi}_{L \subset J} \backslash \rS^{L \subset J}_{K,x}}$. (We will use similar conventions whenever convenient below.) Under the equivalence~\eqref{eq:semi-orbt}, this object corresponds to $\omega_{\mathrm{Stab}_{M_L}(\dot{x}\I_K/\I_K)\backslash \pt}$.
%, the object in LHS corresponding to the trivial representation $\k$ in RHS under the equivalence \eqref{eq:quar-orbt}, \eqref{eq:semi-orbt}, \eqref{eq:fl-orbt}, respectively. 
%For $x\in {}^JW_\ex^{(K,J)}$, resp. $x\in {}^JW_\ex^{K}$, we set 
%\[
%\ell^J_K(x)=\max\{\ell^J(y)\}_{y\in W_Jx\cap W_\ex^{(K,J)}}, \quad \text{resp.}\quad 
%\ell_K(x)=\max\{\ell(y)\}_{y\in W_Jx\cap W_\ex^{K}}. 
%\]
We define the \textit{perverse t-structure} on~\eqref{eq:cat-orbt} to be the cohomological shift of the transfer of the natural t-structure on the $\infty$-category $\Shv \left( \mathrm{Stab}_{M_L}(\dot{x}\I_K/\I_K)\backslash \pt \right)$ such that the object
\begin{equation}
\label{eqn:omega-orbit-perverse}
\omega_{\rS^{L \subset J}_{K,x}}[-\ell(w_Lw_{L\cap ({}^xK)})-\ell^J(x)]
\end{equation}
%belong to the heart in the corresponding category, which is called the \textit{perverse t-structure}. 
belongs to the heart. 
Note that when $J=S$, the forgetful functor $\Shv(\I_L \backslash \rS^{L \subset S}_{K,x}) \to \Shv(\rS^{L \subset S}_{K,x})$ is t-exact with respect to the usual perverse t-structure on $\Shv(\rS^{L \subset S}_{K,x})$, see~\eqref{eqn:dim-par-orbit}. In general we are dealing with schemes of infinite dimension, so that there is no natural choice for the shift appearing in~\eqref{eqn:omega-orbit-perverse}; as in~\cite{adr} the shifts are chosen so as to provide a nice theory. (In particular, when $L=J=\varnothing$, our conventions are consistent with those in~\cite{adr}.)

We set
\begin{align*}
\Delta^{L \subset J,\semi}_{K,x} &:= (\bs^{L \subset J}_{K,x})_! \omega_{\rS^{L \subset J}_{K,x}}[-\ell(w_L w_{L \cap ({}^xK)})-\ell^J(x)], \\ 
\nabla^{L \subset J,\semi}_{K,x} &:= (\bs^{L \subset J}_{K,x})_* \omega_{\rS^{L \subset J}_{K,x}}[-\ell(w_L w_{L \cap ({}^xK)})-\ell^J(x)].
\end{align*}
These objects will be called the standard and costandard sheaves in $\Shv \left( \I^{\semi}_{L \subset J} \backslash \Fl_{K} \right)$ respectively.
We will omit $K$ in the notation for these objects if it is the empty set, and will also omit ``$L \subset J$'' in case $L=J=\varnothing$ (which corresponds to the setting of~\cite{adr}). In case $L=\varnothing$, we will also replace  ``$\varnothing \subset J, \semi$'' by ``$J,\quar$,'' and when $L=J$ we will replace ``$J \subset J, \semi$'' by ``$J,\semi$.'' Finally, when $J=S$ (so that $\I_{L \subset S}^\semi = \I_L$) we will replace ``$L \subset S, \semi$'' by ``$L$'', and will omit $L$ if $L=\varnothing$ (so that we work with usual Iwahori-equivariant sheaves on $\Fl_K$). We will use similar conventions for the variants of these objects introduced below.

\begin{lem}
\label{lem:Av-quar-semi-D-N}
Let $L \subset J \subset S$, and let $K \subset S_\af$ be finitary.
Let $x\in {}^LW_\ex^{(K,J)}$. 
\begin{enumerate}
\item 
\label{it:Av-quar-semi-D-N-1}
Let $w\in W_L^{L\cap ({}^x K)}$. 
The object 
\[
\Av_!^{\I_{L \subset J}^{\semi}|\I_J^{\quar}}(\Delta^{J,\quar}_{K,wx}), \quad 
\text{resp.}\quad \Av_*^{\I_{L \subset J}^{\semi}|\I_J^{\quar}}(\nabla^{J,\quar}_{K,wx}),
\]
is an iterated extension of objects
\[
\Delta^{L \subset J,\semi}_{K,x}[\ell(w_{L\cap ({}^xK)}u)+\ell(w_Lwu)], \
\text{resp.}\ \nabla^{L \subset J,\semi}_{K,x}[-\ell(w_{L\cap ({}^xK)}u)-\ell(w_Lwu)],
\]
with $u\in W_{L\cap ({}^xK)}$, each of them appearing once. 
\item 
\label{it:Av-quar-semi-D-N-2}
The object 
\[
\oblv^{\I_{L \subset J}^{\semi}|\I_J^{\quar}}(\Delta^{L \subset J,\semi}_{K,x}), 
\quad \text{resp.}\quad 
\oblv^{\I_{L \subset J}^{\semi}|\I_J^{\quar}}(\nabla^{L \subset J,\semi}_{K,x}),
\]
is an iterated extension of objects
\[
\Delta^{J,\quar}_{K,wx}[\ell(w_L w_{L\cap ({}^xK)})-\ell(w)], 
\ \text{resp.}\ 
\nabla^{J,\quar}_{K,wx}[-\ell(w_L w_{L \cap ({}^xK)})+\ell(w)],
\]
with $w\in W_L^{L\cap ({}^xK)}$, each of them appearing once. 
\end{enumerate}
\end{lem}

\begin{proof}
\eqref{it:Av-quar-semi-D-N-1}
Lemma~\ref{lem:bi-min-av}\eqref{it:bi-min-av-1} allows to reduce the proof to the case $w=1$. Then we use Lemma~\ref{lem:bi-min-av}\eqref{it:bi-min-av-2}. For the first claim we have to compute the $!$-pushforward under that map of $\omega_{\I_{L \subset J}^\semi \times^{\I_J^\quar} \rS_{K,x}^{J}}$, and the statement follows from the fact that the cohomology of $\omega_{M_{L \cap ({}^x K)}/B_{L \cap ({}^x K)}}$ is $\bigoplus_{u \in W_{L \cap ({}^x K)}} \k[2\ell(w_{L \cap ({}^x K)}) - 2\ell(u)]$. For the second claim we use the fact that $\Av_*^{\I_{L \subset J}^{\semi}|\I_J^{\quar}}$ is the composition of the appropriate $*$-pushforward functor with 
%= \Av_!^{\I_{L \subset J}^{\semi}|\I_J^{\quar}}
$[-2\ell(w_L)]$ (see~\S\ref{ss:qs-sheaves}), and write now the cohomology above as $\bigoplus_{u \in W_{L \cap ({}^x K)}} \k[2\ell(u)]$.

\eqref{it:Av-quar-semi-D-N-2}
Using the decomposition from Lemma~\ref{lem:bi-min-av}\eqref{it:bi-min-av-3}, we obtain that $\omega_{\rS^{L \subset J}_{K,x}}$ is an iterated extension of objects $f_! f^* \omega_{\rS^{L \subset J}_{K,x}}$, resp.~$f_* f^! \omega_{\rS^{L \subset J}_{K,x}}$, where $f$ runs over the locally closed immersions $\rS^{J}_{K,wx} \hookrightarrow \rS^{L \subset J}_{K,x}$ ($w \in W_L^{L \cap ({}^xK)}$). The second assertion immediately follows. For the first one, we need to compute the complexes $f^* \omega_{\rS^{L \subset J}_{K,x}}$ for $f$ as above.
Using Lemma~\ref{lem:bi-min-av}\eqref{it:bi-min-av-1},
we write the embedding $\rS^{J}_{K,wx} \hookrightarrow \rS^{L \subset J}_{K,x}$ as the composition
\[
\rS^{J}_{K,wx} \hookrightarrow \I_{L \subset J}^\semi \times^{\I_J^\quar} \rS_{K,x}^J \to \rS^{L \subset J}_{K,x}.
\]
The second map is smooth of relative dimension $\ell(w_{L \cap ({}^x K)})$ by Lemma~\ref{lem:bi-min-av}\eqref{it:bi-min-av-2}, so the $*$-pullback under this map of $\omega_{\rS^{L \subset J}_{K,x}}$ is $\omega_{\I_{L \subset J}^\semi \times^{\I_J^\quar} \rS_{K,x}^J}[-2\ell(w_{L \cap ({}^x K)})]$. On the other hand, the $*$-pullback under the first map of $\omega_{\I_{L \subset J}^\semi \times^{\I_J^\quar} \rS_{K,x}^J}$ is $\omega_{\rS^{J}_{K,wx}}[2\ell(w_L)-2\ell(w)]$. The claim follows.
%It follows from a standard computation using 
%This follows from Lemma~\ref{lem:bi-min-av} and standard arguments. 
\end{proof}

\begin{lem}
\label{lem:push-pull-pK-D-N-semi}
Let $L \subset J \subset S$, and let $K \subset S_\af$ be finitary.
Let $x\in {}^L W_\ex^{(K,J)}$. 
\begin{enumerate}
\item 
\label{it:push-pull-pK-D-N-semi-1}
Let $w\in {}^{(L^x)\cap K}W_K$. 
The object 
\[
(p_K)_*(\Delta^{L \subset J,\semi}_{xw})[\ell(w_K)], 
\quad \text{resp.}\quad 
(p_K)_*(\nabla^{L \subset J, \semi}_{xw})[-\ell(w_K)], 
\]
is an iterated extension of objects
\[
\Delta^{L \subset J,\semi}_{K,x}[\ell(uw_{(L^x)\cap K})+\ell(uww_K)], \
\text{resp.}\ \nabla^{L \subset J,\semi}_{K,x}[-\ell(uw_{(L^x)\cap K})-\ell(uww_K)],
\]
with $u\in W_{(L^x)\cap K}$, each appearing once. 
\item
\label{it:push-pull-pK-D-N-semi-2}
The object 
\[
(p_K)^\dag(\Delta^{L \subset J,\semi}_{K,x}), 
\quad \text{resp.}\quad 
(p_K)^\dag(\nabla^{L \subset J,\semi}_{K,x}), 
\]
is an iterated extension of objects
\[
\Delta^{L \subset J, \semi}_{xw}[\ell(w_{(L^x)\cap K}w_K)-\ell(w)], \
\text{resp.}\ \nabla^{L \subset J, \semi}_{xw}[-\ell(w_{(L^x)\cap K}w_K)+\ell(w)],
\]
with $w\in {}^{(L^x)\cap K}W_K$, each appearing once. 
\end{enumerate}
\end{lem}

\begin{proof}
\eqref{it:push-pull-pK-D-N-semi-1}
Using Lemma~\ref{lem:bi-min-push}\eqref{it:bi-min-push-1}, we reduce the proof to the case $w=1$. In that case we use Lemma~\ref{lem:bi-min-push}\eqref{it:bi-min-push-2} to see that $(p_K)_*(\Delta^{L \subset J,\semi}_{x})$ is an iterated extension of objects $\Delta^{L \subset J,\semi}_{K,x}[-\ell(w_{L \cap ({}^x K)}) + 2\ell(w_{(L^x) \cap K})-2\ell(u)]$ with $u\in W_{(L^x)\cap K}$. 
The first assertion follows, using the fact that $\ell(w_{L \cap ({}^x K)}) = \ell(w_{(L^x) \cap K})$. 
The second assertion follows from the same considerations (and replacing $u$ by $u w_{(L^x) \cap K}$ in the formula above).

\eqref{it:push-pull-pK-D-N-semi-2}
By the base change theorem and smoothness of $p_K$, the complex $(p_K)^\dag(\Delta^{L \subset J,\semi}_{K,x})$, resp.~$(p_K)^\dag(\nabla^{L \subset J,\semi}_{K,x})$, is the $!$-pushforward, resp.~$*$-pushforward, of
\[
\omega_{(p_K)^{-1}(\rS^{L \subset J}_{K,x})}[-\ell(w_L)+\ell(w_{L \cap ({}^x K)})-\ell^J(x)-\ell(w_K)]
\]
along the immersion $(p_K)^{-1}(\rS^{L \subset J}_{K,x}) \hookrightarrow \Fl$. Now, by the decomposition in Lem\-ma~\ref{lem:bi-min-push}\eqref{it:bi-min-push-3}, $\omega_{(p_K)^{-1}(\rS^{L \subset J}_{K,x})}$ is an iterated extension of $*$-pushforwards of complexes $\omega_{\rS^{L \subset J}_{xw}}$ with $w \in {}^{(L^x) \cap K} W_K$, which implies the second assertion using the formula $\ell(w_{L \cap ({}^x K)}) = \ell(w_{(L^x) \cap K})$. Similarly, $\omega_{(p_K)^{-1}(\rS^{L \subset J}_{K,x})}$ is an iterated extension of objects $f_! f^* \omega_{(p_K)^{-1}(\rS^{L \subset J}_{K,x})}$ where $f$ runs over the immersions $\rS^{L \subset J}_{xw} \hookrightarrow (p_K)^{-1}(\rS^{L \subset J}_{K,x})$, and arguments similar to those encountered in the proof of Lemma~\ref{lem:Av-quar-semi-D-N}\eqref{it:Av-quar-semi-D-N-2}
imply the first assertion.
\end{proof}

%---------------------------------------------------------------
\subsection{Whittaker sheaves} 
\label{ss:Whittaker}
%---------------------------------------------------------------

In this subsection we assume that $p>0$, and that $\k$ contains a nontrivial $p$-th root of unity.

%===============================
\subsubsection{Definition} 
\label{sss:Whittaker}
%===============================

%Let us
Let $A\subset S_\af$ be a finitary subset.
%\footnote{One should be careful that the subset ``$K$'' used here in the construction of Whittaker sheaves is unrelated to the subsets also denoted $K$ above.} 
As in~\cite[\S 3.4]{AR23}, for any $s \in S_\af$ we have an associated (negative, affine) root subgroup $U_s^- \subset \sL G$ and a (noncanonical, but fixed) isomorphism $\Ga \simeq U_s^-$.
%Recall that $w_K$ is the longest element in $W_K$, and $\dot{w}_K$ is a lift of $w_K$ in $G$. 
Recall from~\S\ref{sss:def-Fl} that we have a parahoric subgroup $\I_A \subset \sL G$ associated with $A$, with pro-unipotent radical $\Iu \cap \Iu^A$ and Levi factor $M_A$, which therefore identifies with the quotient $\I_A / (\Iu \cap \Iu^A)$. The image $U_A^-$ of $\Iu^A$ in this quotient is the unipotent radical of a Borel subgroup, and as in~\cite[\S 3.4]{AR23} we consider the composition
\[
\Iu^A \to U_A^- \to U_A^- / [U_A^-, U_A^-] \simeq \prod_{s \in A} U_s^- \simeq (\Ga)^{\# A} \xrightarrow{\mathrm{sum}} \Ga
\]
where the first two maps are the natural quotient morphisms.
%obtained by projection to the quotient of $U_K^-$ by its derived subgroup, which identifies (after some choices) with a product of $\# K$ copies of the additive group, followed by the sum morphism $(\Ga)^{\# K} \to \Ga$.
%We set ${}^KU=\Ad_{\dot{w}_K}(U)$ and ${}^K\Iu=\Ad_{\dot{w}_K}(\Iu)$, and consider the morphisms of algebraic groups 
%\[
%{}^K\Iu\twoheadrightarrow {}^KU \twoheadrightarrow \prod_{\alpha\in -\Phis_K} U_\alpha \simeq \prod_{\alpha\in -\Phis_K} \Ga \xrightarrow{\text{sum}} \Ga.
%\]
%(Here, the first morphism is induced by the projection $\Iu \to U$, and for the second one we consider the decomposition of ${}^K U$ as a product of root subgroups, and the projection to those corresponding to roots on $-\Phis_K$.)
Under our assumptions on $p$ and $\k$ the local system $\mathrm{AS}_* \underline{\k}_{\Ga}$, where $\mathrm{AS} : \Ga \to \Ga$ is the Artin--Schreier map sending $x$ to $x^p-x$, splits as a direct sum of $p$ rank-$1$ local systems, parametrized by morphisms from $(\mathbb{F}_p,+)$ to $\k^\times$. We choose one nontrivial irreducible direct summand,
%We fix an Artin--Schreier local system on $\Ga$ (i.e.~a nonconstant irreducible direct summand of the local system $\mathrm{AS}_* \underline{\k}_{\Ga}$, where $\mathrm{AS} : \Ga \to \Ga$ is the Artin--Schreier map sending $x$ to $x^p-x$), 
and denote by $\chi_A$, resp.~$\chi'_A$, its pullback to $\Iu^A$, resp.~to $U_A^-$. 

Combining the constructions from~\S\ref{ss:qs-sheaves} with the consideration of Whittaker shea\-ves as in~\cite[End of~\S 4.2]{zhu},
one can define for any subsets $L \subset J \subset S$ the $\infty$-category
\begin{equation}
\label{eqn:Whit-cat}
%\Shv(\I_J\backslash \sL G/({}^K\Iu,\chi_K)), \ \ \Shv(\I_J^{\quar}\backslash \sL G/({}^K\Iu,\chi_K)) \ \ \text{and}\ \ \Shv(\I_J^{\semi}\backslash \sL G/({}^K\Iu,\chi_K)). 
\Shv \left( \I_{L \subset J}^{\semi}\backslash \sL G/(\Iu^A,\chi_A) \right). 
\end{equation}
%Let us explain the definition of $\Shv(\I_J^{\quar}\backslash \sL G/({}^K\Iu,\chi_K))$ as an example. 
%Let $\theta\in \Xd_J^0 \cap \Xd_{J^\comp}^{++}$.
%% be $J^c$-dominant enough. 
%We set 
%\[
%\I_{J,(n)}^{\quar}=\I_{M_J}\ltimes \Ad_{z^{-n\theta}}(\sL^+U_{J}),
%%=\Ad_{z^{-n\theta}}(\I_{M_J}\ltimes \sL^+U_{J})
%\]
%then $\{\I_{J,(n)}^{\quar}\}_n$ forms an exhaustive filtration of $\I_{J}^{\quar}$. 
%We set $\widetilde{\Gr}=\sL G/\sL^{++}G$, which admits an $\sL G$-action from the left and a ${}^KU$-action from the right. 
%For any $n$ one can define the category 
%\[
%\Shv(\I_{J,(n)}^{\quar}\backslash \sL G/ ({}^K\Iu,\chi_K))
%=\Shv\big( \I_{J,(n)}^{\quar} \backslash \widetilde{\Gr} / ({}^KU,\chi'_K) \big).
%\]
%Then we define 
%\[
%\Shv(\I_J^{\quar}\backslash \sL G/({}^K\Iu,\chi_K)) = \varprojlim_{n\geq 0} \Shv(\I_{J,(n)}^{\quar}\backslash \sL G/ ({}^K\Iu,\chi_K)), 
%\]
%where the transition maps are given by the functor $\oblv^{\I_{J,(n+1)}^{\quar}\mid \I_{J,(n)}^{\quar}}$. 
More specifically, let $\widetilde{\Iu^A} := \Iu^A \times_{\Ga} \Ga$, where the map $\Ga \to \Ga$ is $\mathrm{AS}$. Then one can consider the stack $\I_{L \subset J}^\semi \backslash \sL G / \widetilde{\Iu^A}$. The $\infty$-category
\[
\Shv \left( \I_{L \subset J}^\semi \backslash \sL G / \widetilde{\Iu^A} \right)
\]
breaks as a direct sum indexed by the $p$ different irreducible direct summands of $\mathrm{AS}_* \underline{\k}_{\Ga}$, and $\Shv \left( \I_{L \subset J}^{\semi}\backslash \sL G/(\Iu^A,\chi_A) \right)$ is defined as the direct summand corresponding to our choice of Artin--Schreier local system. As in~\S\ref{ss:qs-sheaves} we have functors
\[
\oblv^{\I_{L_2 \subset J}^{\semi}|\I_{L_1 \subset J}^{\semi}}, \Av_!^{\I_{L_2 \subset J}^{\semi}|\I_{L_1 \subset J}^{\semi}}, \Av_*^{\I_{L_2 \subset J}^{\semi}|\I_{L_1 \subset J}^{\semi}}
\]
relating the $\infty$-categories~\eqref{eqn:Whit-cat} when $L$ varies.

It is convenient to introduce the ind-scheme 
\[
\uFl_A := \sL G / (\Iu \cap \Iu^A),
\]
which admits an $\sL G$-action from the left and an $M_A$-action from the right; then we have
\[
\Shv \left( \I_{L \subset J}^{\semi}\backslash \sL G/(\Iu^A,\chi_A) \right) \simeq \Shv \left( \I_{L \subset J}^{\semi} \backslash \uFl_A /(U_A^-,\chi'_A) \right)
\]
where the right-hand side is defined as above.
%and similarly for the other variants.
Denote by $r_A: \uFl_A \rightarrow \Fl$ the natural projection. 
By~\eqref{eq:decomp} we have a decomposition
%of $\widetilde{\Gr}$ into $(\I_{J}^{\quar}\times {}^KU)$-, $(\I_{J}^{\semi}\times {}^KU)$-, and $(\I_{J}\times {}^KU)$-orbits 
\begin{equation}
\label{eq:decomp-tGr}
\uFl_A= \bigsqcup_{x\in {}^L W_\ex} (r_A)^{-1}(\rS^{L \subset J}_{x}) \cdot \dot{w}_A
%\widetilde{\Gr}= \bigsqcup_{x\in W_\ex} \widetilde{q}^{-1}({}^{J}\rQ_x)\dot{w}_K
%= \bigsqcup_{x\in {}^JW_\ex} \widetilde{q}^{-1}({}^{J}\rS_x)\dot{w}_K 
%= \bigsqcup_{x\in {}^JW_\ex} \widetilde{q}^{-1}({}^{J}\Fl_x)\dot{w}_K. 
\end{equation}
into $(\I_{L \subset J}^\semi \times (U_A^-)^\mathrm{op})$-orbits.

\begin{lem}
\label{lem:whit}
For $x\in {}^L W_\ex$, $!$-restriction to the point $\dot{x}\dot{w}_A$ induces an equivalence 
\begin{equation*}
%\label{eq:semi-orbt-whit}
\Shv\left( \I_{L \subset J}^{\semi} \backslash (r_A)^{-1}(\rS^{L \subset J}_x)\dot{w}_A / (U_A^-,\chi'_A) \right)
\simeq
\begin{cases}
\Vect_\k & \text{if $x\in {}^L W_\ex^{(A,J),\max}$;} 
\\ 0 & \text{otherwise.} 
\end{cases}
\end{equation*}
\end{lem}

\begin{proof}
This follows from standard arguments; namely, the same arguments as in~\cite[Proposition~1.3.5]{Gaits18} allow to reduce the proof to a finite-dimensional setting, and then one uses the same arguments as in~\cite[Lemma~4.2.1]{by}. In the end we use Lemma~\ref{lem:quar-order} and Lemma~\ref{lem:quar-min-max}, which ensure that $x \in {}^L W_\ex^{(A,J),\max}$ if and only if we have $\Ad_{\dot{x}\dot{w}_A}(U_s^-) \cap \I_{L \subset J}^{\semi} = \{1\}$ for any $s \in A$.
%We only explain the proof of~\eqref{it:lem-whit-1}; those of the other statements are similar. 
%Let $\theta\in \Xd_J^0 \cap \Xd_{J^\comp}^{++}$, and for $n \geq 0$ set 
%\[
%\I_{L \subset J,(n)}^{\semi}=\I_{L \subset J}\ltimes \Ad_{z^{-n\theta}}(\sL^+U_{J});
%\]
%then $( \I_{L \subset J,(n)}^{\semi} : n \geq 0)$ forms an exhaustive filtration of $\I_{L \subset J}^{\semi}$. 
%The $\infty$-category in the left-hand side of~\eqref{eq:quar-orbt-whit} is equivalent to the limit $\lim_{n} \sC_n$, where 
%\[
%\sC_n:= \Shv\big(\I_{J,(n)}^{\quar} \backslash (\I_{J,(n)}^{\quar}\dot{x}\dot{w}_K \sL^{++}G/\sL^{++}G) / ({}^KU,\chi'_K) \big).
%\]
%It is therefore enough to show that for $n \gg 0$, $!$-restriction to the point $\dot{x}\dot{w}_K$ induces an equivalence 
%\begin{equation}
%\label{eq:Cn}
%\sC_n \simeq
%\begin{cases}
%\Vect_\k & \text{if $x\in W_\ex^{(K,J),\mathrm{w}}$;} 
%\\ 0 & \text{otherwise.} 
%\end{cases}
%\end{equation}
%By Lemma~\ref{lem:min-length}\eqref{it:min-length-2}, $x\in W_\ex^{(K,J),\max}$, equivalently $xw_K\in W_\ex^{(K,J)}$, if and only if $xw_K\preceq^{J}xw_K s_\alpha$ for any $\alpha\in \Phis_K$, which by Lemma~\ref{lem:quar-order} is equivalent to the condition that 
%\[
%\Ad_{\dot{x}\dot{w}_K}(U_{-\alpha}) \cap \I_J^{\quar} = \{1\}, \quad \forall \alpha\in \Phis_K. 
%\]
%In this case $\Ad_{\dot{x}\dot{w}_K}(U_{-\alpha}) \cap \I_{J,(n)}^{\quar} = \{1\}$ for any $n\in \Z$. 
%Hence we obtain \eqref{eq:Cn}, provided $n \gg 0$. 
\end{proof}

If $x \in {}^L W_\ex^{(A,J),\max}$, we will denote by 
\[
\omega_{(r_A)^{-1}(\rS^{L \subset J}_x)\dot{w}_A}^{\chi_A} \in \Shv\left( \I_{L \subset J}^{\semi} \backslash (r_A)^{-1}(\rS^{L \subset J}_x)\dot{w}_A / (U_A^-,\chi'_A) \right)
\]
the object corresponding to $\k\in \Vect_\k$ under the equivalence of Lemma~\ref{lem:whit}.
% under the equivalence \eqref{eq:quar-orbt-whit}. 
Let $\widetilde{\bs}^{L \subset J}_{A,x}$ be the embedding of $(r_A)^{-1}(\rS^{L \subset J}_x)\dot{w}_A$ in $\uFl_A$. 
We define the standard and costandard sheaves in $\Shv \left( \I_{L \subset J}^{\semi}\backslash \sL G/(\Iu^A,\chi_A) \right)$ as 
\begin{align*}
\Delta_{\chi_A,x}^{L \subset J, \semi} &:= (\widetilde{\bs}^{L \subset J}_{A,x})_! \omega_{(r_A)^{-1}(\rS^{L \subset J}_x)\dot{w}_A}^{\chi_A} [-\ell^J(x)-\ell(w_L)],
\\
\nabla_{\chi_A,x}^{L \subset J, \semi} &:= (\widetilde{\bs}^{L \subset J}_{A,x})_* \omega_{(r_A)^{-1}(\rS^{L \subset J}_x)\dot{w}_A}^{\chi_A} [-\ell^J(x)-\ell(w_L)].
\end{align*}
%(Note that here we have $L \cap ({}^x K) = \varnothing$ by Remark~\ref{rmk:intersection-max-elt}, so the shift follows the same pattern as in the non-Whittaker case.)
%As above, we will replace ``$J \subset J, \semi$'' by ``$J,\semi$'' or ``$\varnothing \subset J, \semi$'' by ``$J, \quar$'' or ``$J \subset S, \semi$'' by ``$J$'' in this notation whenever convenient.

It follows from Lemma~\ref{lem:quar-min-max} that if $x\in W_\ex^{(A,J),\max}$, there exists at most one element $w \in W_L$ such that $wx \in {}^L W_\ex^{(A,J),\max}$; on the other hand, if $x\in {}^L W_\ex^{(A,J),\max}$, then $wx \in W_\ex^{(A,J),\max}$ for any $w \in W_L$.

\begin{lem}
\label{lem:av-oblv-whit}
Let $L \subset J \subset S$, and let $A \subset S_\af$ be a finitary subset.
\begin{enumerate}
\item 
\label{it:av-oblv-whit-1}
Let $x\in W_\ex^{(A,J),\max}$. 
We have 
\[
\Av_!^{\I_{L \subset J}^{\semi}|\I_J^{\quar}}(\Delta^{J,\quar}_{\chi_A,x})\simeq
\begin{cases}
\Delta^{L \subset J,\semi}_{\chi_A, wx}[\ell(w_Lw)] & \text{if $w\in W_L$ satisfies $wx\in {}^LW_\ex^{(A,J),\max}$;} \\ 
0 & \text{if $W_L x \cap {}^LW_\ex^{(A,J),\max} = \varnothing$,}
\end{cases}
\]
and 
\[
\Av_*^{\I_{L \subset J}^{\semi}|\I_J^{\quar}}(\nabla^{J,\quar}_{\chi_A,x})\simeq
\begin{cases}
\nabla^{L \subset J,\semi}_{\chi_A, wx}[-\ell(w_Lw)] & \text{if $w\in W_L$ satisfies $wx\in {}^LW_\ex^{(A,J),\max}$;} \\ 
0 & \text{if $W_L x \cap {}^LW_\ex^{(A,J),\max} = \varnothing$.} 
\end{cases}
\]
\item 
\label{it:av-oblv-whit-2}
Let $x\in {}^L W_\ex^{(A,J),\max}$. 
Then $\oblv^{\I_{L \subset J}^{\semi}|\I_J^{\quar}}(\Delta^{L \subset J,\semi}_{\chi_A,x})$, resp. $\oblv^{\I_{L \subset J}^{\semi}|\I_J^{\quar}}(\nabla^{L \subset J,\semi}_{\chi_A,x})$, is an iterated extension of objects $\Delta^{J,\quar}_{\chi_A, wx}[\ell(w_Lw)]$, resp. $\nabla^{J,\quar}_{\chi_A,wx}[-\ell(w_Lw)]$, with $w\in W_L$, each appearing once. 
\end{enumerate}
\end{lem}

\begin{proof}
\eqref{it:av-oblv-whit-1}
We only consider the first assertion; the second one can be justified similarly.
The case when $W_L x \cap {}^L W_\ex^{(A,J),\max}=\varnothing$ directly follows from Lemma~\ref{lem:whit}. Now, let us assume that there exists $w\in W_L$ such that $wx\in {}^L W_\ex^{(A,J),\max}$. Using Lemma~\ref{lem:bi-min-av}\eqref{it:bi-min-av-1} we obtain that
\[
\Av_!^{\I_{L \subset J}^{\semi}|\I_J^{\quar}}(\Delta^{J,\quar}_{\chi_A,x}) \simeq \Av_!^{\I_{L \subset J}^{\semi}|\I_J^{\quar}}(\Delta^{J,\quar}_{\chi_A,wx})[-\ell(w)],
\]
so we can assume $w=1$. In this case, the claim follows from Lemma~\ref{lem:bi-min-av}\eqref{it:bi-min-av-2} (applied with $K=\varnothing$, so that the flag variety of the Levi subgroup under consideration is a point).

\eqref{it:av-oblv-whit-2}
The proof is the same as that of Lemma~\ref{lem:Av-quar-semi-D-N}\eqref{it:Av-quar-semi-D-N-2}, using
%This follows from 
Lemma~\ref{lem:bi-min-av}\eqref{it:bi-min-av-3}.
\end{proof}

%===============================
\subsubsection{Whittaker averaging functors} 
\label{sss:av-fun}
%===============================

%Let $A=\I_J, \I_J^{\quar}$ or $\I_J^{\semi}$. 
%We continue to assume that $p>0$ and that $\k$ contains a nontrivial $p$-th root of unity, and
Let $L$, $J$, $A$ be as above.
We have a forgetful functor 
\[
\oblv^{(\Iu^A,\chi_A)\mid \Iu^A \cap \Iu}: 
\Shv \left( \I_{L \subset J}^{\semi}\backslash \sL G/(\Iu^A,\chi_A) \right) \rightarrow \Shv \left( \I_{L \subset J}^{\semi}\backslash \sL G/(\Iu^A \cap \Iu) \right)
\]
which admits left and right adjoints 
\begin{multline*}
\Av^{(\Iu^A,\chi_A)\mid \Iu^A \cap \Iu}_!, \Av^{(\Iu^A,\chi_A)\mid \Iu^A \cap \Iu}_* : \Shv \left( \I_{L \subset J}^{\semi}\backslash \sL G/(\Iu^A \cap \Iu) \right) \\
\to \Shv \left( \I_{L \subset J}^{\semi}\backslash \sL G/(\Iu^A,\chi_A) \right).
\end{multline*}
Similarly we have the forgetful functor 
\[
\oblv^{\Iu\mid \Iu^A \cap \Iu}: \Shv \left( \I_{L \subset J}^{\semi} \backslash \tFl \right) \rightarrow \Shv \left( \I_{L \subset J}^{\semi}\backslash \sL G/(\Iu^A \cap \Iu) \right),
\]
and its left and right adjoints 
\[
\Av^{\Iu\mid \Iu^A \cap \Iu}_!, \Av^{\Iu\mid \Iu^A \cap \Iu}_* : \Shv \left( \I_{L \subset J}^{\semi}\backslash \sL G/(\Iu^A \cap \Iu) \right) \to \Shv \left( \I_{L \subset J}^{\semi} \backslash \tFl \right).
\]
We set 
\begin{align*}
\Av^A_{\chi,!} &:= \Av^{(\Iu^A,\chi_A)\mid \Iu^A \cap \Iu}_! \circ \oblv^{\Iu\mid \Iu^A \cap \Iu}[-\ell(w_A)], \\ 
\Av^A_{\chi,*} &:= \Av^{(\Iu^A,\chi_A)\mid \Iu^A \cap \Iu}_* \circ \oblv^{\Iu\mid \Iu^A \cap \Iu}[\ell(w_A)], \\
\Av^A_{!} &:= \Av^{\Iu\mid \Iu^A \cap \Iu}_! \circ \oblv^{(\Iu^A,\chi_A)\mid \Iu^A \cap \Iu}[-\ell(w_A)], \\ 
\Av^A_{*} &:= \Av^{\Iu\mid \Iu^A \cap \Iu}_* \circ \oblv^{(\Iu^A,\chi_A)\mid \Iu^A \cap \Iu}[\ell(w_A)].
\end{align*}
Then we have adjoint pairs $(\Av^A_{\chi,!},\Av^A_{*})$ and $(\Av^A_{!},\Av^A_{\chi,*})$. 

The following lemma follows from standard arguments, as e.g.~in~\cite[Lemma~3.2]{AR23}.

\begin{lem}
There exists a canonical isomorphism of functors
\[
\Av^A_{\chi,!} \simeq \Av^A_{\chi,*}.
\]
\end{lem}

In view of this lemma, we will denote both $\Av^A_{\chi,!}$ and $\Av^A_{\chi,*}$ by $\Av^A_\chi$.
We will say more about these averaging functors in~\S\ref{sss:prop-av} below. 

When $J=S$, we are in the traditional setting of Whittaker sheaves on (partial) affine flag varieties as studied e.g.~in~\cite[\S 3]{AR23} (except that the roles of multiplication on the left and on the right are switched compared to this reference). In particular, it follows from~\cite[Lemma~3.3]{AR23} (or, more explicitly, an immediate generalization of this statement to all parabolic affine flag varieties) that for $x \in {}^L W_\ex$ we have
\begin{equation}
\label{eqn:Av-chi-D-classical}
\Av^A_\chi(q^! \Delta^L_{x}) \simeq
\begin{cases}
\Delta^L_{\chi_A,xw} & \text{if $w \in W_A$ satisfies $xw \in {}^L W_\ex^{A,\max}$;} \\
0 & \text{if $xW_A \cap {}^L W_\ex^{A,\max} = \varnothing$,}
\end{cases}
\end{equation}
and similarly for costandard objects.

%---------------------------------------------------------------
\subsection{Convolution actions} 
\label{ss:convolution}
%---------------------------------------------------------------

We now drop the assumptions on $p$ and $\k$ from~\S\ref{ss:Whittaker}.
For any $K \subset S_\af$ finitary,
the $\infty$-category $\Shv(\I_K \backslash \sL G/\I_K)$ admits a canonical structure of algebra object in $\k$-linear presentable stable $\infty$-categories, with binary product given by convolution and denoted $\star^{\I_K}$ (see the discussion and references in~\cite[\S 2.5]{adr}). For any $L \subset J \subset S$, we have a natural right action of $\Shv \left( \I_K \backslash \sL G/\I_K \right)$ on $\Shv \left( \I^\semi_{L \subset J} \backslash \Fl_K \right)$; the corresponding binary operation
\[
\Shv \left( \I^\semi_{L \subset J} \backslash \Fl_K \right) \times \Shv \left( \I_K \backslash \sL G/\I_K \right) \to \Shv \left( \I^\semi_{L \subset J} \backslash \Fl_K \right)
\]
will also be denoted $\star^{\I_K}$.

Let now $X$ be either $\Fl_K$ for some finitary $K \subset S_\af$, $\tFl$, or the symbol $\sL G/(\Iu^A,\chi_A)$ for some finitary $A \subset S_\af$ (in this case we assume that $p>0$ and that $\k$ contains a nontrivial $p$-th root of unity). 
%Then for any $J$ as above, the $\infty$-category $\Shv(\I_J \backslash X)$ is naturally a module for $\Shv(\I_J\backslash \sL G/\I_J)$; the corresponding action bifunctor will also be denoted $\star^{\I_J}$.
For any $L \subset J \subset S$ we have a natural action by convolution from the left 
\[
(-) \star^{\I_{L \subset J}} (-) : \Shv \left( \I_{L \subset J}\backslash \sL{M_J}/\I_{L \subset J} \right) \times \Shv \left( \I_{L \subset J}^{\semi} \backslash X \right) \rightarrow \Shv \left( \I_{L \subset J}^{\semi}\backslash X \right).
\]
%\[
%-\star^{\sL^+{M_J}}-: \Shv(\sL^+{M_J}\backslash \sL{M_J}/\sL^+{M_J})\times \Shv(\I_J^{\semi}\backslash X)\rightarrow \Shv(\I_J^{\semi}\backslash X). 
%\]

If $\theta\in \Xd_J^0$, then $z^\theta$ commutes with $\sL^+ M_J$ and normalizes $\sL U_J$. Hence pushforward under the action of this element on $X$ induces an autoequivalence of $\Shv \left( \I_{L \subset J}^{\semi} \backslash X \right)$, which we will denote $z^\theta \cdot (-)$. We will denote by
$\langle \theta \rangle$ the auto-equivalence given by $\sF\mapsto z^\theta\cdot \sF [-\langle \theta, 2\rho \rangle]$. It is clear (using Remark~\ref{rmk:lJ}\eqref{it:lJ-translation}) that for any $\theta \in \Xd_J^0$ and $x \in {}^L W_\ex^{(K,J)}$ we have
\begin{equation}
\label{eqn:D-N-shift}
\Delta^{L \subset J, \semi}_{K,x} \langle \theta \rangle \simeq \Delta^{L \subset J, \semi}_{K,t_\theta x}, \quad
\nabla^{L \subset J, \semi}_{K,x} \langle \theta \rangle \simeq \nabla^{L \subset J, \semi}_{K,t_\theta x}.
\end{equation}

%%---------------------------------------------------------------
%\subsection{Sheaves on affine flag varieties} 
%\label{subset shvs}
%%---------------------------------------------------------------
%
%We will work in the sheaf-theoretic context explained in \cite[\S\S 2.1--2.2]{adr}, in the case of \'etale sheaves.
%%So let $\k$ be a field, which will play the role of coefficients for sheaves. 
%%%From now on, w
%%We assume that the pair $(\F,\k)$ is one of the following: 
%%\begin{enumerate}
%%\item
%%\label{it:sheaves-etale}
%%$\F$ is arbitrary, 
%So let $\ell$ be a prime number invertible in $\F$, and let $\k$ be either a finite extension or an algebraic closure of $\F_\ell$ or $\Q_\ell$.
%%and we consider \'{e}tale $\k$-sheaves; 
%%\item 
%%$\F=\C$, $\k$ is arbitrary, and we work with sheaves in the analytic topology. 
%%\end{enumerate}
%For a prestack $X$ over $\F$, we have a $\k$-linear presentable stable $\infty$-category $\Shv(X)$ of $\k$-sheaves on $X$. The functorialities that these $\infty$-categories satisfied are briefly discussed in~\cite[\S 2.2]{adr}; in particular, for any map $f : X \to Y$ of $\F$-prestacks we have a $!$-pullback functor $f^! : \Shv(Y) \to \Shv(X)$, which sends the dualizing complex $\omega_Y$ of $Y$ to the dualizing complex $\omega_X$ of $X$. These functors, as well as all the functors between such $\infty$-categories that we will consider below, are continuous, i.e.~colimit-preserving.
%
%\begin{rmk}
%It is likely that all the results below have analogues for sheaves in the analytic topology. However, since a detailed treatment of the construction of such $\infty$-categories is not yet available in the literature, we will stick to \'etale sheaves. (For all the statements involving Whittaker conditions, their counterparts should involve the ``Kirillov model.'')
%\end{rmk}
%
%%===============================
%\subsubsection{Quarter- and semi-infinite sheaves} 
%\label{sss:qs-sheaves}
%%===============================
%
%Fix a subset $J \subset S$, and
%let $X$ be either $\Fl_K$ (for some finitary subset $K \subset S_\af$) or $\tFl$. 
%Below we will study the $\infty$-categories 
%\[
%\Shv(\I_J\backslash X), \quad \Shv(\I_J^{\quar}\backslash X) \quad \text{and}\quad \Shv(\I_J^{\semi}\backslash X).
%\]
%These $\infty$-categories admit a natural description in terms of limits/colimits of $\infty$-cate\-gories of sheaves on algebraic stacks of finite type, which can be obtained by adapting (in the obvious way) the discussion in~\cite[\S 2.5]{adr} for the first case, and that of~\cite[\S 2.6]{adr} in the second and third cases.
%
%%For any $x\in W_\ex^K$, resp. $x\in {}^JW_\ex^K$, we have 
%%\[
%%\Shv(\I_J^{\quar}\backslash {}^J\rQ^K_x) \simeq \Vect_\k, \quad \text{resp.}\quad \Shv(\I_J^{\semi}\backslash {}^J\S^\semi_x) \simeq ?. 
%%\]
%
%Since $\I_J^{\quar}$ is a subgroup of $\I_J^{\semi}$, we
%have a natural forgetful functor
%\[
%\oblv^{\I_J^{\semi}|\I_J^{\quar}}: \Shv(\I_J^{\semi}\backslash X) \rightarrow \Shv(\I_J^{\quar}\backslash X)
%\]
%induced by $!$-pullback along the quotient map $\I_J^{\quar}\backslash X \to \I_J^{\semi}\backslash X$,
%which admits left and right adjoints 
%\[
%\Av_!^{\I_J^{\semi}|\I_J^{\quar}}, \Av_*^{\I_J^{\semi}|\I_J^{\quar}} : \Shv(\I_J^{\quar}\backslash X) \rightarrow \Shv(\I_J^{\semi}\backslash X). 
%\]
%Here $\Av_!^{\I_J^{\semi}|\I_J^{\quar}}$ is given by $!$-pushforward along the map $\I_J^{\quar}\backslash X \to \I_J^{\semi}\backslash X$, and $\Av_*^{\I_J^{\semi}|\I_J^{\quar}}$ is given by $*$-pushforward along this map, followed by $[-2\ell(w_J)]$. In fact, since the projection map is proper we have
%\[
%\Av_*^{\I_J^{\semi}|\I_J^{\quar}} = \Av_!^{\I_J^{\semi}|\I_J^{\quar}}[-2\ell(w_J)].
%\]
%Similarly we have a forgetful functor
%\[
%\oblv^{\I_J|\I} : \Shv(\I_J \backslash X) \to \Shv(\I \backslash X),
%\]
%and its left and right adjoints $\Av_!^{\I_J|\I}$, $\Av_*^{\I_J|\I}$, which differ by the same shift as above.
%
%For $A=\I_J, \I_J^{\quar}$ or $\I_J^{\semi}$, we also have pullback and pushforward functors 
%\begin{gather*}
%(p_K)^*, (p_K)^!: \Shv(A \backslash \Fl_K) \rightarrow \Shv(A \backslash \Fl), \ 
%(p_K)_* = (p_K)_!: \Shv(A \backslash \Fl) \rightarrow \Shv(A \backslash \Fl_K), \\
%%\]
%%\[
%%\text{and}\quad 
%q^*, q^!: \Shv(\mathrm{J}\backslash \Fl) \rightarrow \Shv(\mathrm{J}\backslash \tFl), \ 
%q_*,q_!: \Shv(\mathrm{J}\backslash \tFl) \rightarrow \Shv(\mathrm{J}\backslash \Fl). 
%\end{gather*}
%%Denote by $n_K$ the dimension of the flag variety of $M_K$. Then
%We have $(p_K)^! \simeq (p_K)^*[2 \ell(w_K)]$, and we set
%\[
%(p_K)^\dag = (p_K)^*[\ell(w_K)].
%\]
%Similarly we have $q^! \simeq q^*[2\dim(T)]$, and we set
%\[
%q^\dag = q^*[\dim(T)].
%\]
%
%%===============================
%\subsubsection{Standard and costandard sheaves} 
%%===============================
%
%For $J,K$ as above and $x \in {}^JW_\ex^{(K,J)}$,
%we will also consider the $\infty$-categories 
%\begin{equation}
%\label{eq:cat-orbt}
%\Shv(\I^{\quar}_J\backslash {}^J\rQ^K_x),\quad 
%\Shv(\I^{\semi}_J\backslash {}^J\rS^K_x),\quad 
%\Shv(\I_J\backslash {}^J\Fl^K_x). 
%\end{equation}
%As in~\cite[Proposition~1.3.5]{Gaits18} or \cite[Lemma~3.3.2]{DL23}, $!$-restriction to the point $\dot{x}\I_K/\I_K$ induces equivalences 
%\begin{gather}
%\label{eq:quar-orbt}
%\Shv(\I^{\quar}_J\backslash {}^J\rQ^K_x) \xs \Shv(T\backslash \pt), \\
%\label{eq:semi-orbt}
%\Shv(\I^{\semi}_J\backslash {}^J\rS^K_x) \xs \Shv(\mathrm{Stab}_{M_J}(\dot{x}\I_K/\I_K)\backslash \pt), 
%\\
%\label{eq:fl-orbt}
%\Shv(\I_J\backslash {}^J\Fl^K_x) \xs \Shv(\mathrm{Stab}_{M_J}(\dot{x}\I_K/\I_K)\backslash \pt). 
%\end{gather}
%To simplify notation we will write $\omega_{{}^J\rQ^K_x}$, $\omega_{{}^J\rS^K_x}$ and $\omega_{{}^J\Fl^K_x}$ for $\omega_{\I_J^\quar \backslash {}^J\rQ^K_x}$, $\omega_{\I_J^\semi \backslash {}^J\rS^K_x}$, $\omega_{\I_J \backslash {}^J\Fl^K_x}$ respectively. Under the equivalences \eqref{eq:quar-orbt}, \eqref{eq:semi-orbt}, \eqref{eq:fl-orbt}, these object correspond to the dualizing sheaves of the stacks appearing on the right-hand side.
%%, the object in LHS corresponding to the trivial representation $\k$ in RHS under the equivalence \eqref{eq:quar-orbt}, \eqref{eq:semi-orbt}, \eqref{eq:fl-orbt}, respectively. 
%%For $x\in {}^JW_\ex^{(K,J)}$, resp. $x\in {}^JW_\ex^{K}$, we set 
%%\[
%%\ell^J_K(x)=\max\{\ell^J(y)\}_{y\in W_Jx\cap W_\ex^{(K,J)}}, \quad \text{resp.}\quad 
%%\ell_K(x)=\max\{\ell(y)\}_{y\in W_Jx\cap W_\ex^{K}}. 
%%\]
%We define the \textit{perverse t-structure} on each of the $\infty$-categories appearing in~\eqref{eq:cat-orbt} to be the cohomological shift of the transfer of the natural t-structure on the corresponding right-hand side such that the objects
%%renormalize the t-structure on the categories in \eqref{eq:cat-orbt} such that 
%\begin{gather*}
%\omega_{{}^J\rQ^K_x}[-\ell^J(x)], \quad
%\omega_{{}^J\rS^K_x}[-\ell(w_Jw_{J\cap ({}^xK)})-\ell^J(x)], \\ 
%%\quad \text{and}\quad 
%\omega_{{}^J\Fl^K_x}[-\ell(w_Jw_{J\cap ({}^xK)})-\ell(x)]
%\end{gather*}
%%belong to the heart in the corresponding category, which is called the \textit{perverse t-structure}. 
%belong to the corresponding heart. Here the shift appearing in the third case is chosen so that the forgetful functor $\Shv(\I_J\backslash {}^J\Fl^K_x) \to \Shv({}^J\Fl^K_x)$ is t-exact with respect to the usual perverse t-structure on $\Shv({}^J\Fl^K_x)$, see~\eqref{eqn:dim-par-orbit}. In the other cases we are dealing with schemes of infinite dimension, so that there is no natural choice for the shift; as in~\cite{adr} the shifts are chosen so as to provide a nice theory.
%
%We define the standard and costandard sheaves on $\Fl_K$ as follows:
%\begin{itemize}
%\item 
%in $\Shv(\I_J^\quar \backslash \Fl_K)$, for $x\in W_\ex^{(K,J)}$ we set
%\[
%\left\{
%\begin{array}{c}
%{}^J\Delta^{K,\quar}_x:= ({}^J\bq^K_x)_! \omega_{{}^J\rQ^K_x}[-\ell^J(x)]  \\ 
%{}^J\nabla^{K,\quar}_x:= ({}^J\bq^K_x)_* \omega_{{}^J\rQ^K_x}[-\ell^J(x)]
%\end{array}
%\right. ;
%%, \quad \forall x\in W_\ex^{(K,J)},
%\]
%\item
%in $\Shv(\I_J^\semi \backslash \Fl_K)$, for $x\in {}^JW_\ex^{(K,J)}$ we set
%\[
%\left\{
%\begin{array}{c}
%{}^J\Delta^{K,\semi}_x := ({}^J\bs^K_x)_! \omega_{{}^J\rS^K_x}[-\ell(w_Jw_{J\cap ({}^xK)})-\ell^J(x)] \\ 
%{}^J\nabla^{K,\semi}_x:= ({}^J\bs^K_x)_* \omega_{{}^J\rS^K_x}[-\ell(w_Jw_{J\cap ({}^xK)})-\ell^J(x)]
%\end{array}
%\right. ;
%%, \quad \forall x\in {}^JW_\ex^{(K,J)}, 
%\]
%\item
%in $\Shv(\I_J \backslash \Fl_K)$, for $x\in {}^JW_\ex^K$ we set
%\[
%\left\{
%\begin{array}{c}
%{}^J\Delta^K_x:= ({}^J\bi^K_x)_! \omega_{{}^J\Fl^K_x}[-\ell(w_Jw_{J\cap ({}^xK)})-\ell(x)] \\ 
%{}^J\nabla^K_x:= ({}^J\bi^K_x)_* \omega_{{}^J\Fl^K_x}[-\ell(w_Jw_{J\cap ({}^xK)})-\ell(x)]
%\end{array}
%\right. .
%%, \quad \forall x\in {}^JW_\ex^K. 
%\]
%\end{itemize}
%We will omit $J$, resp. $K$, in the notations above if it is the empty set, e.g. $\Delta_x={}^\varnothing \Delta^{\varnothing}_x$, $\Delta^{K,\semi}_x={}^\varnothing \Delta^{K,\semi}_x$ and ${}^J\Delta^{\quar}_x={}^J \Delta^{\varnothing,\quar}_x$. 
%%We will also replace $\Delta^{K,\semi}_x={}^{\emptyset}\Delta^{K,\quar}_x$ and $\nabla^{K,\semi}_x={}^{\emptyset}\nabla^{K,\quar}_x$. 
%%Suppose that either $(J,K)=(J,\emptyset)$ or $(\emptyset,K)$. 
%%We define the standard and costandard sheaves on $\Fl$ as follows 
%%\[
%%{}^J\Delta_x:= ({}^J\iota_x)_! \k_{{}^J\Fl_x}[\ell(w_Jx)], \quad 
%%{}^J\nabla_x:= ({}^J\iota_x)_* \k_{{}^J\Fl_x}[\ell(w_Jx)], \quad \forall x\in {}^JW_\ex,
%%\]
%%\[
%%{}^J\Delta^{\quar}_x:= ({}^J\bi_x)_! \omega_{{}^J\rQ_x}[-\ell^J(x)], \quad 
%%{}^J\nabla^{\quar}_x:= ({}^J\bi_x)_* \omega_{{}^J\rQ_x}[-\ell^J(x)], \quad \forall x\in W_\ex,
%%\]
%%\[
%%{}^J\Delta^{\semi}_x:= ({}^J\bj_x)_! \omega_{{}^J\S_x}[-\ell^J(w_Jx)], \quad 
%%{}^J\nabla^{\semi}_x:= ({}^J\bj_x)_* \omega_{{}^J\S_x}[-\ell^J(w_Jx)], \quad \forall x\in {}^JW_\ex,
%%\]
%%and define standard and costandard sheaves on $\Fl_K$, 
%%\[
%%\Delta^K_x:= (\iota^K_x)_! \k_{\Fl^K_x}[\ell(x)], \quad 
%%\nabla^K_x:= (\iota^K_x)_* \k_{\Fl^K_x}[\ell(x)], \quad \forall x\in W^K_\ex,
%%\]
%%\[
%%{}^J\Delta^{K,\quar}_x:= ({}^J\bi^K_x)_! \omega_{{}^J\rQ^K_x}[-\ell^J(x)], \quad 
%%{}^J\nabla^{K,\quar}_x:= ({}^J\bi^K_x)_* \omega_{{}^J\rQ^K_x}[-\ell^J(x)], \quad \forall x\in W_\ex^{(K,J)}.
%%\]
%
%\begin{lem}
%\label{lem:Av-quar-semi-D-N}
%Let $x\in {}^JW_\ex^{(K,J)}$. 
%\begin{enumerate}
%\item 
%\label{it:Av-quar-semi-D-N-1}
%Let $w\in W_J^{J\cap ({}^x K)}$. 
%The object 
%\[
%\Av_!^{\I_J^{\semi}|\I_J^{\quar}}({}^J\Delta^{K,\quar}_{wx}), \quad 
%\text{resp.}\quad \Av_*^{\I_J^{\semi}|\I_J^{\quar}}({}^J\nabla^{K,\quar}_{wx}),
%\]
%is an iterated extension of objects
%\[
%{}^J\Delta^{K,\semi}_x[\ell(w_{J\cap ({}^xK)}u)+\ell(w_Jwu)], \
%\text{resp.}\ {}^J\nabla^{K,\semi}_x[-\ell(w_{J\cap ({}^xK)}u)-\ell(w_Jwu)],
%\]
%with $u\in W_{J\cap ({}^xK)}$, each of them appearing once. 
%\item 
%\label{it:Av-quar-semi-D-N-2}
%The object 
%\[
%\oblv^{\I_J^{\semi}|\I_J^{\quar}}({}^J\Delta^{K,\semi}_x), 
%\quad \text{resp.}\quad 
%\oblv^{\I_J^{\semi}|\I_J^{\quar}}({}^J\nabla^{K,\semi}_x),
%\]
%is an iterated extension of objects
%\[
%{}^J\Delta^{K,\quar}_{wx}[\ell(w_Jw_{J\cap ({}^xK)})-\ell(w)], 
%\ \text{resp.}\ 
%{}^J\nabla^{K,\quar}_{wx}[-\ell(w_Jw_{J\cap ({}^xK)})+\ell(w)],
%\]
%with $w\in W_J^{J\cap ({}^xK)}$, each of them appearing once. 
%\end{enumerate}
%\end{lem}
%
%\begin{proof}
%\eqref{it:Av-quar-semi-D-N-1}
%Lemma~\ref{lem:bi-min-av}\eqref{it:bi-min-av-1} allows to reduce the proof to the case $w=1$. Then we use Lemma~\ref{lem:bi-min-av}\eqref{it:bi-min-av-2}. For the first claim we have to compute the $!$-pushforward under that map of $\omega_{\I_J^\semi \times^{\I_J^\quar} {}^J \rQ_x^K}$, and the claim follows from the fact that the cohomology of $\omega_{M_{J \cap ({}^x K)}}$ is $\bigoplus_{u \in W_{J \cap ({}^x K)}} \k[2\ell(w_{J \cap ({}^x K)}) - 2\ell(u)]$. For the second claim we use the fact that $\Av_*^{\I_J^{\semi}|\I_J^{\quar}} = \Av_!^{\I_J^{\semi}|\I_J^{\quar}}[-2\ell(w_J)]$ (see~\S\ref{sss:qs-sheaves}), and write now the cohomology above as $\bigoplus_{u \in W_{J \cap ({}^x K)}} \k[2\ell(u)]$.
%
%\eqref{it:Av-quar-semi-D-N-2}
%Using the decomposition from Lemma~\ref{lem:bi-min-av}\eqref{it:bi-min-av-3}, we obtain that $\omega_{{}^J \rS^K_x}$ is an iterated extension of objects $f_! f^* \omega_{{}^J \rS^K_x}$, resp.~$f_* f^! \omega_{{}^J \rS^K_x}$, where $f$ runs over the locally closed immersions ${}^J \rQ^K_{wx} \hookrightarrow {}^J \rS^K_x$ ($w \in W_J^{J\cap ({}^xK)}$). The second assertion immediately follows. For the first one, we need to compute the complexes $f^* \omega_{{}^J \rS^K_x}$ for $f$ as above.
%We write the embedding ${}^J \rQ^K_{wx} \hookrightarrow {}^J \rS^K_x$ as the composition
%\[
%{}^J \rQ^K_{wx} \hookrightarrow \I_J^\semi \times^{\I_J^\quar} {}^J \rQ_x^K \to {}^J \rS^K_x.
%\]
%The second map is smooth of relative dimension $\ell(w_{J \cap ({}^x K)})$ by Lemma~\ref{lem:bi-min-av}\eqref{it:bi-min-av-2}, so the $*$-pullback under this map of $\omega_{{}^J \rS^K_x}$ is $\omega_{\I_J^\semi \times^{\I_J^\quar} {}^J \rQ_x^K}[-2\ell(w_{J \cap ({}^x K)})]$. On the other hand, using Lemma~\ref{lem:bi-min-av}\eqref{it:bi-min-av-1} we obtain that the $*$-pullback under the first map of $\omega_{\I_J^\semi \times^{\I_J^\quar} {}^J \rQ_x^K}$ is $\omega_{{}^J \rQ^K_{wx}}[2\ell(w_J)-2\ell(w)]$. The claim follows.
%%It follows from a standard computation using 
%%This follows from Lemma~\ref{lem:bi-min-av} and standard arguments. 
%\end{proof}
%
%\begin{lem}
%\label{lem:push-pull-pK-D-N-quar}
%Let $x\in W_\ex^{(K,J)}$. 
%\begin{enumerate}
%\item 
%\label{it:push-pull-pK-D-N-quar-1}
%For any $w\in W_K$, we have 
%\[
%(p_K)_*({}^J\Delta^{\quar}_{xw})\simeq {}^J\Delta^{K,\quar}_x[-\ell(w)], \quad 
%(p_K)_*({}^J\nabla^{\quar}_{xw})\simeq {}^J\nabla^{K,\quar}_x[\ell(w)]. 
%\]
%\item
%\label{it:push-pull-pK-D-N-quar-2}
%The object $(p_K)^\dag({}^J\Delta^{K,\quar}_x)$, resp. $(p_K)^\dag({}^J\nabla^{K,\quar}_x)$, is an iterated extension of objects ${}^J\Delta^{\quar}_{xw}[\ell(w_K)-\ell(w)]$ ($w\in W_K$), resp. ${}^J\nabla^{\quar}_{xw}[\ell(w)-\ell(w_K)]$ ($w\in W_K$), each of which appearing once. 
%\end{enumerate}
%\end{lem}
%
%\begin{proof}
%\eqref{it:push-pull-pK-D-N-quar-1}
%Using Lemma~\ref{lem:min-push}, the claims follows from the facts that the cohomology with compact supports and cohomology of $\omega_{\mathbb{A}^{\ell(w)}}[-\ell(w)]$ is $\k[-\ell(w)]$ and $\k[\ell(w)]$ respectively.
%
%\eqref{it:push-pull-pK-D-N-quar-2}
%Using Lemma~\ref{lem:min-push}, the claims follow from the fact that $\omega_{\I_K/\I}[-\ell(w_K)]$ is an iterated extension of objects $\Delta_w[\ell(w_K)-\ell(w)]$ ($w \in W_K$), and also an iterated extension of objects $\nabla_w[\ell(w)-\ell(w_K)]$ ($w \in W_K$).
%\end{proof}
%
%\begin{lem}
%\label{lem:push-pull-pK-D-N-semi}
%Let $x\in {}^JW_\ex^{(K,J)}$. 
%\begin{enumerate}
%\item 
%\label{it:push-pull-pK-D-N-semi-1}
%Let $w\in {}^{J^x\cap K}W_K$. 
%The object 
%\[
%(p_K)_*({}^J\Delta^{\semi}_{xw})[\ell(w_K)], 
%\quad \text{resp.}\quad 
%(p_K)_*({}^J\nabla^{\semi}_{xw})[-\ell(w_K)], 
%\]
%is an iterated extension of objects
%\[
%{}^J\Delta^{K,\semi}_x[\ell(uw_{(J^x)\cap K})+\ell(uww_K)], \
%\text{resp.}\ {}^J\nabla^{K,\semi}_x[-\ell(uw_{(J^x)\cap K})-\ell(uww_K)],
%\]
%with $u\in W_{(J^x)\cap K}$, each appearing once. 
%\item
%\label{it:push-pull-pK-D-N-semi-2}
%The object 
%\[
%(p_K)^\dag({}^J\Delta^{K,\semi}_x), 
%\quad \text{resp.}\quad 
%(p_K)^\dag({}^J\nabla^{K,\semi}_x), 
%\]
%is an iterated extension of objects
%\[
%{}^J\Delta^{\semi}_{xw}[\ell(w_{(J^x)\cap K}w_K)-\ell(w)], \
%\text{resp.}\ {}^J\nabla^{\semi}_{xw}[-\ell(w_{(J^x)\cap K}w_K)+\ell(w)],
%\]
%with $w\in {}^{(J^x)\cap K}W_K$, each appearing once. 
%\end{enumerate}
%\end{lem}
%
%\begin{proof}
%\eqref{it:push-pull-pK-D-N-semi-1}
%Using Lemma~\ref{lem:bi-min-push}\eqref{it:bi-min-push-1}, we reduce the proof to the case $w=1$. In that case we use Lemma~\ref{lem:bi-min-push}\eqref{it:bi-min-push-2} to see that $(p_K)_*({}^J\Delta^{\semi}_{x})$ is an iterated extension of objects ${}^J\Delta^{K,\semi}_{x}[\ell(w_{J \cap ({}^x K)})-2\ell(u)]$. The first assertion follows, using the fact that $\ell(w_{J \cap ({}^x K)}) = \ell(w_{(J^x) \cap K})$. The second assertion follows from the same considerations (and replacing $u$ by $u w_{(J^x) \cap K}$ in the formula above).
%
%\eqref{it:push-pull-pK-D-N-semi-2}
%By the base change theorem and smoothness of $p_K$, the complex $(p_K)^\dag({}^J\Delta^{K,\semi}_x)$, resp.~$(p_K)^\dag({}^J\nabla^{K,\semi}_x)$, is the $!$-pushforward, resp.~$*$-pushforward, of
%\[
%\omega_{(p_K)^{-1}({}^J \rS^K_x)}[-\ell(w_J)+\ell(w_{J \cap ({}^x K)})-\ell^J(x)-\ell(w_K)]
%\]
%along the immersion $(p_K)^{-1}({}^J \rS^K_x) \hookrightarrow \Fl$. Now, by the decomposition in Lem\-ma~\ref{lem:bi-min-push}\eqref{it:bi-min-push-3}, $\omega_{(p_K)^{-1}({}^J \rS^K_x)}$ is an iterated extension of $*$-pushforwards of complexes $\omega_{{}^J \rS_{xw}}$ ($w \in {}^{(J^x) \cap K} W_K$), which implies the second assertion using the formula $\ell(w_{J \cap ({}^x K)}) = \ell(w_{(J^x) \cap K})$. Similarly, $\omega_{(p_K)^{-1}({}^J \rS^K_x)}$ is an iterated extension of objects $f_! f^* \omega_{(p_K)^{-1}({}^J \rS^K_x)}$ where $f$ runs over the immersions ${}^J \rS_{xw} \hookrightarrow (p_K)^{-1}({}^J \rS^K_x)$, and arguments similar to those encountered in the proof of Lemma~\ref{lem:Av-quar-semi-D-N}\eqref{it:Av-quar-semi-D-N-2}
%imply the first assertion.
%%This follows from standard arguments using Lemma~\ref{lem:bi-min-push}. 
%\end{proof}
%
%%===============================
%\subsubsection{Whittaker sheaves} 
%\label{sss:Whittaker}
%%===============================
%
%%Let us assume that $(\F,\K)$ is in situation \eqref{it:sheaves-etale}. 
%
%Let $K\subset S$. 
%Recall that $w_K$ is the longest element in $W_K$, and $\dot{w}_K$ is a lift of $w_K$ in $G$. 
%We set ${}^KU=\Ad_{\dot{w}_K}(U)$ and ${}^K\Iu=\Ad_{\dot{w}_K}(\Iu)$, and consider the morphisms of algebraic groups 
%\[
%{}^K\Iu\twoheadrightarrow {}^KU \twoheadrightarrow \prod_{\alpha\in -\Phis_K} U_\alpha \simeq \prod_{\alpha\in -\Phis_K} \Ga \xrightarrow{\text{sum}} \Ga.
%\]
%(Here, the first morphism is induced by the projection $\Iu \to U$, and for the second one we consider the decomposition of ${}^K U$ as a product of root subgroups, and the projection to those corresponding to roots on $-\Phis_K$.)
%We fix an Artin--Schreier local system on $\Ga$ (i.e.~a nonconstant irreducible direct summand of the local system $\mathrm{AS}_* \underline{\k}_{\Ga}$, where $\mathrm{AS} : \Ga \to \Ga$ is the Artin--Schreier map), and denote by $\chi_K$, resp.~$\chi'_K$, its pullback to ${}^K\Iu$, resp.~to ${}^KU$. 
%Combining the constructions from~\S\ref{sss:qs-sheaves} with the consideration of Whittaker sheaves as in~\cite[End of~\S 4.2]{zhu},
%one can define the $\infty$-categories
%\[
%\Shv(\I_J\backslash \sL G/({}^K\Iu,\chi_K)), \ \ \Shv(\I_J^{\quar}\backslash \sL G/({}^K\Iu,\chi_K)) \ \ \text{and}\ \ \Shv(\I_J^{\semi}\backslash \sL G/({}^K\Iu,\chi_K)). 
%\]
%
%Let us explain the definition of $\Shv(\I_J^{\quar}\backslash \sL G/({}^K\Iu,\chi_K))$ as an example. 
%%Let $\theta\in \Xd_J^0 \cap \Xd_{J^\comp}^{++}$.
%%% be $J^c$-dominant enough. 
%%We set 
%%\[
%%\I_{J,(n)}^{\quar}=\I_{M_J}\ltimes \Ad_{z^{-n\theta}}(\sL^+U_{J}),
%%%=\Ad_{z^{-n\theta}}(\I_{M_J}\ltimes \sL^+U_{J})
%%\]
%%then $\{\I_{J,(n)}^{\quar}\}_n$ forms an exhaustive filtration of $\I_{J}^{\quar}$. 
%%We set $\widetilde{\Gr}=\sL G/\sL^{++}G$, which admits an $\sL G$-action from the left and a ${}^KU$-action from the right. 
%%For any $n$ one can define the category 
%%\[
%%\Shv(\I_{J,(n)}^{\quar}\backslash \sL G/ ({}^K\Iu,\chi_K))
%%=\Shv\big( \I_{J,(n)}^{\quar} \backslash \widetilde{\Gr} / ({}^KU,\chi'_K) \big).
%%\]
%%Then we define 
%%\[
%%\Shv(\I_J^{\quar}\backslash \sL G/({}^K\Iu,\chi_K)) = \varprojlim_{n\geq 0} \Shv(\I_{J,(n)}^{\quar}\backslash \sL G/ ({}^K\Iu,\chi_K)), 
%%\]
%%where the transition maps are given by the functor $\oblv^{\I_{J,(n+1)}^{\quar}\mid \I_{J,(n)}^{\quar}}$. 
%Let $\widetilde{{}^K\Iu} = {}^K \Iu \times_{\Ga} \Ga$, where the map $\Ga \to \Ga$ is $\mathrm{AS}$. Then one can consider the stack $\I_J^\quar \backslash \sL G / \widetilde{{}^K\Iu}$. The $\infty$-category
%\[
%\Shv(\I_J^\quar \backslash \sL G / \widetilde{{}^K\Iu})
%\]
%breaks as a direct sum indexed by the $p$ different irreducible direct summands of the pushforward of $\mathrm{AS}_* \underline{\k}_{\Ga}$, and $\Shv(\I_J^{\quar}\backslash \sL G/({}^K\Iu,\chi_K))$ is defined as the direct summand corresponding to our choice of Artin--Schreier local system.
%
%It is convenient to introduce the ind-scheme $\widetilde{\Gr}=\sL G/\sL^{++}G$, which admits an $\sL G$-action from the left and a ${}^KU$-action from the right; then we have
%\[
%\Shv(\I_J^{\quar}\backslash \sL G/({}^K\Iu,\chi_K)) = \Shv(\I_J^{\quar}\backslash \widetilde{\Gr} /({}^K U,\chi'_K)),
%\]
%and similarly for the other variants.
%Denote by $\widetilde{q}: \widetilde{\Gr}\rightarrow \Fl$ the natural projection. 
%By \eqref{eq:decomp} we have decompositions of $\widetilde{\Gr}$ into $(\I_{J}^{\quar}\times {}^KU)$-, $(\I_{J}^{\semi}\times {}^KU)$-, and $(\I_{J}\times {}^KU)$-orbits 
%\begin{equation}\label{eq:decomp-tGr}
%\widetilde{\Gr}= \bigsqcup_{x\in W_\ex} \widetilde{q}^{-1}({}^{J}\rQ_x)\dot{w}_K
%= \bigsqcup_{x\in {}^JW_\ex} \widetilde{q}^{-1}({}^{J}\rS_x)\dot{w}_K 
%= \bigsqcup_{x\in {}^JW_\ex} \widetilde{q}^{-1}({}^{J}\Fl_x)\dot{w}_K. 
%\end{equation}
%
%\begin{lem}
%\label{lem:whit}
%\begin{enumerate}
%\item
%\label{it:lem-whit-1}
%For $x\in W_\ex$, $!$-restriction to the point $\dot{x}\dot{w}_K$ induces an equivalence 
%\begin{equation}
%\label{eq:quar-orbt-whit}
%\Shv\big(\I_{J}^{\quar} \backslash \widetilde{q}^{-1}({}^{J}\rQ_x)\dot{w}_K / ({}^KU,\chi'_K) \big)
%\simeq
%\begin{cases}
%\Vect_\k & \text{if $x\in W_\ex^{(K,J),\mathrm{w}}$;} \\ 
%0 & \text{otherwise.}
%\end{cases}
%\end{equation}
%\item 
%For $x\in {}^JW_\ex$, $!$-restriction to the point $\dot{x}\dot{w}_K$ induces an equivalence 
%\begin{equation}
%\label{eq:semi-orbt-whit}
%\Shv\big( \I_{J}^{\semi} \backslash \widetilde{q}^{-1}({}^{J}\rS_x)\dot{w}_K / ({}^KU,\chi'_K) \big)
%\simeq
%\begin{cases}
%\Vect_\k & \text{if $x\in {}^JW_\ex^{(K,J),\mathrm{w}}$;} \\ 0 & \text{otherwise.} 
%\end{cases}
%\end{equation}
%\item 
%For $x\in {}^JW_\ex$, $!$-restriction to the point $\dot{x}\dot{w}_K$ induces an equivalence 
%\begin{equation}
%\label{eq:fl-orbt-whit}
%\Shv\big(\I_{J} \backslash \widetilde{q}^{-1}({}^{J}\Fl_x)\dot{w}_K / ({}^KU,\chi'_K) \big)
%\simeq
%\begin{cases}
%\Vect_\k & \text{if $x\in {}^JW_\ex^{K,\mathrm{w}}$;} \\ 0 & \text{otherwise.} 
%\end{cases}
%\end{equation}
%\end{enumerate}
%\end{lem}
%
%\begin{proof}
%We only explain the proof of~\eqref{it:lem-whit-1}; those of the other statements are similar. 
%Let $\theta\in \Xd_J^0 \cap \Xd_{J^\comp}^{++}$, and for $n \geq 0$ set 
%\[
%\I_{J,(n)}^{\quar}=\I_{M_J}\ltimes \Ad_{z^{-n\theta}}(\sL^+U_{J});
%\]
%then $( \I_{J,(n)}^{\quar} : n \geq 0)$ forms an exhaustive filtration of $\I_{J}^{\quar}$. 
%The $\infty$-category in the left-hand side of~\eqref{eq:quar-orbt-whit} is equivalent to the limit $\lim_{n} \sC_n$, where 
%\[
%\sC_n:= \Shv\big(\I_{J,(n)}^{\quar} \backslash (\I_{J,(n)}^{\quar}\dot{x}\dot{w}_K \sL^{++}G/\sL^{++}G) / ({}^KU,\chi'_K) \big).
%\]
%It is therefore enough to show that for $n \gg 0$, $!$-restriction to the point $\dot{x}\dot{w}_K$ induces an equivalence 
%\begin{equation}
%\label{eq:Cn}
%\sC_n \simeq
%\begin{cases}
%\Vect_\k & \text{if $x\in W_\ex^{(K,J),\mathrm{w}}$;} 
%\\ 0 & \text{otherwise.} 
%\end{cases}
%\end{equation}
%By Lemma~\ref{lem:min-length}\eqref{it:min-length-2}, $x\in W_\ex^{(K,J),\mathrm{w}}$, equivalently $xw_K\in W_\ex^{(K,J)}$, if and only if $xw_K\preceq^{J}xw_K s_\alpha$ for any $\alpha\in \Phis_K$, which by Lemma~\ref{lem:quar-order} is equivalent to the condition that 
%\[
%\Ad_{\dot{x}\dot{w}_K}(U_{-\alpha}) \cap \I_J^{\quar} = \{1\}, \quad \forall \alpha\in \Phis_K. 
%\]
%In this case $\Ad_{\dot{x}\dot{w}_K}(U_{-\alpha}) \cap \I_{J,(n)}^{\quar} = \{1\}$ for any $n\in \Z$. 
%Hence we obtain \eqref{eq:Cn}, provided $n \gg 0$. 
%\end{proof}
%
%We denote by $\omega_{\widetilde{q}^{-1}({}^{J}\rQ_x)\dot{w}_K}^{\chi_K}$ the object in the left hand side of~\eqref{eq:quar-orbt-whit} corresponding to $\k\in \Vect_\k$.
%% under the equivalence \eqref{eq:quar-orbt-whit}. 
%Let ${}^J\widetilde{\bi}^{K}_x$ be the embedding of $\widetilde{q}^{-1}({}^{J}\rQ_x)\dot{w}_K$ in $\widetilde{\Gr}$. 
%We define the standard and costandard sheaves in $\Shv(\I_J^{\quar}\backslash \sL G/({}^K\Iu,\chi_K))$ as 
%\[
%{}^J\Delta^{\chi_K,\quar}_{x}:= ({}^J\widetilde{\bi}^{K}_x)_! \omega_{\widetilde{q}^{-1}({}^{J}\rQ_x)\dot{w}_K}^{\chi_K}[-\ell^J(x)],
%\]
%\[
%{}^J\nabla^{\chi_K,\quar}_{x}:= ({}^J\widetilde{\bi}^{K}_x)_*\omega_{\widetilde{q}^{-1}({}^{J}\rQ_x)\dot{w}_K}^{\chi_K}[-\ell^J(x)]. 
%\]
%We can similarly define the standard and costandard sheaves 
%\begin{gather*}
%{}^J\Delta^{\chi_K,\semi}_{x}, \
%%\quad \text{and}\quad 
%{}^J\nabla^{\chi_K,\semi}_{x} 
%%\quad 
%\in \Shv(\I_J^{\semi}\backslash \sL G/({}^K\Iu,\chi_K)), \\
%%\]
%%\[
%{}^J\Delta^{\chi_K}_{x}, \
%%\quad \text{and}\quad 
%{}^J\nabla^{\chi_K}_{x} 
%%\quad 
%\in \Shv(\I_J\backslash \sL G/({}^K\Iu,\chi_K)). 
%\end{gather*}
%
%
%\begin{lem}
%\label{lem:av-oblv-whit}
%\begin{enumerate}
%\item 
%\label{it:av-oblv-whit-1}
%Let $x\in W_\ex^{(K,J),\mathrm{w}}$. 
%We have 
%\[
%\Av_!^{\I_J^{\semi}|\I_J^{\quar}}({}^J\Delta^{\chi_K,\quar}_{x})\simeq
%\begin{cases}
%{}^J\Delta^{\chi_K,\semi}_{wx}[\ell(w_Jw)] & \text{if $w\in W_J$ satisfies $wx\in {}^JW_\ex^{(K,J),\mathrm{w}}$;} \\ 
%0 & \text{otherwise},
%\end{cases}
%\]
%and 
%\[
%\Av_*^{\I_J^{\semi}|\I_J^{\quar}}({}^J\nabla^{\chi_K,\quar}_{x})\simeq
%\begin{cases}
%{}^J\nabla^{\chi_K,\semi}_{wx}[-\ell(w_Jw)] & \text{if $w\in W_J$ satisfies $wx\in {}^JW_\ex^{(K,J),\mathrm{w}}$;} \\ 
%0 & \text{otherwise.} 
%\end{cases}
%\]
%\item 
%\label{it:av-oblv-whit-2}
%Let $x\in {}^JW_\ex^{(K,J),\mathrm{w}}$. 
%Then $\oblv^{\I_J^{\semi}|\I_J^{\quar}}({}^J\Delta^{\chi_K,\semi}_x)$, resp. $\oblv^{\I_J^{\semi}|\I_J^{\quar}}({}^J\nabla^{\chi_K,\semi}_x)$, is an iterated extension of objects ${}^J\Delta^{\chi_K,\quar}_{wx}[\ell(w_Jw)]$, resp. ${}^J\nabla^{\chi_K,\quar}_{wx}[-\ell(w_Jw)]$, with $w\in W_J$, each of appearing once. 
%\end{enumerate}
%\end{lem}
%
%\begin{proof}
%We show the first assertion in~\eqref{it:av-oblv-whit-1}; the other ones can be obtained similarly. 
%Let $w\in W_J$ be such that $wx\in {}^JW_\ex$. 
%Consider the functor 
%\begin{equation}
%\label{eq:av-orbt}
%\begin{aligned}
%\Av_!^{\I_J^{\semi}|\I_J^{\quar}}\circ ({}^J\widetilde{\bi}^{K}_x)_! :\
%\Shv\big(\I_{J}^{\quar} 
%& \backslash \widetilde{q}^{-1}({}^{J}\rQ_x)\dot{w}_K  / ({}^KU,\chi'_K) \big) \\
%&\rightarrow \Shv\big(\I_{J}^{\semi} \backslash \widetilde{q}^{-1}({}^{J}\rS_{wx})\dot{w}_K / ({}^KU,\chi'_K) \big),
%\end{aligned}
%\end{equation}
%where we still write ${}^J\widetilde{\bi}^{K}_x$ for the embedding of $\widetilde{q}^{-1}({}^{J}\rQ_x)\dot{w}_K$ in $\widetilde{q}^{-1}({}^{J}\rS_{wx})\dot{w}_K$. 
%It is enough to show that under the equivalences~\eqref{eq:quar-orbt-whit} and~\eqref{eq:semi-orbt-whit}, the functor~\eqref{eq:av-orbt} sends $\underline{\k}$ to $\underline{\k}[\ell(w_Jw)]$ if $wx\in {}^JW_\ex^{(K,J),\mathrm{w}}$ and to 0 otherwise. 
%This follows from a standard computation. 
%\end{proof}
%
%%===============================
%\subsubsection{Averaging functors} 
%\label{sss:av-fun}
%%===============================
%
%Let $A=\I_J, \I_J^{\quar}$ or $\I_J^{\semi}$. 
%Consider the forgetful functor 
%\[
%\oblv^{({}^K\Iu,\chi_K)\mid {}^K\Iu\cap \Iu}: \Shv(A \backslash \sL G/ ({}^K\Iu,\chi_K)) \rightarrow \Shv(A \backslash \sL G/ {}^K\Iu\cap \Iu)
%\]
%which admits left and right adjoints 
%\[
%\av^K_{\chi_K,!},\ \av^K_{\chi_K,*}: \Shv(A \backslash \sL G/ {}^K\Iu\cap \Iu)\rightarrow \Shv(A \backslash \sL G/ ({}^K\Iu,\chi_K)). 
%\]
%Similarly we have the forgetful functor 
%\[
%\oblv^{\Iu\mid {}^K\Iu\cap \Iu}: \Shv(A \backslash \tFl) \rightarrow \Shv(A \backslash \sL G/ {}^K\Iu\cap \Iu),
%\]
%and its left and right adjoints 
%\[
%\av^K_{!},\ \av^K_{*}: \Shv(A \backslash \sL G/ {}^K\Iu\cap \Iu) \rightarrow \Shv(A \backslash \tFl). 
%\]
%We set 
%\[
%\Av^K_{\chi_K,!}:=\av^K_{\chi_K,!}\circ \oblv^{\Iu\mid {}^K\Iu\cap \Iu}[-\ell(w_K)], \ 
%\Av^K_{\chi_K,*}:=\av^K_{\chi_K,*}\circ \oblv^{\Iu\mid {}^K\Iu\cap \Iu}[\ell(w_K)],
%\]
%\[
%\Av^K_{!}:=\av^K_{!}\circ \oblv^{({}^K\Iu,\chi_K)\mid {}^K\Iu\cap \Iu}[-\ell(w_K)], \ 
%\Av^K_{*}:=\av^K_{*}\circ \oblv^{({}^K\Iu,\chi_K)\mid {}^K\Iu\cap \Iu}[\ell(w_K)].
%\]
%Then we have adjoint pairs $(\Av^K_{\chi_K,!},\Av^K_{*})$ and $(\Av^K_{!},\Av^K_{\chi_K,*})$. 
%
%We will say more about the averaging functors in \S\ref{sss:prop-av} below. 
%
%%===============================
%\subsubsection{Convolution actions} 
%%===============================
%The category $\Shv(\I_J\backslash \sL G/\I_J)$ admits a monoidal structure given by convolution product, and we denote the monoidal product by $\star^{\I_J}$. 
%
%Let $X=\Fl_K$ for some $K$, $\tFl$, or the symbol $\sL G/({}^K\Iu,\chi_K)$ when $(\F,\K)$ is in situation \eqref{it:sheaves-etale}. 
%We have natural actions by convolution from the left 
%\[
%-\star^{\I_{M_J}}-: \Shv(\I_{M_J}\backslash \sL{M_J}/\I_{M_J})\times \Shv(\I_J^{\quar}\backslash X)\rightarrow \Shv(\I_J^{\quar}\backslash X),
%\]
%\[
%-\star^{\sL^+{M_J}}-: \Shv(\sL^+{M_J}\backslash \sL{M_J}/\sL^+{M_J})\times \Shv(\I_J^{\semi}\backslash X)\rightarrow \Shv(\I_J^{\semi}\backslash X). 
%\]
%
%For $\theta\in \Xd_J^0$, we denote by $\langle \theta \rangle$ the auto-equivalence of $\Shv(\I_J^{\quar}\backslash X)$, resp.~$\Shv(\I_J^{\semi}\backslash X)$, given by $\sF\mapsto z^\theta\cdot \sF [\ell(t_\theta)]$. 

%%%%%%%%%%%%%%%%%%%%%%%%%
\section{Raskin's equivalences} 
\label{sec:Raskin}
%%%%%%%%%%%%%%%%%%%%%%%%%

Let $X$ be either $\Fl_K$ for some finitary $K \subset S_\af$, $\tFl$, or the symbol $\sL G/(\Iu^A,\chi_A)$ for some finitary $A \subset S_\af$ (assuming in this case that $p>0$ and that $\k$ contains a nontrivial $p$-th root of unity).
% when $(\F,\K)$ is in situation \eqref{it:sheaves-etale}. 

%---------------------------------------------------------------
\subsection{Construction and functorialities} 
%---------------------------------------------------------------

%We have the following equivalences whose construction is inspired by work of Raskin.
%, which holds in a general framwork for categories equipped with an action of $\sL{G}$, see discussions in 
The following proposition generalizes constructions
discussed in~\cite[\S 5.2]{Gaits18},~\cite[Proposition~3.1.7]{DL23} and~\cite[\S2.7]{adr}, which themselves generalize an earlier construction of Raskin. These equivalences sometimes allow one to reduce the proof of properties of semiinfinite sheaves to the more familiar case of sheaves equivariant with respect to a parahoric subgroup.

\begin{prop}
\label{prop:Raskin}
For any subsets $L \subset J \subset S$, there exists a canonical equivalence of $\infty$-categories
\[
\Ras^{\I_{L \subset J}^{\semi} \mid \I_L}: \Shv \left( \I_{L \subset J}^{\semi} \backslash X \right) \simto \Shv \left( \I_L \backslash X \right).
\]
%\]
%and
%\[
%\Ras^{\I_{J}^{\semi}\mid \I_{J}}: \Shv(\I_{J}^{\semi}\backslash X) \simto \Shv(\I_{J}\backslash X)
%\]
%and a commutative diagram
%\[
%\begin{tikzcd}[column sep=huge]
%\Shv(\I_{J}^{\semi}\backslash X) \ar[r, "\Ras^{\I_{J}^{\semi}\mid \I}", "\sim"'] \ar[d, "\oblv^{\I_{J}^{\semi}|\I_{J}^{\quar}}"'] & \Shv(\I_J \backslash X) \ar[d, "\oblv^{\I_J|\I}"] \\
%\Shv(\I_{J}^{\quar}\backslash X) \ar[r, "\Ras^{\I_{J}^{\quar}\mid \I_{J}}", "\sim"'] & \Shv(\I \backslash X).
%\end{tikzcd}
%\]
If $L_1 \subset L_2 \subset J \subset S$ we have a commutative diagram
\begin{equation}
\label{eq:comm-diag-ras-oblv}
\begin{tikzcd}[column sep=3cm]
\Shv \left( \I_{L_2 \subset J}^{\semi}\backslash X \right) \ar[r, "\Ras^{\I_{L_2 \subset J}^{\semi} \mid \I_{L_2}}", "\sim"'] \ar[d, "\oblv^{\I_{L_2 \subset J}^{\semi}|\I_{L_1 \subset J}^{\semi}}"'] & \Shv(\I_{L_2} \backslash X) \ar[d, "\oblv^{\I_{L_2}|\I_{L_1}}"] \\
\Shv \left( \I_{L_1 \subset J}^{\semi}\backslash X \right) \ar[r, "\Ras^{\I_{L_1 \subset J}^{\semi} \mid \I_{L_1}}", "\sim"'] & \Shv(\I_{L_1} \backslash X).
\end{tikzcd}
\end{equation}
\end{prop}

For $L \subset J \subset S$,
we will denote by $\Ras^{\I_L \mid \I_{L \subset J}^{\semi}}$ the inverse of $\Ras^{\I_{L \subset J}^{\semi} \mid \I_L}$.
If $J_1, J_2 \subset S$ and $J_1 \supset L \subset J_2$, we also set
%For any $J_1,J_2\subset S$, we abbreviate\footnote{Can we define $\Ras^{\I_{J_1}^{\quar}\mid \I_{J_2}^{\quar}}$ directly for any $J_1\subset J_2\subset S$, and then prove the compatibility?}
\[
\Ras^{\I_{L \subset J_1}^{\semi} \mid \I_{L \subset J_2}^{\semi}} := \Ras^{\I_L \mid \I_{L \subset J_2}^{\semi}} \circ \Ras^{\I_{L \subset J_1}^{\semi} \mid \I_L} : \Shv \left( \I_{L \subset J_1}^{\semi} \backslash X \right) \simto \Shv \left( \I_{L \subset J_2}^\semi \backslash X \right).
\] 

The proof of Proposition~\ref{prop:Raskin} is similar to that of its counterparts in~\cite{Gaits18,DL23,adr}, so we will not repeat it, but only explain some aspects that will be relevant below.
%Let us explain the construction of the equivalence $\Ras^{\I_{L \subset J}^{\semi} \mid \I_L}$. 
Consider the natural forgetful functors 
\[
\Shv \left( \I_{L \subset J}^{\semi} \backslash X \right) \xrightarrow{\oblv^{\I_{L \subset J}^{\semi} \mid \I_{L \subset J}}} 
\Shv \left( \I_{L \subset J} \backslash X \right) \xleftarrow{\oblv^{\I_L \mid \I_{L \subset J}}}
\Shv \left( \I_L \backslash X \right). 
\]
The functor $\oblv^{\I_L \mid \I_{L \subset J}}$ admits a continuous right adjoint 
\[
\Av^{\I_L \mid \I_{L \subset J}}_*: \Shv(\I_{L \subset J}\backslash X)\rightarrow \Shv(\I_L \backslash X),
\]
and the equivalence $\Ras^{\I_{L \subset J}^{\semi} \mid \I_L}$ is given by 
\[
\Ras^{\I_{L \subset J}^{\semi} \mid \I_L} := \Av^{\I_L \mid \I_{L \subset J}}_*  \circ  \oblv^{\I_{L \subset J}^{\semi} \mid \I_{L \subset J}}. 
\]
In the other direction, the functor $\oblv^{\I_{L \subset J}^{\semi} \mid \I_{L \subset J}}$ admits a (continuous) left adjoint 
\[
\Av_!^{\I_{L \subset J}^{\semi} | \I_{L \subset J}}: \Shv \left( \I_{L \subset J}\backslash X \right) \to \Shv \left( \I_{L \subset J}^{\semi} \backslash X \right),
\]
%which is defined on the image of $\oblv^{\I_{J}\mid \I_{M_J}}$. 
and we have 
\[
\Ras^{\I_L \mid \I_{L \subset J}^{\semi}} = 
\Av^{\I_{L \subset J}^{\semi} \mid \I_{L \subset J}}_! \circ \oblv^{\I_L \mid \I_{L \subset J}}. 
\]
More explicitly, 
let $\theta\in \Xd_J^0 \cap \Xd_{J^\comp}^{++}$, and for $n \geq 0$ set 
\[
\I_{L \subset J,(n)}^{\semi}=\I_{L \subset J}\ltimes \Ad_{z^{-n\theta}}(\sL^+U_{J});
\]
then $( \I_{L \subset J,(n)}^{\semi} : n \geq 0)$ forms an exhaustive filtration of $\I_{L \subset J}^{\semi}$ by subgroup schemes, so that as in~\cite[\S 2.6]{adr} we have
\[
\Shv \left( \I_{L \subset J}^{\semi} \backslash X \right) = \lim_n \Shv \left( \I_{L \subset J, (n)}^{\semi} \backslash X \right).
\]
%recall the subgroup $\I_{J,(n)}^{\quar}$ considered in the proof of Lemma~\ref{lem:whit}.
%=\I_{M_J}\ltimes \Ad_{z^{-n\theta}}(\sL^+U_{J})$ of $\I_{J}^{\quar}$, where $\theta\in \Xd_{J}$ is $J^\comp$-dominant enough. 
For any $\sF\in \Shv(\I_L \backslash X)$, we have 
\begin{equation}
\label{eq:Raskin-formula}
\begin{aligned}
\Ras^{\I_L \mid \I_{L \subset J}^{\semi}}(\sF)
&\simeq \colim_{n} \, \Av^{\I_{L \subset J,(n)}^{\semi} \mid \I_{L \subset J}}_!\circ \oblv^{\I_L \mid \I_{L \subset J}}(\sF) \\
&\simeq \colim_{n} \, z^{-n\theta} \Delta^L_{L,t_{n\theta}} \star^{\I_L} \sF [n\cdot \ell(t_{\theta})]
\end{aligned}
\end{equation}
(where the second identification can be justified e.g.~using~\eqref{eqn:decomposition-IL-J} below).

\begin{rmk}
\phantomsection
\label{rmk:Raskin}
\begin{enumerate}
\item
\label{it:ttheta-minimal}
For any $\theta \in \Xd_J^0$, we have $t_\theta \in {}^J W_\ex^J$. Hence one can consider $\Delta^J_{J,t_\theta}$, and we have $(p_J)_* \Delta_{t_\theta} \simeq \Delta^J_{J,t_\theta}$. In particular, the object $\Delta^L_{L,t_{n\theta}}$ considered in~\eqref{eq:Raskin-formula} indeed makes sense, and for $\sF$ as above we have $\oblv^{\I_L \mid \I}(\Delta^L_{L,t_{n\theta}} \star^{\I_L} \sF) \simeq \Delta_{t_{n\theta}} \star^{\I} \oblv^{\I_L \mid \I}(\sF)$.
\item
\label{it:rmk-Raskin-2}
The functor $\Ras^{\I_{L \subset J}^{\semi} \mid \I_L}$ can also be described as the composition
\begin{multline*}
\Shv \left( \I_{L \subset J}^{\semi} \backslash X \right) \xrightarrow{\oblv^{\I_{L \subset J}^{\semi} \mid \I_{L \subset J} \ltimes \sL^+ U_J}} 
\Shv \left( \I_{L \subset J} \ltimes \sL^+ U_J \backslash X \right) \\
\xrightarrow{\Av_*^{\I_L \mid \I_{L \subset J} \ltimes \sL^+ U_J}}
\Shv \left( \I_L \backslash X \right). 
\end{multline*}
where the first functor is the forgetful functor and the second functor is the right adjoint to the corresponding forgetful functor. In fact, this follows from the fact that the forgetful functor $\Shv \left( \I_{L \subset J} \ltimes \sL^+ U_J \backslash X \right) \to \Shv \left( \I_{L \subset J} \backslash X \right)$ is fully faithful, since $\sL^+ U_J$ is unipotent. Similarly, $\Ras^{\I_L \mid \I_{L \subset J}^{\semi}}$ identifies with the composition
\begin{multline*}
\Shv \left( \I_L \backslash X \right) \xrightarrow{\oblv^{\I_L \mid \I_{L \subset J} \ltimes \sL^+ U_J}}
\Shv \left( \I_{L \subset J} \ltimes \sL^+ U_J \backslash X \right) \\
\xrightarrow{\Av_!^{\I_{L \subset J}^{\semi} \mid \I_{L \subset J} \ltimes \sL^+ U_J}} 
\Shv \left( \I_{L \subset J}^{\semi} \backslash X \right). 
\end{multline*}
\item
\label{it:rmk-Raskin-3}
Denote by $U_J^-$ the unipotent radical of the parabolic subgroup of $G$ opposite to $P_J$ with respect to $T$. Then
a version of the Iwahori decomposition gives that multiplication induces an isomorphism
\begin{equation}
\label{eqn:decomposition-IL-J}
\sL^{++} U_J^- \times \I_{L \subset J} \times \sL^+ U_J \simto \I_L.
\end{equation}
Combining this with the considerations in~\eqref{it:rmk-Raskin-2}, we obtain that for any $\sF$ in $\Shv \left( \I_{L \subset J}^{\semi} \backslash X \right)$ we have
\[
\Ras^{\I_{L \subset J}^{\semi} \mid \I_L}(\sF) \simeq \Av^{\sL^{++} U_J^-}_*(\sF)
\]
where we omit the forgetful functors, and as usual $\Av^{\sL^{++} U_J^-}_*$ is the right adjoint to the corresponding forgetful functor. Similarly, for $\sF$ in $\Shv \left( \I_L \backslash X \right)$ we have
\[
\Ras^{\I_L \mid \I_{L \subset J}^{\semi}}(\sF) \simeq \Av^{\sL^{--} U_J}_!(\sF).
\]
\item
Consider subsets $L \subset J_1 \subset J_2 \subset S$. Then by~\eqref{it:rmk-Raskin-3} $\Ras^{\I_{L \subset J_2}^{\semi} \mid \I_L}$ is the restriction of $\Av^{\sL^{++} U_{J_2}^-}_*$, and $\Ras^{\I_L \mid \I_{L \subset J_1}^{\semi}}$ is the restriction of $\Av^{\sL^{--} U_{J_1}}_!$; as a consequence $\Ras^{\I_{L \subset J_2}^{\semi} \mid \I_{L \subset J_1}^{\semi}}$ is the restriction of $\Av^{\sL^{--} U_{J_1}}_! \circ \Av^{\sL^{++} U_{J_2}^-}_*$. Now we have $U_{J_1} = (B_{J_2} \cap U_{J_1}) \ltimes U_{J_2}$, hence $\sL^{--}U_{J_1} = (\sL^{--} (B_{J_2} \cap U_{J_1})) \ltimes (\sL^{--} U_{J_2})$. Since $\Av_!^{\sL^{--} U_{J_2}} \circ \Av^{\sL^{++} U_{J_2}^-}_* \simeq \id$ on $\Shv \left( \I^{\semi}_{L \subset J_2} \backslash X \right)$, it follows that $\Ras^{\I_{L \subset J_2}^{\semi} \mid \I_{L \subset J_1}^{\semi}}$ is the restriction of $\Av_!^{\sL^{--} (B_{J_2} \cap U_{J_1})}$. Using this one sees that
\[
\Ras^{\I_{L \subset J_2}^{\semi} \mid \I_{L \subset J_1}^{\semi}} = \Av_!^{\I^{\semi}_{L \subset J_1} | \I_{L \subset J_1} \ltimes \sL^+(B_{J_2} \cap U_{J_1}) \ltimes \sL U_{J_2}} \circ \oblv^{\I^{\semi}_{L \subset J_2} | \I_{L \subset J_1} \ltimes \sL^+(B_{J_2} \cap U_{J_1}) \ltimes \sL U_{J_2}}.
\]
\end{enumerate}
\end{rmk}

%Similarly, the equivalence $\Ras^{\I_{J}^{\semi}\mid \I_{J}}$ is given by 
%\[
%\Shv(\I_{J}^{\semi}\backslash X)\xrightarrow{\oblv^{\I_{J}^{\semi}\mid \sL^+{M_J}}} 
%\Shv(\sL^+{M_J}\backslash X) \xrightarrow{\Av_*^{\I_{J}\mid \sL^+{M_J}}}
%\Shv(\I_{J}\backslash X). 
%\]

The following lemma is clear. 

\begin{lem}
\label{lem:Av-Ras}
Let $L \subset J\subset S$, and let $K \subset S_\af$ and $A \subset S_\af$ be finitary. 
%\begin{enumerate}
%\item We have an equivalence of functors from $\Shv(\I_{J}^{\semi}\backslash X)$ to $\Shv(\I\backslash X)$, 
%\[
%\oblv^{\I_J|\I}\circ \Ras^{\I_J^{\semi}\mid \I_J} 
%\simeq \Ras^{\I_J^{\quar}\mid \I}\circ \oblv^{\I_J^{\semi}\mid \I_J^{\quar}}, 
%\]
%and an equivalence of functors from $\Shv(\I\backslash X)$ to $\Shv(\I_{J}^{\semi}\backslash X)$, 
%\[
%\Ras^{\I_J^{\semi}\mid \I_J}\circ \Av_?^{\I_J|\I}
%\simeq \Av_?^{\I_J^{\semi}\mid \I_J^{\quar}}\circ \Ras^{\I_J^{\quar}\mid \I}, \quad ?=!,*.
%\]
\begin{enumerate}
\item
\label{it:Av-Ras-1}
There are canonical isomorphisms
\[
\Ras^{\I_{L \subset J}^\semi \mid \I_L} \circ (p_K)^\dag \simeq (p_K)^\dag \circ \Ras^{\I_{L \subset J}^\semi \mid \I_L}, \quad
\Ras^{\I_{L \subset J}^\semi \mid \I_L} \circ (p_K)_* \simeq (p_K)_* \circ \Ras^{\I_{L \subset J}^\semi \mid \I_L},
\]
and similarly for $q^!$, $q_*$ and $q_!$.
\item 
\label{it:Av-Ras-2}
Assume that $p>0$ and that $\k$ contains a nontrivial $p$-th root of unity. There is a canonical isomorphism
\[
\Ras^{\I_{L \subset J}^{\semi} \mid \I_L} \circ \Av^A_{\chi} \simeq \Av^A_{\chi} \circ \Ras^{\I_{L \subset J}^{\semi} \mid \I_L}, 
%\quad
%\Ras^{\I_{L \subset J}^{\semi} \mid \I_L} \circ \Av^K_{\chi,*} \simeq \Av^K_{\chi,*} \circ \Ras^{\I_{L \subset J}^{\semi} \mid \I_L},
\]
and similarly for $\Av^A_*$ and $\Av^A_!$.
%Suppose $(\F,\K)$ is in situation \eqref{it:sheaves-etale}. We have an equivalence of functors from $\Shv(\I_{J}^{\quar}\backslash \tFl)$ to $\Shv(\I_{J}^{\quar}\backslash \sL G/({}^K\Iu,\chi_K))$, 
%\[
%\Ras^{\I_{J}^{\quar}\mid \I}\circ \Av^K_{\chi_K,!} \simeq \Av^K_{\chi_K,!} \circ \Ras^{\I_{J}^{\quar}\mid \I}.
%\]
%Similar statement holds for the functors $\Av^K_{\chi_K,*}$, $\Av^K_{!}$, $\Av^K_{*}$, $(p_K)^\dag$, $(p_K)_*$ and $q^\dag$, and also
%%We also have similar statements 
%for the functor $\Ras^{\I_{J}^{\semi}\mid \I_J}$. 
\end{enumerate}
\end{lem}

%---------------------------------------------------------------
\subsection{\texorpdfstring{$J$}{J}-Wakimoto sheaves}
\label{sss:Wakimoto}
%---------------------------------------------------------------

Let $J \subset S$, and
recall the subset $W_\ex^{+,J} \subset W_\ex$ defined in~\S\ref{sss:quar-length}. For $x\in W_\ex$, 
we set
\[
\Xd_{J,x}^{0,+} = \{ \mu\in \Xd_J^0\cap \Xd_{J^{\comp}}^+ \mid t_\mu x\in W_\ex^{+,J} \}.
\]
Using~\cite[Lemma~4.1]{adr} one sees that the object
\[
\nabla_{t_{-\theta}} \star^\I \Delta_{t_\theta x} \quad \in \Shv(\I\backslash \Fl)
\]
does not depend on the choice of element $\theta \in \Xd_{J,x}^{0,+}$, up to canonical isomorphism. This complex will be denoted $\sW^J_x$, and called the \textit{$J$-Wakimoto sheaf} associated with $x$. If $J=\varnothing$, the object $\sW_x := \sW^{\varnothing}_x$ is the usual Wakimoto sheaf associated with $x$ (with the conventions of~\cite{adr}). In case $J=S$, one can easily see that $\sW_x^S = \Delta_x$ for any $x \in W_\ex$. 

If $\theta \in \Xd_J^0$, it is clear that $\sW^J_{t_\theta} = \sW_{t_\theta}$, and it is not difficult to check that for any $x \in W_\ex$ we have
\begin{equation}
\label{eqn:convolution-Waki}
\sW_{t_\theta} \star^\I \sW^J_x \simeq \sW^J_{t_\theta x}.
\end{equation}
%For $x\in W_\ex$,
%we define the \textit{$J$-Wakimoto sheaves} 
%\[
%{}^J\sW_x:= \lim_{\theta\in \Xd_J^0} \nabla_{t_{-\theta}} \star^\I \Delta_{t_\theta x} 
%\quad \in \Shv(\I\backslash \Fl),
%\]
%where the limit actually stabilizes. 
%If $J=\varnothing$, the object $\sW_x={}^{\varnothing}\sW_x$ is the usual Wakimoto sheaf (with the conventions of~\cite{adr}). 
Using~\cite[Lemma~4.1(3)]{adr} one sees that $\sW^J_x$ is a perverse sheaf, supported on $\overline{\I \dot{x} \I/\I}$.

When $p>0$ and $\k$ contains a nontrivial $p$-th root of unity, for $x\in W_\ex^{(A,J),\max}$ we also set
\begin{equation*}
%\label{eq:wak-whit-quar}
\sW^J_{\chi_A,x}:= \Av^A_{\chi}(q^! \sW^J_x) \quad \in \Shv(\I \backslash \sL G / (\Iu^A, \chi_A)).
%\simeq \lim_{\theta\in \Xd_J^0} \nabla_{t_{-\theta}} \star^\I \Delta_{t_\theta x}^{\chi_K},
\end{equation*}
%where we use the fact that $t_\theta x\in W_\ex^{K,\mathrm{w}}$ when $\theta\in \Xd_J^0$ is $J^\comp$-dominant enough for any $x\in W_\ex^{(K,J),\mathrm{w}}$. 
Note that if $\theta \in \Xd_J^0 \cap \Xd_{J^\comp}^+$ satisfies $t_\theta x \in W_\ex^{+,J} \cap W_\ex^{A,\max}$, then by~\eqref{eqn:Av-chi-D-classical} we have
\[
\sW^J_{\chi_A,x} = \nabla_{t_{-\theta}} \star^\I \Delta_{\chi_A,t_\theta x}.
\]
Here again, $\sW^J_{\chi_A,x}$ is a perverse sheaf supported on the closure of the orbit parame\-trized by $x$.

\begin{lem}
\label{lem:Av-chi-Wak}
Assume that $p>0$ and that $\k$ contains a nontrivial $p$-th root of unity. Let $J \subset S$, and let $A \subset S_\af$ be finitary.
\begin{enumerate}
\item
\label{it:Av-chi-Wak-1}
Let $x \in W_\ex$, and let $w \in W_A$ be the unique element such that $xw \in W_\ex^{(A,J), \max}$. Then we have
\[
\Av^A_{\chi}(q^! \sW_x^J) \simeq \sW_{\chi_A,xw}^J.
\]
\item
\label{it:Av-chi-Wak-2}
Let $x \in W_\ex^{(A,J),\max}$. The object $\Av^{A}_{!}(\sW^{J}_{\chi_A, x})$ is an iterated extension of objects $q^! \sW^J_{xw}$ with $w\in W_A$, each appearing once.
\end{enumerate}
\end{lem}

\begin{proof}
\eqref{it:Av-chi-Wak-1}
One can choose $\theta \in \Xd_J^0 \cap \Xd_{J^\comp}^+$ such that $t_\theta x \in W_\ex^{+,J}$ and $t_\theta xw \in W_\ex^{+,J} \cap W_\ex^{A,\max}$. Then we have $\sW_x^J \simeq \nabla_{t_{-\theta}} \star^\I \Delta_{t_\theta x}$, $\sW^J_{\chi_A,xw} = \nabla_{t_{-\theta}} \star^\I \Delta_{\chi_A,t_\theta xw}$, and the claim follows from~\eqref{eqn:Av-chi-D-classical}.

\eqref{it:Av-chi-Wak-2}
The proof is similar, using a variant of~\cite[Lemma~3.6]{ar-model} (for $\Fl$ instead of $\Gr$, and with left and right multiplications switched).
\end{proof}

The following statement is analogous to~\cite[Lemma~5.5]{adr}, and follows from similar arguments. (See also the discussion in~\cite[\S 3.1.10]{DL23}.)

\begin{lem}
\label{lem:Ras-shift-Wak}
Let $L \subset J \subset S$, and let $X$ be as above.
For any $\theta \in \Xd_J^0$, we have a canonical equivalence
\[
\Ras^{\I^\semi_{L \subset J} \mid \I_L} ( - \langle \theta \rangle) \simeq ((p_L)_* \sW_{t_\theta}) \star^{\I_L} \Ras^{\I^\semi_{L \subset J} \mid \I_L}(-).
\]
\end{lem}

%\begin{equation}\label{equ wak-quar}
%\sW^{K,\quar}_x:=(p_K)_*(\sW_x)
%\simeq \varprojlim_{\theta\in \check{\X}_J} \nabla_{t_{-\theta}} \star^\I \Delta^K_{t_\theta x}, 
%\quad x\in W_\ex^{(K,J)}, 
%\end{equation}
%\begin{equation}
%{}^J\sW^{\semi}_x:=\Av^{\I_J\mid \I}_!({}^J\sW^{\quar}_x)
%\simeq \varprojlim_{\theta\in \check{\X}_J} {}^J\nabla_{t_{-\theta}} \star^\I \Delta_{t_\theta x}, 
%\quad x\in {}^JW_\ex, 
%\end{equation}
%\begin{equation}\label{equ wak-whit-semi}
%{}^J\sW^{\chi_K,\semi}_x:=\Av^{\I_J\mid \I}_!({}^J\sW^{\chi_K,\quar}_x)
%\simeq \varprojlim_{\theta\in \check{\X}_J} {}^J\nabla_{t_{-\theta}} \star^\I \Delta_{t_\theta x}^{\chi_K}, 
%\quad x\in {}^JW_\ex^{(K,J),\mathrm{w}}, 
%\end{equation}
%where in \eqref{equ wak-quar}, \eqref{eq:wak-whit-quar} and \eqref{equ wak-whit-semi} we use the fact that for $x\in W_\ex^{(K,J)}$, $t_\theta x\in W_\ex^{(K,S)}$ when $\theta\in \check{\X}_J$ is $J^c$-dominant enough. 

We deduce the following result, which generalizes~\cite[Proposition~5.6]{adr}.

\begin{prop}
\label{prop:Ras-D}
Let $J \subset S$.
\begin{enumerate}
\item
\label{it:Ras-D-1}
For any $x \in W_\ex$ we have an equivalence 
\[
\Ras^{\I_J^{\quar}\mid \I} \left( \Delta^{J,\quar}_x \right) \simeq \sW^J_x.
\]
\item
\label{it:Ras-D-2}
Assume that $p>0$ and that $\k$ contains a nontrivial $p$-th root of unity, and let $A \subset S_\af$ be finitary. For any $x \in W_\ex^{(A,J),\max}$ we have an equivalence
\[
\Ras^{\I_J^{\quar}\mid \I} \left( \Delta^{J,\quar}_{\chi_A,x} \right) \simeq \sW^J_{\chi_A,x}. 
\]
\end{enumerate}
%\[
%\Ras^{\I_J^{\quar}\mid \I}({}^J\Delta^{\quar}_x)={}^J\sW^{\quar}_x, \ \
%\Ras^{\I^{\semi}\mid \I}(\Delta^{K,\quar}_x)=\sW^{K,\quar}_x, \ \
%\Ras^{\I_J^{\semi}\mid \I_J}({}^J\Delta^{\semi}_x)={}^J\sW^{\semi}_x, 
%\]
%and 
%\[
%\Ras^{\I_J^{\quar}\mid \I}({}^J\Delta^{\chi_K,\quar}_x)={}^J\sW^{\chi_K,\quar}_x, \quad 
%\Ras^{\I_J^{\semi}\mid \I_J}({}^J\Delta^{\chi_K,\semi}_x)={}^J\sW^{\chi_K,\semi}_x. 
%\]
\end{prop}

\begin{proof}
\eqref{it:Ras-D-1}
Using Lemma~\ref{lem:Ras-shift-Wak} together with~\eqref{eqn:D-N-shift} and~\eqref{eqn:convolution-Waki}, we can assume that $x \in W_\ex^{+,J}$. In this case, what we need to prove is that $\Ras^{\I \mid \I_J^{\quar}} ( \Delta_x) \simeq \Delta^{J,\quar}_x$.
If $\theta$ is as in~\eqref{eq:Raskin-formula}, for any $n \geq 0$ we have $\ell(t_{n\theta} x) = \ell(x)+\ell(t_{n\theta})$, so
\[
\Delta_{t_{n\theta}} \star^{\I} \Delta_x \simeq \Delta_{t_{n\theta} x}
\]
by~\cite[Lemma~4.1(1)]{adr}.
The desired claim follows using~\eqref{eq:Raskin-formula}, since for $n \gg 0$ we have
\[
z^{-n\theta}\I z^{n\theta} \dot{x} \I/\I = \I^{\quar}_{J,(n)} \dot{x} \I/\I.
\]

\eqref{it:Ras-D-2}
The proof is similar.
%We show the second equivalence; the proof of the first one is similar. 
%It is enough to prove the assertion for $x=t_\lambda w$ where $\lambda\in \Xd_J^0$ is $J^\comp$-dominant enough. 
%In this case we have $\Delta_{x}^{\chi_K}={}^J\sW_x^{\chi_K}$, and $\ell(t_{n\theta} x)=\ell(t_{n\theta})+\ell(x)$ for any $n\geq 0$. 
%By \eqref{eq:Raskin-formula} we have 
%\[
%\Ras^{\I\mid \I_J^{\quar}}(\Delta_{x}^{\chi_K})\simeq \colim_{n} \, z^{-n\theta} \Delta_{t_{n\theta}x}^{\chi_K} [n\ell(t_{\theta})] 
%\simeq {}^J\Delta^{\chi_K,\quar}_x, 
%\]
%where the second equivalence follows from the identity 
%\[
%z^{-n\theta}\I z^{n\theta}x\sL^{++}G/\sL^{++}G = \I^{\quar}_{J,(n)}x\sL^{++}G/\sL^{++}G
%%, \quad  \qedhere
%\]
%when $n \gg 0$
\end{proof}

We now prove a generalization of (the case $L=\varnothing$ of) Lemma~\ref{lem:Ras-shift-Wak}.

\begin{lem}
\label{lem:Ras-convolution-Wak}
For any $\lambda\in \Xd$, we have an equivalence of functors from $\Shv(\I^{\quar}_J\backslash X)$ to $\Shv(\I\backslash X)$, 
\begin{equation*}
%\label{equ 2.2}
\Ras^{\I_J^{\quar}\mid \I} \circ (\sW^{M_J}_{t_\lambda} \star^{\I_{\varnothing \subset J}} -) [-\langle \lambda, 2\rho-2\rho_J \rangle]
\simeq \sW_{t_\lambda} \star^{\I} \Ras^{\I_J^{\quar}\mid \I} (-),
\end{equation*}
where $\sW^{M_J}_{t_\lambda}$ is the Wakimoto sheaf in $\Shv(\I_{\varnothing \subset J}\backslash \sL{M}_J/\I_{\varnothing \subset J})$ attached to $\lambda$. 
%In particular, for any $\theta\in \Xd_J^0$ we have 
%\begin{equation}\label{equ 2.3}
%\Ras^{\I_J^{\quar}\mid \I}\circ \langle \theta \rangle \simeq \sW_{t_\theta} \star^{\I} \Ras^{\I_J^{\quar}\mid \I} (-). 
%\end{equation}
\end{lem}

\begin{proof}
%Since $\sW_{t_\lambda}$ is invertible in $\Shv(\I\backslash \sL{G}/\I)$ with inverse $\sW_{t_{-\lambda}}$, and similarly for $\sW^{M_J}_{t_\lambda}$, it is enough to prove the claim when $\lambda$ is dominant. 
Since every coweight is a difference of two dominant coweights, and both sides are compatible with convolution in the appropriate way, invertibility of the Wakimoto sheaves in $\Shv(\I\backslash \sL{G}/\I)$ and $\Shv(\I_{\varnothing \subset J}\backslash \sL{M}_J/\I_{\varnothing \subset J})$ reduces the proof to the case $\lambda$ is dominant.
In this case $\sW_{t_\lambda}=\Delta_{t_\lambda}=\sW_{t_\lambda}^J$, and $\sW^{M_J}_{t_\lambda}$ is the standard object $\Delta^{M_J}_{t_\lambda}$ in $\Shv(\I_{\varnothing \subset J}\backslash \sL{M}_J/\I_{\varnothing \subset J})$ attached to $\lambda$. 
Once we have fixed an element $\theta$ as in~\eqref{eq:Raskin-formula},
by Proposition~\ref{prop:Ras-D} and~\eqref{eq:Raskin-formula} it is enough to show that there is an equivalence 
%of objects in $\Shv(\I^{\quar}_J\backslash \Fl)$, 
\[
\colim_{n} \, \Delta^{M_J}_{t_\lambda} \star^{\I_{\varnothing \subset J}} z^{-n\theta} \cdot \Delta_{t_{n\theta}} [n\ell(t_\theta) -\langle \lambda, 2\rho-2\rho_J \rangle] \simeq \Delta^{J,\quar}_{t_\lambda} . 
\]
This follows from the isomorphism
\[
 \I_{\varnothing \subset J}z^\lambda \I_{\varnothing \subset J} \times^{\I_{\varnothing \subset J}} z^{-n\theta} \I z^{n\theta}\I/\I 
\xs \big(\I_{\varnothing \subset J}\ltimes \Ad_{z^{\lambda-n\theta}}(\sL^+U_{J})\big) z^{\lambda} \I/\I
%. \qedhere 
\]
induced by multiplication for $n \gg 0$.
\end{proof}

%---------------------------------------------------------------
\subsection{Properties of averaging functors} 
\label{sss:prop-av}
%---------------------------------------------------------------

As an application of Raskin's equivalence, we obtain some properties of the averaging functors defined in \S\ref{sss:av-fun}. We assume that $p>0$ and that $\k$ contains a nontrivial $p$-th root of unity.
%Suppose $(\F,\K)$ is in situation \eqref{it:sheaves-etale}. 

%\begin{lem}
%Let $A=\I_J^\quar$ or $\I_J^\semi$. 
%The natural morphism of functors from $\Shv(A\backslash \tFl)$ to $\Shv(A\backslash \sL G/ ({}^K\Iu,\chi_K))$ 
%\[
%\Av^K_{\chi_K,!} \rightarrow \Av^K_{\chi_K,*} 
%\]
%is an equivalence. 
%\end{lem}
%
%\begin{proof}
%By Lemma~\ref{lem:Av-Ras} it is enough to prove a similar equivalence of functors from $\Shv(A'\backslash \tFl)$ to $\Shv(A'\backslash \sL G/ ({}^K\Iu,\chi_K))$, for $A'=\I$ or $\I_J$. 
%Then this is standard, see e.g.~\cite[Lemma~3.2]{AR23}. 
%(The proof in \textit{loc. cit.} is written in the framework of ordinary bounded derived category of constructible sheaves, but the same arguments apply in our present setting.)
%%we give some details as follows to adapt to our setting. 
%%Firstly we consider the $\infty$-category of monodromy sheaves on $\tFl$, then the functor $\Av^K_{\chi_K,!}$, resp. $\Av^K_{\chi_K,*}$, is equivalent to convolution by certain free monodromy standard sheaf, resp. costandard sheaf. 
%%One uses the fact that the involved standard sheaf and costandard sheaf coincide, thus implies the equivalence of functors. 
%%It then induces an equivalence of functors on the desired categories, since the natural functor to the category of monodromy sheaves is conservative. 
%\end{proof}
%%\noindent 
%
%From now on, because of this lemma we will set
%\[
%\Av^{K}_{\chi_K}:=\Av^K_{\chi_K,!} \simeq \Av^K_{\chi_K,*}.
%\]

\begin{lem}
\label{lem:Av-chi-D}
Let $L \subset J \subset S$, and let $A \subset S_\af$ be finitary.

\begin{enumerate}
\item
\label{it:Av-chi-D-1}
%\label{it:Av-pullback-q-1}
%For any $x\in W_\ex$ we have 
%\[
%\Av^{K}_{\chi_K}(q^\dag ({}^J\Delta^{\quar}_{x}))\simeq {}^J\Delta^{\chi_K,\quar}_{xw},  
%\]
%where $w\in W_K$ such that $xw\in W_\ex^{(K,J),\mathrm{w}}$. 
%%Similar statement holds for ${}^J\nabla^{\quar}_{x}$. 
%%\[
%%\Av^{K}_{\chi_K}({}^J\Delta^{\quar}_{xw})\simeq {}^J\Delta^{\chi_K,\quar}_{x}, \quad 
%%\Av^{K}_{\chi_K}({}^J\nabla^{\quar}_{xw})\simeq {}^J\nabla^{\chi_K,\quar}_{x}. 
%%\]
%\item
%\label{it:Av-pullback-q-2}
For any $x\in {}^L W_\ex$ we have 
\[
\Av^{A}_{\chi}(q^! (\Delta^{L \subset J, \semi}_{x})) \simeq 
\begin{cases}
\Delta^{L \subset J,\semi}_{\chi_A, xw} & \text{if $w\in W_A$ satisfies $xw\in {}^L W_\ex^{(A,J),\max}$;} \\
0 & \text{if $x W_A \cap {}^L W_\ex^{(A,J),\max} = \varnothing$.} 
\end{cases}  
\]
%Similar statement holds for ${}^J\nabla^{\semi}_{x}$. 
\item 
\label{it:Av-chi-D-2}
For any $x\in {}^L W_\ex^{(A,J),\max}$, the object $\Av^{A}_{!}(\Delta^{L \subset J,\semi}_{\chi_A,x})$
%, resp. $\Av^{K}_{*}({}^J\nabla^{\chi_K,\semi}_{x})$, 
is an iterated extension of objects $q^! (\Delta^{L \subset J, \semi}_{xw})$ with $w\in W_A$, 
%resp. ${}^J\nabla^{\semi}_{xw}$ ($w\in W_K$), 
each appearing once.
%For $x\in {}^JW_\ex^{(K,S)}$ and $w\in W_K$ such that $xw\in {}^JW_\ex$, we have 
%\[
%\Av^{K}_{\chi_K}({}^J\Delta_{xw})\simeq {}^J\Delta^{\chi_K}_{x}, \quad 
%\Av^{K}_{\chi_K}({}^J\nabla_{xw})\simeq {}^J\nabla^{\chi_K}_{x}. 
%\]
\end{enumerate}
\end{lem}

\begin{proof}
\eqref{it:Av-chi-D-1}
First, assume that $L=\varnothing$. (In this case, an element $w$ as in the statement always exists.) Then the claim follows from
Lemma~\ref{lem:Av-Ras}, Lemma~\ref{lem:Av-chi-Wak}\eqref{it:Av-chi-Wak-1} and Proposition~\ref{prop:Ras-D}. The general case follows, using Lemma~\ref{lem:Av-quar-semi-D-N}\eqref{it:Av-quar-semi-D-N-1} (applied for the subset $K=\varnothing$), the commutation of $\Av_!^{\I_{L \subset J}^\semi \mid \I_J^\quar}$ with $q^!$ and $\Av^{A}_{\chi}$, and Lemma~\ref{lem:av-oblv-whit}\eqref{it:av-oblv-whit-1}.
%\eqref{it:Av-pullback-q-2}
%This follows from the commutation of $\Av^{K}_{\chi_K}$ and $\Av_!^{\I_J^\quar\mid \I_J^\semi}$, \eqref{it:Av-pullback-q-1}, and Lemmas~\ref{lem:Av-quar-semi-D-N}\eqref{it:Av-quar-semi-D-N-1} and \ref{lem:av-oblv-whit}\eqref{it:av-oblv-whit-1}. 

\eqref{it:Av-chi-D-2}
The proof is similar, using Lemma~\ref{lem:Av-chi-Wak} for the case $L=\varnothing$, and then Lem\-ma~\ref{lem:av-oblv-whit}\eqref{it:av-oblv-whit-1} and Lemma~\ref{lem:Av-quar-semi-D-N}\eqref{it:Av-quar-semi-D-N-1} (with $K = \varnothing$).
\end{proof}

%\begin{lem}
%\label{lem:Av-D}
%Let $L \subset J \subset S$, and let $K \subset S_\af$ be finitary.
%%\begin{enumerate}
%%\item 
%%\label{it:Av-D-1}
%%For any $x\in W_\ex^{(K,J),\mathrm{w}}$, the object $\Av^{K}_{!}({}^J\Delta^{\chi_K,\quar}_{x})$
%%%, resp. $\Av^{K}_{*}({}^J\nabla^{\chi_K,\quar}_{x})$, 
%%is an iterated extension of objects $q^\dag ({}^J\Delta^{\quar}_{xw})$ ($w\in W_K$), 
%%%resp. ${}^J\nabla^{\quar}_{xw}$ ($w\in W_K$), 
%%each appearing once. 
%%\item
%%\label{it:Av-D-2}
%For any $x\in {}^L W_\ex^{(K,J),\max}$, the object $\Av^{K}_{!}(\Delta^{L \subset J,\semi}_{\chi_K,x})$
%%, resp. $\Av^{K}_{*}({}^J\nabla^{\chi_K,\semi}_{x})$, 
%is an iterated extension of objects $q^* (\Delta^{L \subset J, \semi}_{xw})$ with $w\in W_K$, 
%%resp. ${}^J\nabla^{\semi}_{xw}$ ($w\in W_K$), 
%each appearing once. 
%%\end{enumerate}
%\end{lem}

%\begin{proof}
%\eqref{it:Av-D-1}
%By Lemma~\ref{lem:Av-Ras} and Proposition~\ref{prop:Ras-D}, it suffices to show that $\Av^{K}_{!}({}^J\sW^{\chi_K}_{x})$ is an iterated extension of objects ${}^J\sW_{xw}$ ($w\in W_K$), each appearing once. 
%Now we have 
%\[
%\Av^{K}_{!}({}^J\sW^{\chi_K}_{x}) \simeq \lim_{\theta\in \Xd_J^0} \nabla_{t_{-\theta}} \star^\I \Av^{K}_{!}(\Delta_{t_\theta x}^{\chi_K}).
%\]
%Note that $t_\theta x\in W_\ex^{K,\mathrm{w}}$ when $\theta\in \Xd_J^0$ is $J^\comp$-dominant enough. 
%In this case $\Av^{K}_{!}(\Delta_{t_\theta x}^{\chi_K})$ is an iterated extension of objects $q^\dag(\Delta_{t_\theta xw})$ ($w\in W_K$), each appearing once, which proves the desired claim. 
%
%\eqref{it:Av-D-2}
%This follows from~\eqref{it:Av-D-1}, the commutation of $\Av^{K}_{!}$ and $\Av_!^{\I_J^\semi\mid \I_J^\quar}$, and Lemmas~\ref{lem:Av-quar-semi-D-N}\eqref{it:Av-quar-semi-D-N-1} and \ref{lem:av-oblv-whit}\eqref{it:av-oblv-whit-1}. 
%\end{proof}

%%%%%%%%%%%%%%%%%%%%%%%%%
\section{Perverse t-structure} 
\label{sec:perv}
%%%%%%%%%%%%%%%%%%%%%%%%%

Let $X$ be either $\Fl_K$ for some finitary $K \subset S_\af$, $\tFl$, or the symbol $\sL G/(\Iu^A,\chi_A)$ for some finitary $A \subset S_\af$ (assuming that $p>0$ and that $\k$ contains a nontrivial $p$-th root of unity).
%(when $(\F,\K)$ is in situation \eqref{it:sheaves-etale}). 

%---------------------------------------------------------------
\subsection{Definition} 
%---------------------------------------------------------------

%Recall that there is a standard perverse t-structure on the $\infty$-category $\Shv(\I_{L}\backslash X)$. 
If $L \subset J \subset S$, we define a t-structure, called the \textit{perverse t-structure}, on $\Shv\left(\I_{L\subset J}^{\semi}\backslash X\right)$, such that 
\[
\Shv\left(\I_{L\subset J}^{\semi}\backslash X\right)^{\leq 0}  
\]
is the smallest full subcategory that is stable under extensions and colimits and contains the following collection of objects: 
\begin{itemize}
\item $\left( \Delta^{L\subset J,\semi}_{K,x} : x\in {}^LW_\ex^{(K,J)} \right)$ when $X=\Fl_K$; 
\item $\left( q^!\Delta^{L\subset J,\semi}_x : x\in {}^LW_\ex \right)$ when $X=\tFl$; 
\item $\left( \Delta^{L\subset J,\semi}_{\chi_A,x} : x\in {}^LW_\ex^{(A,J),\max} \right)$ when $X = \sL G/(\Iu^A,\chi_A)$.
%(assuming that $p>0$ and that $\k$ contains a nontrivial $p$-th root of unity); 
\end{itemize}
Note that when $J=S$ (so that $\I_{L \subset S}^\semi=\I_L$) and $X=\Fl_K$, this t-structure is the natural one in the sense that the forgetful functor $\oblv^{\I_L} : \Shv(\I_L \backslash \Fl_K) \to \Shv(\Fl_K)$ is t-exact with respect to our perverse t-structure on the domain and the usual perverse t-structure (of middle perversity) on the codomain.

As in~\cite[\S 6.1]{adr}, the nonnegative part $\Shv\left(\I_{L\subset J}^{\semi}\backslash X\right)^{\geq 0}$ can be described as the right orthogonal to $\Shv\left(\I_{L\subset J}^{\semi}\backslash X\right)^{\leq -1}$, or equivalently in terms of $!$-restriction to strata.

As usual, we will say that an object of $\Shv\left(\I_{L\subset J}^{\semi}\backslash X\right)$ \emph{is perverse} to say that it belongs to the heart of this t-structure.

%---------------------------------------------------------------
\subsection{Some t-exactness results}
%---------------------------------------------------------------

We continue with $X$ as above.

\begin{lem}
\label{lem:oblv-semi-quar-texact}
Let $L_1\subset L_2\subset J\subset S$. 
The functor 
\[
\oblv^{\I_{L_2 \subset J}^{\semi}|\I_{L_1 \subset J}^{\semi}}:
\Shv\left(\I_{L_2 \subset J}^{\semi}\backslash X\right)\rightarrow \Shv\left(\I_{L_1 \subset J}^{\semi}\backslash X\right)
\]
is t-exact. 
\end{lem}

\begin{proof}
We explain the proof in case $X=\Fl_K$; the other cases are similar. (For $X=\sL G/(\Iu^A,\chi_A)$, the argument uses Lemma~\ref{lem:av-oblv-whit}.)
We first assume that $L_1=\varnothing$. 
To show right t-exactness of our functor, it is enough to show that $\oblv^{\I_{L_2 \subset J}^{\semi}|\I_{J}^{\quar}}(\Delta^{L_2\subset J,\semi}_{K,x})$ is contained in $\Shv\left(\I_{J}^{\quar}\backslash X\right)^{\leq 0}$ for any $x\in {}^{L_2}W_\ex^{(K,J)}$, which follows from Lemma~\ref{lem:Av-quar-semi-D-N}\eqref{it:Av-quar-semi-D-N-2}. 
We next show left t-exactness. If $\sF\in \Shv\left(\I_{L_2 \subset J}^{\semi}\backslash X\right)^{\geq 0}$, then for any $x\in W_\ex^{(K,J)}$ and $n> 0$ we have 
\begin{equation*}
%\label{eq:Raskin-formula3}
\Hom \left( \Delta^{J,\quar}_{K,x}[n], \oblv^{\I_{L_2 \subset J}^{\semi}|\I_{J}^{\quar}}(\sF) \right) \simeq \Hom \left( \Av_!^{\I_{L_2 \subset J}^{\semi}|\I_{J}^{\quar}}(\Delta^{J,\quar}_{K,x})[n],\sF \right). 
\end{equation*}
By Lemma~\ref{lem:Av-quar-semi-D-N}\eqref{it:Av-quar-semi-D-N-1}, the right-hand side
%RHS of \eqref{eq:Raskin-formula3} 
vanishes, which implies that $\oblv^{\I_{L_2 \subset J}^{\semi}|\I_{J}^{\quar}}(\sF)\in \Shv\left(\I_J^{\quar}\backslash X\right)^{\geq 0}$ as desired.
%The proof of the other cases is similar, where in the case $X=\sL G/({}^K\Iu,\chi_K)$, one applies Lemma~\ref{lem:av-oblv-whit}. 

Now let $L_1$ be general. 
It is clear that we have a natural equivalence of functors 
\[
\oblv^{\I_{L_2 \subset J}^{\semi}|\I_{J}^{\quar}} \simeq 
\oblv^{\I_{L_1 \subset J}^{\semi}|\I_{J}^{\quar}} \circ \oblv^{\I_{L_2 \subset J}^{\semi}|\I_{L_1 \subset J}^{\semi}}. 
\]
For any $L\subset J$ the functor $\oblv^{\I_{L \subset J}^{\semi}|\I_{J}^{\quar}}$ is conservative, and is t-exact by our previous discussion. 
Hence the functor $\oblv^{\I_{L_2 \subset J}^{\semi}|\I_{L_1 \subset J}^{\semi}}$ is t-exact. 
\end{proof}

\begin{lem}
\label{lem:pullback-Av-texact}
Let $L\subset J\subset S$, and let $K\subset S_\af$ and $A \subset S_\af$ be finitary. 
\begin{enumerate}
\item 
\label{it:pullback-q-texact}
The functor $q^!: \Shv\left(\I_{L\subset J}^{\semi}\backslash \Fl\right)\rightarrow \Shv\left(\I_{L\subset J}^{\semi}\backslash \tFl\right)$ is t-exact. 
Moreover, its restriction to the hearts of the perverse t-structures is fully faithful, and its essential image is stable under subquotients. 
\item
\label{it:pullback-p-texact}
The functor $(p_K)^\dag:\Shv\left(\I_{L\subset J}^{\semi}\backslash \Fl_K\right)\rightarrow \Shv\left(\I_{L\subset J}^{\semi}\backslash \Fl\right)$ is t-exact. 
Moreover, its restriction to the hearts of the perverse t-structures is fully faithful, and its essential image is stable under subquotients and extensions. 
\item
\label{it:Av-texact}
If $p>0$ and $\k$ contains a nontrivial $p$-th root of unity,
the functor $\Av^A_{\chi}:\Shv\left(\I_{L\subset J}^{\semi}\backslash \tFl\right)\rightarrow \Shv\left(\I_{L\subset J}^{\semi}\backslash \sL G/(\Iu^A,\chi_A)\right)$ is t-exact. 
%, and so are $\Av^K_!$ and $\Av^K_*$. 
\end{enumerate}
\end{lem}

\begin{proof}
The proofs of~\eqref{it:pullback-q-texact} and~\eqref{it:pullback-p-texact} are similar to those of~\cite[Lemmas 6.6 \& 6.8]{adr}. (In particular, the proof of~\eqref{it:pullback-p-texact} uses Lemma~\ref{lem:push-pull-pK-D-N-semi}).  
The statement in~\eqref{it:Av-texact} follows from similar arguments as those used in the proof of Lemma~\ref{lem:oblv-semi-quar-texact}, using Lemma~\ref{lem:Av-chi-D}. 
%The case $\mathrm{J}=\I_J^{\semi}$ can be reduced to the former one, by the commutativity of $\oblv^{\I_J^{\semi}|\I_J^{\quar}}$ and $\Av^K_{\chi_K}$, and the fact that $\oblv^{\I_J^{\semi}|\I_J^{\quar}}$ is conservative and t-exact (by Lemma~\ref{lem:oblv-semi-quar-texact}). 
\end{proof}


%---------------------------------------------------------------
\subsection{Perversity of standard sheaves} 
%---------------------------------------------------------------

We now establish an analogue of~\cite[Lem\-ma~6.9]{adr}, proving that the standard objects are in the heart of the perverse t-structure in some cases. (This property is not satisfied in general in the other cases; e.g.~in $\Shv(\I_L \backslash \Fl_K)$ with $L$ and $K$ nonempty the orbits are not necessarily affine, so that the standard objects are not perverse in general.)

\begin{lem}
\label{lem:standard-perverse}
Let $L\subset J\subset S$, and let $K\subset S_\af$ and $A \subset S_\af$ be finitary. 
The objects $\Delta^{J,\quar}_{K,x}$ for $x \in W_\ex^{(K,J)}$, $\Delta^{L\subset J,\semi}_{x}$ for $x \in {}^L W_\ex$, and $\Delta^{L\subset J,\semi}_{\chi_A,x}$ for $x \in {}^L W_\ex^{(A,J),\max}$ (when $p>0$ and $\k$ contains a nontrivial $p$-th root of unity) are contained in the heart of the perverse t-structure on the corresponding $\infty$-categories. 
\end{lem}

\begin{proof}
First, we consider the case of the objects $\Delta^{J,\quar}_{K,x}$ ($x \in W_\ex^{(K,J)}$). It is enough to show that these objects belong to the ``$\geq 0$'' part of the t-structure, i.e.~that for any $x,y\in W_\ex^{(K,J)}$ and $n>0$ we have 
\begin{equation*}
%\label{eq:Raskin-formula5}
\Hom(\Delta^{J,\quar}_{K,x}[n],\Delta^{J,\quar}_{K,y})=0. 
\end{equation*}

Let $x \in W_\ex^{(K,J)}$.
By Proposition~\ref{prop:Ras-D}\eqref{it:Ras-D-1}, we have $\Ras^{\I_{J}^{\quar}\mid \I}(\Delta^{J,\quar}_{x})\simeq \sW^{J}_{x}$. Using Lemma~\ref{lem:push-pull-pK-D-N-semi}\eqref{it:push-pull-pK-D-N-semi-1} and Lemma~\ref{lem:Av-Ras}\eqref{it:Av-Ras-1} we deduce that
\[
\Ras^{\I_{J}^{\quar}\mid \I}(\Delta^{J,\quar}_{K,x})\simeq \Ras^{\I_{J}^{\quar}\mid \I}((p_K)_* \Delta^{J,\quar}_{x}) \simeq (p_K)_* \sW^{J}_{x} \simeq \nabla_{t_{-\theta}} \star^{\I} \Delta_{K,t_\theta x}
\]
where $\theta \in \Xd_J^0 \cap \Xd_{J^\comp}^+$ satisfies $t_\theta x \in W_\ex^{+,J} \cap W_\ex^{K}$. If now we consider another element $y \in W_\ex^{(K,J)}$, we can choose $\theta \in \Xd_J^0 \cap \Xd_{J^\comp}^+$ such that $t_\theta x$ and $t_\theta y$ are both in $W_\ex^{+,J} \cap W_\ex^{K}$, and we obtain that
\begin{multline*}
\Hom(\Delta^{J,\quar}_{K,x}[n],\Delta^{J,\quar}_{K,y}) \simeq \Hom(\nabla_{t_{-\theta}} \star^{\I} \Delta_{K,t_\theta x} [n],\nabla_{t_{-\theta}} \star^{\I} \Delta_{K,t_\theta y}) \\
\simeq \Hom(\Delta_{K,t_\theta x} [n], \Delta_{K,t_\theta y})
\end{multline*}
by invertibility of the object $\nabla_{t_{-\theta}}$. We deduce the desired vanishing, since the objects $\Delta_{K,t_\theta x}$ and $\Delta_{K,t_\theta y}$ are perverse.

Now, we treat the objects $\Delta^{L\subset J,\semi}_{x}$ ($x \in {}^L W_\ex$). Here again, what we have to show is that
\[
\Hom(\Delta^{L \subset J,\semi}_{x}[n],\Delta^{L \subset J,\semi}_{y})=0
\]
for $x,y \in {}^L W_\ex$ and $n>0$.

Let $x \in {}^L W_\ex$. By Lemma~\ref{lem:Av-quar-semi-D-N}\eqref{it:Av-quar-semi-D-N-1}, the commutative diagram obtained from~\eqref{eq:comm-diag-ras-oblv} by passing to left adjoints, and Proposition~\ref{prop:Ras-D}\eqref{it:Ras-D-1},
%and \ref{lem:Av-Ras} 
we have equivalences 
\begin{multline*}
\Ras^{\I_{L\subset J}^{\semi}\mid \I_L}(\Delta^{L\subset J,\semi}_{x}) \simeq \Ras^{\I_{L\subset J}^{\semi}\mid \I_L}(\Av_!^{\I^\semi_{L \subset J} \mid \I_J^\quar}(\Delta^{J,\quar}_{x}))[-\ell(w_L)] \\
\simeq \Av^{\I_L\mid \I}_! (\sW^{J}_{x})[-\ell(w_L)] \simeq \Av^{\I_L\mid \I}_!(\nabla_{t_{-\theta}}) \star^\I \Delta_{t_\theta x} [-\ell(w_L)] 
\end{multline*}
where $\theta \in \Xd_J^0 \cap \Xd_{J^\comp}^+$ satisfies $t_\theta x \in W_\ex^{+,J}$.
Now, using Remark~\ref{rmk:Raskin}\eqref{it:ttheta-minimal} one sees that we have canonical isomorphisms
\[
\I_L \times^{\I} \I z^\theta \I / \I \simeq \I_L z^\theta \I/\I = \I_L z^\theta \I_L / \I \simeq \I_L z^\theta \I_L \times^{\I_L} \I_L/\I,
\]
from which we obtain an equivalence
\begin{equation}
\label{eqn:computation-standard-perverse}
\Av^{\I_L\mid \I}_!(\nabla_{t_{-\theta}}) \star^\I \Delta_{t_\theta x} [-\ell(w_L)] \simeq \nabla^L_{L,t_{-\theta}} \star^{\I_L} \Av^{\I_L\mid \I}_!(\Delta_{t_\theta x}) [-\ell(w_L)].
\end{equation}
Using once again Lemma~\ref{lem:Av-quar-semi-D-N}\eqref{it:Av-quar-semi-D-N-1} (for the subset ``$J$'' equal to $S$) we deduce an equivalence
\[
\Av^{\I_L\mid \I}_!(\nabla_{t_{-\theta}}) \star^\I \Delta_{t_\theta x} [-\ell(w_L)] \simeq \nabla^L_{L,t_{-\theta}} \star^{\I_L} \Delta^L_{t_\theta x}.
\]
(Note that here we have $t_\theta x \in {}^L W_\ex$ by Lemma~\ref{lem:min-length}\eqref{it:min-length-0}.)
If now $y \in {}^L W_\ex$ is another element, and if $\theta \in \Xd_J^0 \cap \Xd_{J^\comp}^+$ is such that $t_\theta x$ and $t_\theta y$ are both in $W_\ex^{+,J}$, then we obtain that 
\[
\Hom(\Delta^{L \subset J,\semi}_{x}[n],\Delta^{L \subset J,\semi}_{y})=\Hom(\nabla^L_{L,t_{-\theta}} \star^{\I_L} \Delta^L_{t_\theta x}[n],\nabla^L_{L,t_{-\theta}} \star^{\I_L} \Delta^L_{t_\theta y}).
\]
As in~\cite[Lemma~3.1.5(i)]{DL23}, the object $\nabla^L_{L,t_{-\theta}}$ is invertible in $\Shv(\I_L \backslash \sL G / \I_L)$, so that
\[
\Hom(\nabla^L_{L,t_{-\theta}} \star^{\I_L} \Delta^L_{t_\theta x}[n],\nabla^L_{L,t_{-\theta}} \star^{\I_L} \Delta^L_{t_\theta y}) \simeq \Hom(\Delta^L_{t_\theta x}[n],\Delta^L_{t_\theta y}).
\]
One then concludes as above using the fact that $\Delta^L_{t_\theta x}$ and $\Delta^L_{t_\theta y}$ are perverse.

Finally, the case of the objects $\Delta^{L\subset J,\semi}_{\chi_A,x}$ ($x \in {}^L W_\ex^{(A,J),\max}$) follows from the preceding case, using Lemma~\ref{lem:pullback-Av-texact}\eqref{it:pullback-q-texact}--\eqref{it:Av-texact} and Lemma~\ref{lem:Av-chi-D}\eqref{it:Av-chi-D-1}.
%
%We explain the proof for the case of the objects $\Delta^{L\subset J,\semi}_{\chi_A,x}$; the other cases can be treated similarly. 
%It is enough to show that these objects belong to the ``$\geq 0$'' part of the t-structure, i.e.~that for any $x,y\in {}^LW_\ex^{(A,J),\max}$ and $n>0$ we have 
%\begin{equation}
%\label{eq:Raskin-formula5}
%\Hom(\Delta^{L\subset J,\semi}_{\chi_A,x}[n],\Delta^{L\subset J,\semi}_{\chi_A,y})=0. 
%\end{equation}
%To show this, we will show that the object $\Ras^{\I_{L\subset J}^{\semi}\mid \I_L}(\Delta^{L\subset J,\semi}_{\chi_A,x})$ is perverse for any $x \in {}^LW_\ex^{(A,J),\max}$.
%
%By Proposition~\ref{prop:Ras-D}, we have $\Ras^{\I_{J}^{\quar}\mid \I}(\Delta^{J,\quar}_{\chi_A,x})\simeq \sW^{J}_{\chi_A,x}$. 
%By Lemma~\ref{lem:av-oblv-whit}, the commutative diagram obtained from~\eqref{eq:comm-diag-ras-oblv} by passing to left adjoints, and Proposition~\ref{prop:Ras-D}\eqref{it:Ras-D-2},
%%and \ref{lem:Av-Ras} 
%we deduce equivalences 
%\begin{align*}
%\Ras^{\I_{L\subset J}^{\semi}\mid \I_L}(\Delta^{L\subset J,\semi}_{\chi_A,x}) &\simeq \Ras^{\I_{L\subset J}^{\semi}\mid \I_L}(\Av_!^{\I^\semi_{L \subset J} \mid \I_J^\quar}(\Delta^{J,\quar}_{\chi_A,x}))[-\ell(w_L)] \\
%&\simeq \Av^{\I_L\mid \I}_! (\sW^{J}_{\chi_A,x})[-\ell(w_L)] \\
%&\simeq \Av^{\I_L\mid \I}_!(\nabla_{t_{-\theta}}) \star^\I \Delta_{\chi_A,t_\theta x} \ [-\ell(w_L)] 
%\end{align*}
%where $\theta \in \Xd_J^0 \cap \Xd_{J^\comp}^+$ satisfies $t_\theta x \in W_\ex^{+,J} \cap W_\ex^{A,\max}$. Now we have an equivalence $\Av^{\I_L\mid \I}_!\simeq \Av^{\I_L\mid \I}_*[2\ell(w_L)]$ (see the discussion in~\S\ref{ss:qs-sheaves}), and $t_{-\theta} \in {}^L W_\ex$ (see Remark~\ref{rmk:Raskin}\eqref{it:ttheta-minimal}), hence we obtain an equivalence
%\[
%\Ras^{\I_{L\subset J}^{\semi}\mid \I_L}(\Delta^{L\subset J,\semi}_{\chi_A,x})
%\simeq \nabla^L_{t_{-\theta}} \star^\I \Delta_{\chi_A,t_\theta x}.
%\]
%%where the last equivalence we use the fact that $\Av^{\I_J\mid \I}_!\simeq \Av^{\I_J\mid \I}_*[2\ell(w_J)]$, and the limits above stabilize. 
%By a standard argument (as in~\cite[Lemma~4.1(3)]{adr}), the complex $\nabla^L_{t_{-\theta}} \star^\I \Delta_{\chi_A,t_\theta x}$ is perverse in the $\infty$-category $\Shv\left(\I_L\backslash \sL{G}/ (\Iu^A,\chi_A)\right)$, which finishes the proof.
%%
%%Using this claim also for $y$, and the fact that $\Ras^{\I_{L\subset J}^{\semi}\mid \I_L}$ is an equivalence, we deduce~\eqref{eq:Raskin-formula5}. 
\end{proof}

%-----------------------------------------------------------
\subsection{Finitely generated complexes}
%-----------------------------------------------------------

%===============================
\subsubsection{Restriction of the t-structure to finitely generated complexes} 
\label{sss:restriction-fg}
%===============================

For $L\subset J\subset S$ and $X$ as above,
let us now consider the small stable subcategory
\[
%\Shvfg(\I \backslash X) \subset \Shv(\I \backslash X), \quad
\Shvfg(\I_L\backslash X) \subset \Shv(\I_L\backslash X)
\]
of ``finitely generated'' complexes, see the discussions in~\cite[\S 2.2 \& \S 2.5]{adr}. (The homotopy category of $\Shvfg(\I_L\backslash X)$ is what is usually called the $\I_L$-equivariant derived category of constructible sheaves on $X$, see e.g.~\cite{ARbook, ab, by}.)
We will denote by 
\[
%\Shv_{\omega}(\I_J^{\quar}\backslash X), \quad 
\Shv_{\omega}\left(\I_{L\subset J}^{\semi}\backslash X\right) \subset \Shv\left(\I_{L\subset J}^{\semi}\backslash X\right)
\]
the image of $\Shvfg(\I_L\backslash X)$ under $\Ras^{\I_L | \I^\semi_{L \subset J}}$. 
We think of this subcategory as consisting of the ``constructible'' objects.
As in~\cite[\S 6.4]{adr}, the objects in $\Shv_{\omega}\left(\I_{L\subset J}^{\semi}\backslash X\right)$ admit the following characterization.

\begin{lem}
\label{lem:constr-complx}
An object $\sF$ of $\Shv\left(\I_{L\subset J}^{\semi}\backslash \Fl_K\right)$ is contained in $\Shv_{\omega}\left(\I_{L\subset J}^{\semi}\backslash \Fl_K\right)$ if and only if 
\begin{enumerate}
\item $\sF$ is supported on a finite union of closures of orbits $\fS^{L\subset J}_{K,x}$; 
\item the set $\{x\in {}^LW_\ex^{(K,J)}\mid (\bs^{L\subset J}_{K,x})^* \sF\neq 0\}$ is finite; 
\item for any $x\in {}^LW_\ex^{(K,J)}$, the image of $(\bs^{L\subset J}_{K,x})^* \sF$ under the equivalence~\eqref{eq:semi-orbt} belongs to $\Shvc \left( \mathrm{Stab}_{M_L}(\dot{x}\I_K/\I_K)\backslash \pt \right)$. 
\end{enumerate}
We have a similar description of $\Shv_{\omega}\left(\I_{L\subset J}^{\semi} \backslash X\right)$ for the other choices of $X$. 
\end{lem}

%The same arguments as for~\cite[Lemma~6.10]{adr} imply the following statement.

%We now prove that, as in~\cite[Lemma~6.10]{adr}, our perverse t-structures restrict to t-structures on the corresponding subcategories of ``constructible complexes.''

The following statement generalizes~\cite[Lemma~6.10]{adr}.

\begin{lem}
\label{lem:tstr-fg-objects}
The perverse t-structure on $\Shv\left(\I_{L\subset J}^{\semi} \backslash X\right)$ restricts to a t-structure on $\Shv_\omega\left(\I_{L\subset J}^{\semi} \backslash X\right)$. 
\end{lem}

\begin{proof}
In the cases where Lemma~\ref{lem:standard-perverse} applies and for $\Shv \left( \I^\semi_{L \subset J} \backslash \tFl \right)$, the claim follows from the same arguments as for the proof of~\cite[Lemma~6.10]{adr}. The cases that remain are those of the $\infty$-categories $\Shv\left(\I_{L\subset J}^{\semi} \backslash \Fl_K \right)$ when $L$ and $K$ are nonempty. In this case we have to show that for any $\sF$ in $\Shv_\omega\left(\I_{L\subset J}^{\semi} \backslash \Fl_K \right)$ and $n \in \Z$, the complex $\tau_{<n}(\sF)$ belongs to $\Shv_\omega\left(\I_{L\subset J}^{\semi} \backslash \Fl_K \right)$. This follows from the corresponding claim in $\Shv_\omega\left(\I_{J}^{\quar} \backslash \Fl_K \right)$ and the t-exactness of the functor $\oblv^{\I_{L\subset J}^{\semi} \mid \I_J^\quar}$ proved in Lemma~\ref{lem:oblv-semi-quar-texact}.
\end{proof}

%===============================
\subsubsection{Description of the t-structure on finitely generated complexes} 
%===============================

Let $X$ be as above and $J \subset S$.
Thanks to Lem\-ma~\ref{lem:tstr-fg-objects} we have a perverse t-structure on the $\infty$-category $\Shv_\omega\left(\I_{J}^{\quar} \backslash X\right)$.
%, which can be transported via $\Ras^{\I\mid \I_J^{\quar}}$ to a t-structure on $\Shvfg(\I\backslash X)$. 
The following proposition describes this t-structure in terms of the standard perverse t-structure on $\Shv(\I\backslash X)$.

\begin{prop}
\label{prop:tstr-convolution-Wak}
For any $\sF\in \Shvfg(\I\backslash X)$, the following conditions are equivalent:
\begin{enumerate}
\item 
$\Ras^{\I\mid \I_J^{\quar}}(\sF) \in \Shv\left(\I^{\quar}_J\backslash X\right)^{\leq 0}$; 
\item 
for any $\lambda\in \Xd_J^0$,
$\sW_{t_\lambda}\star^\I \sF\in \Shv(\I\backslash X)^{\leq 0}$; 
\item there exists $\nu\in \Xd_J^0 \cap\Xd^{++}_{J^\comp}$ such that $\sW_{t_{n\nu}} \star^\I \sF\in \Shv(\I\backslash X)^{\leq 0}$ for any integer $n >0$.
%$\lambda\in \nu+(\Xd_J^0 \cap\Xd^{+}_{J^\comp})$. 
\end{enumerate}
\end{prop}

\begin{proof}
The proof is similar to that of~\cite[Lemma~6.11]{adr}. 
\end{proof}

\begin{corollary}
\label{cor:Ras-right-texact}
%Let $X=\Fl$, $\tFl$ or the symbol $\sL G/({}^K\Iu,\chi_K)$ (when $(\F,\K)$ is in situation \eqref{item (1)}). 
For any $L\subset J_1\subset J_2\subset S$, the equivalence $\Ras^{\I_{L\subset J_1}^{\semi}\mid \I_{L\subset J_2}^{\semi}}$ is right t-exact. 
\end{corollary}

\begin{proof}
By \eqref{eq:comm-diag-ras-oblv} and Lemma~\ref{lem:oblv-semi-quar-texact}, together with the fact that the functor $\oblv^{\I^\semi_{L\subset J}\mid \I^\quar_J}$ is conservative for any $L\subset J \subset S$, it suffices to prove the claim when $L=\varnothing$. 
We explain the proof in the case when $X=\Fl$; the other cases are similar. 
What we have to show is that $\Ras^{\I_{J_1}^{\quar}\mid \I_{J_2}^{\quar}}(\Delta^{J_1,\quar}_x)$ is contained in $\Shv \left( \I_{J_2}^{\quar}\backslash \Fl \right)^{\leq 0}$ for any $x\in W_\ex$. 
Since $\Delta^{J_1,\quar}_x\in \Shv_{\omega} \left( \I_{J_1}^{\quar}\backslash \Fl \right)$, this property follows from the criterion in Proposition~\ref{prop:tstr-convolution-Wak} since $\Xd_{J_2}^0\subset \Xd_{J_1}^0$.
\end{proof}

Let us also note the following consequence of Proposition~\ref{prop:tstr-convolution-Wak}, which will be used in Section~\ref{sec:t-exactness}. Recall, for $L \subset J \subset S$, the right action of $\Shv(\I \backslash \Fl)$ on $\Shv(\I_{L \subset J}^{\semi} \backslash \Fl)$, see~\S\ref{ss:convolution}.

\begin{corollary}
\label{cor:right-t-exactness-convolution-nabla}
Let $L \subset J \subset S$.
For any $w \in W_\ex$, the functor 
\[
(-) \star^\I \nabla_w : \Shv\left(\I^{\semi}_{L \subset J} \backslash \Fl\right) \to \Shv\left(\I^{\semi}_{L \subset J}\backslash \Fl\right)
\]
is right t-exact for the perverse t-structure.
\end{corollary}

\begin{proof}
Since the functor $\oblv^{\I^{\semi}_{L \subset J} \mid \I^{\quar}_J}$ is conservative and t-exact (see Lemma~\ref{lem:oblv-semi-quar-texact}) and commutes with $(-) \star^\I \nabla_w$, it suffices to treat the case $L=\varnothing$. For this it suffices to show the right t-exactness of the functor
\[
(-) \star^\I \nabla_w : \Shv_\omega\left(\I^{\quar}_J\backslash \Fl\right) \to \Shv_\omega\left(\I^{\quar}_J\backslash \Fl\right).
\]
This property follows from Proposition~\ref{prop:tstr-convolution-Wak}, the commutation of $\Ras^{\I\mid \I_J^{\quar}}$ with $(-) \star^\I \nabla_w$, and the standard fact that $(-) \star^\I \nabla_w : \Shvfg(\I\backslash \Fl) \to \Shvfg(\I\backslash \Fl)$ is right t-exact. (We recall that this fact can be proved after interpreting this functor as the composition of exterior product with a perversely shifted constant sheaf with $*$-pushforward under an affine morphism, which is right t-exact by Artin's theorem, see~\cite[Th\'eor\`eme~4.1.1]{bbdg}. See also~\cite[\S 4.1.3]{ARbook} for details.)
\end{proof}

%---------------------------------------------------------------
\subsection{Perversity of costandard sheaves}
\label{ss:perv-costandards}
%---------------------------------------------------------------

Let $J\subset S$. We will denote by $\unlhd^J$ the preorder on $\Xd$ given by 
\[
\mu \unlhd^J \lambda \quad 
\text{if and only if} \quad 
\lambda-\mu \in \Xd_J^0\cap \Xd_{J^{\comp}}^+. 
\]

Let $x\in W_\ex$, and recall the set $\Xd_{J,x}^{0,+}$ defined in~\S\ref{sss:Wakimoto}.
For $x\in W_\ex$ and $\mu\in \Xd_{J,x}^{0,+}$ we set 
\[
\sW^{J,\vee,\mu}_x:=\nabla_{t_{-\mu}} \star^{\I} \nabla_{t_\mu x} . 
\]
For any $\mu, \lambda\in \Xd_{J,x}^{0,+}$ such that $\mu \unlhd^J \lambda$, there is a canonical morphism 
\[
\sW^{J,\vee,\mu}_x \rightarrow \sW^{J,\vee,\lambda}_x 
\]
defined as follows. Since $\lambda-\mu$ is dominant, using~\cite[Lemma~4.1]{adr} we obtain a canonical isomorphism $\nabla_{t_{-\mu}} \simeq \nabla_{t_{-\lambda}} \star^{\I} \Delta_{t_{\lambda-\mu}}$. On the other hand, since $t_\mu x \in W_\ex^{+,J}$ we have $\ell(t_\lambda x) = \ell(t_{\lambda-\mu}) + \ell(t_\mu x)$, hence a canonical isomorphism $\nabla_{t_\lambda x} \simeq \nabla_{t_{\lambda-\mu}} \star^{\I} \nabla_{t_\mu x}$ (again, by~\cite[Lemma~4.1]{adr}). Then the desired map is the composition
\begin{multline*}
\sW^{J,\vee,\mu}_x 
=\nabla_{t_{-\mu}} \star^{\I} \nabla_{t_\mu x} 
\simeq  \nabla_{t_{-\lambda}} \star^{\I} \Delta_{t_{\lambda-\mu}} \star^{\I} \nabla_{t_\mu x} \\ 
\rightarrow  \nabla_{t_{-\lambda}} \star^{\I} \nabla_{t_{\lambda-\mu}} \star^{\I} \nabla_{t_\mu x}
\simeq \nabla_{t_{-\lambda}} \star^{\I} \nabla_{t_\lambda x} 
=\sW^{J,\vee,\lambda}_x ,
\end{multline*}
where the middle map is induced by the canonical morphism $\Delta_{t_{\lambda-\mu}}\rightarrow \nabla_{t_{\lambda-\mu}}$. 
As in \cite[Lemma 4.5]{adr} one checks that the assignment $\mu \mapsto \sW^{J,\vee,\mu}_x$ induces a functor (in the sense of $\infty$-categories) from $(\Xd_{J,x}^{0,+}, \unlhd^J)$ to $\Shv(\I \backslash \Fl)$, so that we can define
the \emph{dual $J$-Wakimoto object} associated with $x$ as
\[
\sW^{J,\vee}_x:= \colim_{\mu\in \Xd_{J,x}^{0,+}} \sW^{J,\vee,\mu}_x .
\]

If $p>0$ and $\k$ contains a nontrivial $p$-th root of unity, for $A \subset S_\af$ a finitary subset and $x \in W_\ex^{(A,J),\max}$ we also set $\Xd^{0,+}_{J,x,A} = \{ \mu \in \Xd^{0,+}_{J,x} \mid t_\mu x \in W_\ex^{A, \max}\}$, and define
\[
\sW^{J,\vee}_{\chi_A,x} := \colim_{\mu\in \Xd_{J,x,A}^{0,+}} \nabla_{t_{-\mu}} \star^{\I} \nabla_{\chi_A,t_\mu x}.
\]
Using an analogue of~\eqref{eqn:Av-chi-D-classical} for costandard objects, one checks that for all $x \in W_\ex^{(A,J),\max}$ we have
\begin{equation}
\label{eqn:Av-dual-Waki}
\Av^A_\chi(q^! \sW^{J,\vee}_x) \simeq \sW^{J,\vee}_{\chi_A,x}.
\end{equation}

The following is a generalization of \cite[Corollary 6.14]{adr}, which follows from similar arguments. 

\begin{lem}
\phantomsection
\label{lem:cost-dwaki}
\begin{enumerate}
\item
\label{it:cost-dwaki-1}
For any $x\in W_\ex$,
we have an equivalence
\[
\Ras^{\I^\quar_J|\I}(\nabla_x^{J,\quar}) \simeq \sW^{J,\vee}_x. 
\]
\item
\label{it:cost-dwaki-2}
Assume that $p>0$ and $\k$ contains a nontrivial $p$-th root of unity, and let $A \subset S_\af$ be finitary.
For any $x\in W_\ex^{(A,J),\max}$,
we have an equivalence
\[
\Ras^{\I^\quar_J|\I}(\nabla_{\chi_A,x}^{J,\quar}) \simeq \sW^{J,\vee}_{\chi_A,x}. 
\]
\end{enumerate}
\end{lem}

In the setting of Lemma~\ref{lem:cost-dwaki}\eqref{it:cost-dwaki-1}, using Lemma~\ref{lem:Av-Ras} and~\eqref{eqn:Av-dual-Waki}, Lemma~\ref{lem:cost-dwaki} implies that for any $x\in W_\ex^{(A,J),\max}$ we have
\begin{equation}
\label{eqn:Avchi-nabla-quar}
\Av^A_\chi(q^! \nabla_x^{J,\quar}) \simeq \nabla_{\chi_A,x}^{J,\quar}.
\end{equation}

Let $L\subset J\subset S$. 
By Lemma~\ref{lem:cost-dwaki}\eqref{it:cost-dwaki-1}, Lemma~\ref{lem:Av-Ras}\eqref{it:Av-Ras-1}, the commutative diagram obtained from~\eqref{eq:comm-diag-ras-oblv} by passing to right adjoints, and Lemmas~\ref{lem:Av-quar-semi-D-N}\eqref{it:Av-quar-semi-D-N-1} and \ref{lem:push-pull-pK-D-N-semi}\eqref{it:push-pull-pK-D-N-semi-1}, for any finitary $K\subset S_\af$ and $x\in {}^LW^{(K,J)}_\ex$ we have an equivalence 
\begin{equation}
\label{equ:costd-K-colim}
\Ras^{\I^\semi_{L\subset J}\mid \I_L}(\nabla_{K,x}^{L\subset J,\semi}) 
\simeq 
\colim_{\{\mu\in \Xd_{J,x}^{0,+} \mid t_\mu x \in W_\ex^K\}} \big( 
\nabla^L_{t_{-\mu}} \star^{\I} \nabla_{K,t_\mu x} 
\big).
\end{equation}
%(Here, $t_\mu x$ belongs to $W_\ex^K$ for $\mu$ sufficiently large, so that $\nabla_{K,t_\mu x}$ makes sense.)
%Similarly (using Lemma~\ref{lem:av-oblv-whit}\eqref{it:av-oblv-whit-2}), for any finitary $A\subset S_\af$ and $x\in {}^LW^{(A,J),\max}_\ex$ we have an equivalence 
%\begin{equation}
%\label{equ:costd-A-colim}
%\Ras^{\I^\semi_{L\subset J}\mid \I_L}(\nabla_{\chi_A,x}^{L\subset J,\semi}) 
%\simeq 
%\colim_{\mu\in \Xd_{J,x}^{0,+}} \big( 
%\nabla^L_{t_{-\mu}} \star^{\I} \nabla_{\chi_A,t_\mu x} 
%\big). 
%\end{equation}
%(Here again, $\nabla_{\chi_A,t_\mu x}$ makes sense for $\mu$ large enough.)

The following statement generalizes \cite[Proposition~6.12]{adr}. 

\begin{lem}
\label{lem:costandard-perverse}
Let $L\subset J\subset S$ and let $A\subset S_\af$ be finitary. 
The costandard objects $\nabla_x^{L\subset J,\semi}$ for $x \in {}^L W_\ex$ and $\nabla_{\chi_A,x}^{L\subset J,\semi}$ for $x \in {}^L W_\ex^{(A,J),\max}$ (when $p>0$ and $\k$ contains a nontrivial $p$-th root of unity) are contained in the heart of the perverse t-structure on the corresponding $\infty$-categories. 
\end{lem}

\begin{proof}
We show the assertion for the objects $\nabla_x^{L\subset J,\semi}$; the other case follows using Lemma~\ref{lem:pullback-Av-texact}\eqref{it:pullback-q-texact}--\eqref{it:Av-texact}, Lemma~\ref{lem:av-oblv-whit}\eqref{it:av-oblv-whit-1} and~\eqref{eqn:Avchi-nabla-quar}.
%, or can be deduced using Lemmas~\ref{lem:Av-chi-D}\eqref{it:Av-chi-D-1} and~\ref{lem:pullback-Av-texact}\eqref{it:Av-texact}. 
It is clear from definitions that our object belongs to $\Shv\left(\I_{L\subset J}^\semi \backslash \Fl \right)^{\geq 0}$, so what we need to show is that it also belongs to $\Shv\left(\I_{L\subset J}^\semi \backslash \Fl \right)^{\leq 0}$. 
Since our t-structure is compatible with colimits, by \eqref{equ:costd-K-colim} it is enough to show that for any $\mu\in \Xd_{J,x}^{0,+}$ we have
\[
\Ras^{\I_L\mid \I^\semi_{L\subset J}}(\nabla^L_{t_{-\mu}} \star^{\I} \nabla_{t_\mu x}) 
 \in
\Shv\left(\I_{L\subset J}^\semi \backslash \Fl \right)^{\leq 0}. 
\]
For this, by Lemma~\ref{lem:oblv-semi-quar-texact} and Proposition~\ref{prop:tstr-convolution-Wak}, it suffices to show that for any $\lambda, \mu \in \Xd_{J,x}^{0,+}$ such that $\mu \unlhd^J \lambda$ we have
\begin{equation}
\label{equ:costd-perv}
\Delta_{t_\lambda} \star^\I \oblv^{\I_L\mid \I}(\nabla^L_{t_{-\mu}}) \star^{\I} \nabla_{t_\mu x}
 \in
\Shv(\I\backslash \Fl)^{\leq 0}. 
\end{equation}
Now we have equivalences 
\begin{multline*}
\Delta_{t_\lambda} \star^\I \oblv^{\I_L\mid \I}(\nabla^L_{t_{-\mu}})
\simeq \Delta_{L,t_\lambda} \star^{\I_L} \nabla^L_{t_{-\mu}} 
\simeq \oblv^{\I_L\mid \I}(\Delta^L_{L,t_\lambda} \star^{\I_L} \nabla^L_{t_{-\mu}}) \\ 
\simeq \oblv^{\I_L\mid \I}(\Delta^L_{L,t_\lambda} \star^{\I_L} \nabla^L_{L,t_{-\mu}} \star^{\I_L} \nabla^L_e) 
\simeq \oblv^{\I_L\mid \I}(\Delta^L_{L,t_{\lambda-\mu}} \star^{\I_L} \nabla^L_e) \\
\simeq \oblv^{\I_L\mid \I}(\Delta^L_{t_{\lambda-\mu}}),
\end{multline*}
where the second equivalence follows from the fact that $\I z^\lambda \I_L/\I_L= \I_L z^\lambda \I_L/\I_L$, the third and the last one follow from the isomorphism induced by multiplication 
\[
\I_L z^\theta \I_L \times^{\I_L} \I_L/\I \xs \I_L z^\theta \I_L/\I
\]
for all $\theta\in \Xd_J^0$,
and the fourth equivalence follows from \cite[Lemma 3.1.5]{DL23} (which also holds in our setting, with the same proof).
Therefore, the left-hand side of~\eqref{equ:costd-perv} identifies with $\oblv^{\I_L\mid \I}(\Delta^L_{t_{\lambda-\mu}}) \star^\I \nabla_{t_\mu x}$.
Now $\oblv^{\I_L\mid \I}(\Delta^L_{t_{\lambda-\mu}})$ is perverse, so that by Corollary~\ref{cor:right-t-exactness-convolution-nabla} the object $\oblv^{\I_L\mid \I}(\Delta^L_{t_{\lambda-\mu}}) \star^\I \nabla_{t_\mu x}$ belongs to $\Shv(\I\backslash \Fl)^{\leq 0}$, which finishes the proof.
% which by a standard argument (see e.g.~the proof of~\cite[Lemma~4.1.7]{ARbook}) is contained in $\Shv(\I\backslash \Fl)^{\leq 0}$. 
\end{proof}

\begin{rmk}
We have not included in Lemma~\ref{lem:costandard-perverse} the case of the parabolic affine flag varieties $\Fl_K$ with $K \neq \varnothing$. In fact, when $J=S$, it is a standard fact that the objects $\nabla^{L}_{K,x}$ are not always perverse if $L$ and $K$ are both nonempty. 
It is also observed in~\cite[\S 6.6]{adr} that the objects $\nabla_{S,x}^{\semi}$ are never perverse if the characteristic of $\k$ is good for $G$. (More specifically, this reference treats the case $x=e$. The general case follows using~\cite[Eqn.~(5.6)]{adr}.) We do not know what happens in other cases.
%We expect a similar behavior for general $J$.
\end{rmk}

%\begin{rmk}
%It is also an interesting question to determine when the costandard objects $\nabla^{L \subset J,\semi}_{K,x}$ and $\nabla^{L \subset J,\semi}_{\chi_A,x}$ belong to the heart of the perverse t-structure. This case is more delicate than that of standard objects, in particular because costandard objects are not finitely generated (in the sense discussed in~\S\ref{sss:restriction-fg}) in general. When $J=S$, it is standard that the objects $\nabla^{L}_{x}$, $\nabla_{K,x}$ and $\nabla^{L \subset J,\semi}_{\chi_A,x}$ are perverse, but that this is not the case in general for $\nabla^{L}_{K,x}$ if $L$ and $K$ are both nonempty. The case $L=J=\varnothing$ and $K \in \{\varnothing, S\}$ is studied in~\cite{adr}, where it is shown that the objects $\nabla_x^{\semi}$ are always perverse (see~\cite[Proposition~6.12]{adr}) but that the objects $\nabla_{S,x}^{\semi}$ are never perverse if the characteristic of $\k$ is good for $G$. Using Lemma~\ref{lem:Av-chi-D}\eqref{it:Av-chi-D-1} and Lemma~\ref{lem:pullback-Av-texact}\eqref{it:Av-texact}, one can deduce that the objects $\nabla^\semi_{\chi_A,x}$ are also always perverse.
%
%Given these examples, it seems reasonable to expect that the objects $\nabla^{L \subset J, \semi}_x$ are always perverse (which would imply a similar property for the objects $\nabla^{L \subset J, \semi}_{\chi_A,x}$), but we will not study this question in this paper.
%\end{rmk}

%---------------------------------------------------------------
\subsection{t-exactness of convolution actions}
%---------------------------------------------------------------

%===============================
\subsubsection{Right convolution with central sheaves}
%===============================

Consider the abelian category
\[
\Shvfg(\sL^+G \backslash \Gr)^\heartsuit,
\]
i.e.~the heart of the t-structure of middle perversity on $\Shvfg(\sL^+G \backslash \Gr)$. This subcategory is stable under the convolution product $\star^{\sL^+ G}$, and the geometric Satake equivalence~\cite{mv} provides an equivalence of monoidal abelian categories
\begin{equation}
\label{eqn:geom-Satake}
(\Shvfg(\sL^+G \backslash \Gr)^\heartsuit, \star^{\sL^+ G}) \simeq (\Rep^{\heartsuit,\fd}(G^\vee_\k), \otimes) 
\end{equation}
where $G^\vee_\k$ is the split reductive group scheme over $\k$ which is Langlands dual to $G$ and $\Rep^{\heartsuit,\fd}(G^\vee_\k)$ is its abelian category of finite-dimensional representations.
Recall that for any finitary subset $K \subset S_\af$ we have a ``central functor''
\[
\ZCen_K : \Shvfg(\sL^+G \backslash \Gr)^\heartsuit \to \Shvfg(\I_K \backslash \Fl_K)^\heartsuit
\]
defined by nearby cycles along a suitable degeneration of $\Gr$ into $\Fl_K$. The case $K=\varnothing$ is studied in detail in~\cite{ARbook}, and the general case is similar. When $K=S$ we have $\ZCen_S=\id$, and for any $K$ we have $(p_K)_* \circ \ZCen_\varnothing \simeq \oblv^{\I_K | \I} \circ \ZCen_K$.

This functor is ``monoidal'' in the sense that for $\sF,\sG$ in $\Shvfg(\sL^+G \backslash \Gr)^\heartsuit$ we have a canonical identification
\[
\ZCen_K(\sF \star^{\sL^+G} \sG) \simeq \ZCen_K(\sF) \star^{\I_K} \ZCen_K(\sG).
\]
(See~\cite[Theorem~3.4.1]{ARbook} for the case $K=\varnothing$; the general case is similar. Here we write monoidal in quotation marks because $\Shvfg(\I_K \backslash \Fl_K)^\heartsuit$ has no natural monoidal structure.)
Recall also that the central sheaves are ``convolution exact'' in the sense that for any $\sF$ in $\Shvfg(\sL^+G \backslash \Gr)^\heartsuit$ the functor
\[
(-) \star^{\I_K} \ZCen_K(\sF) : \Shv(\Fl_K) \to \Shv(\Fl_K)
\]
of right convolution with $\ZCen_K(\sF)$
is t-exact with respect to the perverse t-structure. (See~\cite[Corollary~3.2.5]{ARbook} for the case $K=\varnothing$; the general case is similar.) 
We will now establish a similar property for semiinfinite sheaves, which generalizes~\cite[Proposition~3.3.15]{DL23}. (That reference treats the case when $L=J$, $K=S$ and $\k$ has characteristic $0$. Note that, even specializing to this case, our proof is different.)

\begin{prop}
\label{prop:right-conv-t-exact}
Let $L \subset J \subset S$, and let $K \subset S_\af$ be finitary. For any $\sF$ in $\Shvfg(\sL^+G \backslash \Gr)^\heartsuit$ the functor
\[
(-) \star^{\I_K} \ZCen_K(\sF) : \Shv \left( \I_{L \subset J}^\semi \backslash \Fl_K \right) \to \Shv \left( \I_{L \subset J}^\semi \backslash \Fl_K \right)
\]
is t-exact with respect to the perverse t-structure.
\end{prop}

\begin{proof}
Using Lemma~\ref{lem:oblv-semi-quar-texact} one reduces the proof to the case $L=\varnothing$.
In this case, first we show right t-exactness. For this we need to show that for any $x \in W_\ex^{(K,J)}$ we have $\Delta^{J,\quar}_{K,x} \star^{\I_K} \ZCen_K(\sF) \in \Shv \left( \I_{J}^\quar \backslash \Fl_K \right)^{\leq 0}$. Since right convolution commutes with Raskin's equivalence, this object is finitely generated, and the desired property is clear from the characterization of the t-structure in Proposition~\ref{prop:tstr-convolution-Wak} and convolution exactness of central sheaves as recalled above.

Left t-exactness follows, since the functor under consideration has left adjoint given by $(-) \star^{\I_K} \ZCen_K(\sF^\vee)$ where $\sF^\vee$ is the dual to $\sF$ in the rigid monoidal category $(\Shvfg(\sL^+G \backslash \Gr)^\heartsuit, \star^{\sL^+ G})$, which is right t-exact by the case treated above.
\end{proof}

%===============================
\subsubsection{Left convolution with central sheaves}
%===============================

Now we consider convolution on the left. Recall from~\S\ref{ss:convolution} that we have an action of $\Shv \left( \I_{L \subset J} \backslash \sL M_J / \I_{L \subset J} \right)$ on $\Shv \left( \I_{L \subset J}^\semi \backslash \Fl_K \right)$ by convolution on the left. We have the ``central functor'' associated with the reductive group $M_J$ and the parahoric subgroup defined by $L$, which provides a ``monoidal'' functor
\begin{equation}
\label{eqn:ZCen-MJ}
\ZCen_{L \subset J} : \Shvfg(\sL^+ M_J \backslash \sL M_J / \sL^+ M_J)^\heartsuit \to \Shvfg \left( \I_{L \subset J} \backslash \sL M_J / \I_{L \subset J} \right)^\heartsuit.
\end{equation}
Recall also that the connected components of the affine Grassmannian $\sL M_J / \sL^+ M_J$ are naturally parametrized by $\Xd / \Z \check{\Phi}_J$ where $\check{\Phi}_J$ is the coroot system of $(M_J,T)$, and that the pairing $\langle -,2\rho-2\rho_J \rangle : \Xd \to \Z$ factors through $\Xd / \Z \check{\Phi}_J$.

The following statement generalizes~\cite[Proposition~3.3.20(ii)]{DL23}.

\begin{prop}
\label{prop:exactness-left-conv}
If $\sF \in \Shvfg(\sL^+ M_J \backslash \sL M_J / \sL^+ M_J)^\heartsuit$ is supported on the component corresponding to $\zeta \in \Xd / \Z \check{\Phi}_J$, then the functor
\[
\ZCen_{L \subset J}(\sF) \star^{\I_{L \subset J}} (-)[-\langle \zeta, 2\rho-2\rho_J \rangle] : \Shv \left( \I_{L \subset J}^\semi \backslash \Fl_K \right) \to \Shv \left( \I_{L \subset J}^\semi \backslash \Fl_K \right)
\]
is t-exact.
\end{prop}

\begin{proof}
The proof is an easy adaptation of that of~\cite[Proposition~3.3.20(ii)]{DL23}, and boils down in the end to convolution exactness of central sheaves. Namely, using Lemma~\ref{lem:pullback-Av-texact}\eqref{it:pullback-p-texact} it suffices to treat the case $K=\varnothing$, and using Lemma~\ref{lem:oblv-semi-quar-texact} one can also assume that $L=\varnothing$. By the same argument as in the proof of Proposition~\ref{prop:right-conv-t-exact} it suffices to prove right t-exactness of our functor, which amounts to proving that for all $w \in W_\ex$ we have
\begin{equation}
\label{eqn:proof-exactness-left-conv}
\ZCen_{\varnothing \subset J}(\sF) \star^{\I_{\varnothing \subset J}} \Delta^{J, \quar}_w [-\langle \zeta, 2\rho-2\rho_J \rangle] \in \Shv \left( \I_{J}^\quar \backslash \Fl \right)^{\leq 0}.
\end{equation}

Fix $\eta \in \Xd_J^0 \cap \Xd_{J^\comp}^{++}$, and consider the action of $\Gm$ on $\Fl$ defined by $\eta$. Then the fixed points $\Fl^0$ for this action identify with a disjoint union of copies of the affine flag variety $\sL M_J / \I_{\varnothing \subset J}$ parametrized by the minimal representatives for the cosets in $W_J \backslash W$. (Here the copy parametrized by $x$ is the $\sL M_J$-orbit of $\dot{x} \I/\I$, and will be denoted $\Fl^0_x$.) Correspondingly, the attractor $\Fl^+$ for this action (in the sense discussed e.g.~in~\cite[\S 1.1.2]{ARbook}) admits a decomposition $\Fl^+ = \sqcup_x \Fl^+_x$ where each component is $\I_{J}^\quar$-stable. For any $x$ we have a ``limit at $0$'' morphism $a_x : \Fl^+_x \to \Fl^0_x$ and a canonical section $b_x : \Fl^0_x \to \Fl^+_x$, and the natural morphism $c_x : \Fl^+_x \to \Fl$ is a locally closed immersion. As in~\cite[Lemma~3.3.2]{DL23}, the $!$-pullback functor associated with $a_x$ induces an equivalence
\[
\Shv(\I_{\varnothing \subset J} \backslash \Fl_x^0)  \simto  \Shv \left( \I_{J}^\quar \backslash \Fl^+_x \right)
\]
with inverse given by $(a_x)_!$ and also by $(b_x)^!$.

Let now $w$ be as above, and write $w=yx$ with $y$ in the extended affine Weyl group $W_{J,\ex}$ of $(M_J,T)$ and $x \in W$ minimal in $W_J x$. Write also $y=t_\lambda z$ with $\lambda \in \Xd$ and $z \in W_J$. 
Then using Remark~\ref{rmk:lJ}\eqref{it:lJ-WJext} we see that $\Delta^{J, \quar}_w$ is the image under the functor
\[
(c_x)_! (a_x)^![- \langle \lambda, 2\rho-2\rho_J \rangle - \ell(x)]
\]
of the standard perverse sheaf in $\Shv(\I_{\varnothing \subset J} \backslash \Fl_x^0)$ associated with $y$ (where $\Fl^0_x$ is identified with the affine flag variety of $M_J$ as above). Hence the object in~\eqref{eqn:proof-exactness-left-conv} is the image under $(c_x)_! (a_x)^![-\langle \lambda, 2\rho-2\rho_J \rangle - \ell(x) -\langle \zeta, 2\rho-2\rho_J \rangle]$ of the convolution of $\ZCen_{\varnothing \subset J}(\sF)$ with that standard perverse sheaf. By convolution exactness of central sheaves (now for $M_J$) the latter convolution is perverse. Taking also into account its support, it is therefore a finite extension of nonnegative cohomological shifts of standard perverse sheaves associated with elements of the form $t_\mu u$ with $u \in W_J$ and $\mu \in \lambda+ \zeta \subset \Xd$. Here we have $\langle \lambda, 2\rho-2\rho_J \rangle + \langle \zeta, 2\rho-2\rho_J \rangle = \langle \mu, 2\rho-2\rho_J \rangle$, so these considerations show that the object in~\eqref{eqn:proof-exactness-left-conv} is a finite extension of nonnegative cohomological shifts of standard objects, hence belongs to $\Shv \left( \I_{J}^\quar \backslash \Fl \right)^{\leq 0}$ as desired.
\end{proof}

%%%%%%%%%%%%%%%%%%%%%%%%%
\section{A t-exactness theorem} 
\label{sec:t-exactness}
%%%%%%%%%%%%%%%%%%%%%%%%%

%---------------------------------------------------------------
\subsection{Statement} 
%---------------------------------------------------------------

Let $X$ be either $\Fl_K$ for some finitary $K \subset S_\af$, $\tFl$, or the symbol $\sL G/(\Iu^A,\chi_A)$ for some finitary $A \subset S_\af$ (assuming that $p>0$ and that $\k$ contains a nontrivial $p$-th root of unity). Our goal in this section is to prove the following statement.

\begin{thm}
\label{thm:texactness}
The composition of functors 
\[
\mathsf{F}:
\Shv\left(\I_J^{\semi}\backslash X\right) \xrightarrow{\oblv^{\I_J^{\semi}|\I_J^{\quar}}} \Shv\left(\I_J^{\quar}\backslash X\right) 
\xrightarrow{\mathrm{Ras}^{\I_J^{\quar}\mid \I^\semi}} \Shv\left(\I^{\semi}\backslash X\right)
\] 
is t-exact with respect to the perverse t-structure on both sides. 
\end{thm}

Standard arguments allow to reduce the proof to the case where for any $\alpha\in \Phis$ there exists a corresponding fundamental coweight 
%$(\check{\varpi}_{\alpha})_{\alpha\in \Phis}$, i.e. 
$\check{\varpi}_{\alpha}\in \Xd$, i.e.~an element which satisfies $\langle \check{\varpi}_{\alpha}, \beta \rangle =\delta_{\alpha,\beta}$ for any $\beta\in \Phis$. We will therefore assume that this condition is satisfied, and fix choices for these fundamental coweights. 
By other standard arguments (see e.g.~\cite[\S 6]{bbdg}), one can also assume that $p>0$ and $\k$ contains a nontrivial $p$-th root of unity, which we will do.

%---------------------------------------------------------------
\subsection{Reduction to the Iwahori--Whittaker case} 
\label{subsect:redtn-IW}
%---------------------------------------------------------------

We write $\mathsf{F}$ as $\mathsf{F}_{D}$ if $X=\sL{G}/D$, where $D=\I_K$, $\Iu$ or the symbol $(\Iu^A,\chi_A)$. 
If $K=S$, then $\Iu^-:=\Iu^S=\Ad_{\dot{w}_0}(\Iu)$ is the pro-unipotent radical of the Iwahori subgroup defined by the Borel subgroup of $G$ which is opposite to $B$ with respect to $T$, and $(\Iu^S,\chi_S)$ is known as the Iwahori--Whittaker condition. 
We will abbreviate $\IW=(\Iu^S,\chi_S)$ and $\Av_{\IW}:=\Av^S_{\chi} \circ q^!$. 

We consider the affine braid group $\mathfrak{B}_\af$ associated with $W_\af$, i.e.~the group generated by symbols $\widetilde{w}$ for $w \in W_\af$, with relations $\widetilde{w} \cdot \widetilde{y} = \widetilde{wy}$ if $\ell(wy)=\ell(w)+\ell(y)$.
We also let $\mathfrak{B}^+_\af \subset \mathfrak{B}_\af$ be the monoid of positive elements, i.e.~the submonoid generated by the elements $\widetilde{w}$ for $w \in W_\af$.
%For $\alpha \in \Phis$, we denote by $\widetilde{s}_\alpha\in \mathfrak{B}_\af$ the canonical lift of $s_\alpha\in S_\af$. 
We will denote by $\ell : \mathfrak{B}^+_\af \to \Z_{\geq 0}$ the unique monoid map such that $\ell(\widetilde{w})=\ell(w)$ for any $w \in W_\af$, and will call $\ell(b)$ the length of $b$.

It is a standard fact (which follows e.g.~from~\cite[Lemma~4.1]{adr}) that the assignment $\widetilde{w} \mapsto \nabla_w$ extends to a group morphism from $\mathfrak{B}_\af$ to the group of isomorphism classes of invertible objects in $\Shv(\I\backslash \sL{G}/\I)$. We will denote by $\nabla_b$ the image of $b \in \mathfrak{B}_\af$. In particular, for $w \in W_\af$ we have $\nabla_{\widetilde{w}^{-1}} = \Delta_{w^{-1}}$. Similarly, the assignment $\widetilde{w} \mapsto \Delta_w$ extends to a group morphism from $\mathfrak{B}_\af$ to the group of isomorphism classes of invertible objects in $\Shv(\I\backslash \sL{G}/\I)$, and we will denote by $\Delta_b$ the image of $b \in \mathfrak{B}_\af$.
%Consider the assignment 
%\[
%b\in \mathfrak{B}^+_\af \quad \mapsto \quad \nabla_b\in \Shv(\I\backslash \sL{G}/\I),
%\] 
%where $\nabla_b:=\nabla_{\widetilde{s}_{\alpha_1}}\star^\I \cdots \star^\I \nabla_{\widetilde{s}_{\alpha_n}}$ if $b=\widetilde{s}_{\alpha_1}\cdots \widetilde{s}_{\alpha_n}$ in $\mathfrak{B}^+_\af$. (It is a standard fact that this object does not depend on the choice of decomposition of $b$.)

The following statement is an extension to general coefficients of~\cite[Proposition~3]{BLin}, whose proof (which is similar to that in~\cite{BLin}) will be explained in~\S\ref{ss:pf-lem}.

\begin{lem}
\label{lem:bez-lin}
Let $\sG\in \Shvfg(\I\backslash \Fl)$. 
We have $\sG\in \Shv(\I\backslash \Fl)^{\leq 0}$ if and only if
\[
\Av_\IW(\sG \star^\I \nabla_b)\in \Shv(\I\backslash \sL{G}/\IW)^{\leq 0}
\]
for any $b\in \mathfrak{B}^+_\af$. 
\end{lem}

We deduce the following corollary.

\begin{corollary}
\label{cor:tstr-criterion-IW}
Let $\sG\in \Shv_{\omega}\left(\I^\semi\backslash \Fl\right)$. 
We have $\sG\in \Shv\left(\I^\semi\backslash \Fl\right)^{\leq 0}$ if and only if $\Av_\IW(\sG \star^\I \nabla_b)\in \Shv\left(\I^\semi\backslash \sL{G}/\IW\right)^{\leq 0}$ for any $b\in \mathfrak{B}^+_\af$. 
\end{corollary}

\begin{proof}
Since the equivalence $\Ras^{\I^\semi \mid \I}$ commutes with $\Av_\IW$ and right convolution with any $\nabla_b$, the assertion follows from Proposition~\ref{prop:tstr-convolution-Wak} and Lemma~\ref{lem:bez-lin}. 
\end{proof}

The following lemma will be used to reduce the proof of Theorem~\ref{thm:texactness} to the case $X=\sL{G}/\IW$. 

\begin{lem}
\phantomsection
\label{lem:reduction-texactness}
\begin{enumerate}
\item 
\label{it:reduction-texactness-1}
If $\mathsf{F}_\I$ is right t-exact, then so are $\mathsf{F}_{\Iu}$, $\mathsf{F}_{\I_K}$ and $\mathsf{F}_{(\Iu^A,\chi_A)}$ for any finitary $K,A \subset S_\af$. 
\item 
\label{it:reduction-texactness-2}
If $\mathsf{F}_{\IW}$ is right t-exact, then so is $\mathsf{F}_\I$. 
\end{enumerate}
\end{lem}

\begin{proof}
%By Corollary~\ref{cor:Ras-right-texact} the functor $\mathsf{F}$ is left t-exact. 
\eqref{it:reduction-texactness-1}
To show that $\mathsf{F}_{\Iu}$, resp. $\mathsf{F}_{(\Iu^A,\chi_A)}$, is right t-exact, it suffices to show that $\mathsf{F}_{\Iu}(q^!(\Delta_x^{J,\semi}))$, resp. $\mathsf{F}_{(\Iu^A,\chi_A)}(\Delta^{J,\semi}_{\chi_A,x})$, is contained in the ``$\leq 0$'' part of the perverse t-structure for any $x\in {}^JW_\ex$, resp. $x\in {}^JW_\ex^{(A,J),\max}$. 
This follows from Lemma~\ref{lem:Av-Ras}, Lemma~\ref{lem:Av-chi-D}\eqref{it:Av-chi-D-1}, Lemma~\ref{lem:pullback-Av-texact}\eqref{it:pullback-q-texact}--\eqref{it:Av-texact}, and the assumption that $\mathsf{F}_{\I}(\Delta_x^{J,\semi})$ is contained in $\Shv(\I^\semi\backslash \Fl)^{\leq 0}$. 

The assertion for $\mathsf{F}_{\I_K}$ follows from Lemma~\ref{lem:Av-Ras}\eqref{it:Av-Ras-1}, Lemma~\ref{lem:pullback-Av-texact}\eqref{it:pullback-p-texact} and the fact that the functor $(p_K)^\dag$ is conservative. 

\eqref{it:reduction-texactness-2}
Let us assume that $\mathsf{F}_{\IW}$ is known to be right t-exact. To show that $\mathsf{F}_{\I}$ is right t-exact,
it suffices to show that for any $\sF\in \Shv_{\omega}\left(\I_J^\semi\backslash \Fl\right)^{\leq 0}$ (in particular, when $\sF=\Delta_x^{J,\semi}$ for some $x \in {}^J W_\ex$) we have $\mathsf{F}_{\I}(\sF)\in \Shv \left(\I^\semi\backslash \Fl\right)^{\leq 0}$. 
By Corollary~\ref{cor:tstr-criterion-IW}, it is equivalent to show that 
\[
\mathsf{F}_{\IW}(\Av_\IW(\sF\star^\I \nabla_b)) \simeq \Av_\IW(\mathsf{F}_{\I}(\sF) \star^\I \nabla_b) \in \Shv\left(\I^\semi\backslash \sL{G}/\IW\right)^{\leq 0} 
\]
for all $b\in \mathfrak{B}^+_\af$.
Since by our assumption $\mathsf{F}_{\IW}$ is right t-exact, and since by Lem\-ma~\ref{lem:pullback-Av-texact}\eqref{it:pullback-q-texact}--\eqref{it:Av-texact} the functor $\Av_\IW$ is t-exact, it is enough to show that the functor 
\[
(-) \star^\I \nabla_b : \Shv\left(\I^{\semi}_J\backslash \Fl\right) \to \Shv\left(\I^{\semi}_J\backslash \Fl\right)
\]
is right t-exact for any $b \in \mathfrak{B}_\af^+$, which follows from Corollary~\ref{cor:right-t-exactness-convolution-nabla}.
%Since the functor $\oblv^{\I^{\semi}_J \mid \I^{\quar}_J}$ is conservative and t-exact, it suffices to show the right t-exactness of the functor
%\[
%(-) \star^\I \nabla_b : \Shv\left(\I^{\quar}_J\backslash \Fl\right) \to \Shv\left(\I^{\quar}_J\backslash \Fl\right)
%\]
%for any $b \in \mathfrak{B}_\af^+$,
%or equivalently of the restrictions of these functors to $\Shv_{\omega}\left(\I^{\quar}_J\backslash \Fl\right)$. 
%This property follows from Proposition~\ref{prop:tstr-convolution-Wak} and the standard fact (which follows e.g.~from~\cite[Lemma~4.1(3)]{adr}) that $(-) \star^\I \nabla_b : \Shvfg(\I\backslash \Fl) \to \Shvfg(\I\backslash \Fl)$ is right t-exact for any $b \in \mathfrak{B}_\af^+$. 
\end{proof}

%SR: END OF PROOFREADING

%---------------------------------------------------------------
\subsection{Proof of Theorem~\ref{thm:texactness}} 
%---------------------------------------------------------------

For $\lambda\in \Xd_J^{++}$, we abbreviate 
\[
\Delta^{J,\semi}_{\IW,\lambda}:=\Delta^{J,\semi}_{\chi_S,w_Jt_{\lambda}w_0}.
\]
By Lemma~\ref{lem:JSJw},
%and~\eqref{eq:decomp-tGr},
%Lemma~\ref{lem:whit}(2),
these objects are exactly the
%give a complete collection of 
standard perverse sheaves in the $\infty$-category $\Shv \left( \I^\semi_J\backslash \sL{G}/\IW \right)$. By Remark~\ref{rmk:JSJw-length}, the cohomological shift appearing in the definition of this object is $-\langle \lambda, 2\rho \rangle - \ell(w_0)$.

For $\lambda\in \Xd_J^+$, we will denote by 
\[
\sA^{M_J}_{\lambda} \in \Shv(\sL^+M_J\backslash \sL{M_J}/\sL^+M_J)^\heartsuit 
\]
the 0-th perverse cohomology of the standard complex attached to the $(\sL^+M_J\times \sL^+M_J)$-orbit of $z^\lambda$. (These objects are the perverse sheaves corresponding to the Weyl modules for the reductive group Langlands dual to $M_J$ under the geometric Satake equivalence~\eqref{eqn:geom-Satake} for $M_J$; see e.g.~\cite[\S 1.5.1]{ARbook}.)

We set $\zeta_J=\sum_{\alpha\in \Phis_J} \check{\varpi}_{\alpha}$; then $\Xd_J^{++}=\zeta_J+\Xd_J^{+}$. 

\begin{lem}
\label{lem:D-convolution}
For any $\lambda\in \Xd_J^{+}$, we have
\[
\Delta^{J,\semi}_{\IW,\lambda+\zeta_J} \simeq \sA^{M_J}_{\lambda} \star^{\sL^+M_J} \Delta^{J,\semi}_{\IW,\zeta_J} [-\langle \lambda, 2\rho-2\rho_J \rangle]. 
\]
\end{lem}

\begin{proof}
We consider 
%the Iwahori--Whittaker condition 
%$\I_{\varnothing\subset J,\sW}$, defined on 
%with respect to 
the pro-unipotent radical $\I^-_{\varnothing\subset J,\mathrm{u}}$ of the Iwahori subgroup in $\sL{M_J}$ defined by the Borel subgroup of $M_J$ which is opposite to $B_J$ with respect to $T$, and fix a morphism $\I^-_{\varnothing\subset J,\mathrm{u}} \to \Ga$ as in~\S\ref{sss:Whittaker}. Considering the associated Iwahori--Whittaker condition on sheaves on the ind-scheme $\sL^+ M_J \backslash \sL M_J$ we obtain the $\infty$-category $\Shv(\sL^+M_J\backslash \sL{M_J}/\IW_J)$.

On the other hand, considering the pullback under the induced morphism $\I^-_{\varnothing\subset J,\mathrm{u}}\ltimes \sL{U_J} \to \Ga$ of our fixed Artin--Schreier local system on $\Ga$ we obtain a twisted equivariance condition on sheaves on $\sL G$, which we will denote
%is denoted by 
$\I^\quar_{J,\sW}$. 
We can form the ``bi-Whittaker" category 
%as in \textsection\ref{subset shvs} 
\[
\Shv \left( \I^\quar_{J,\sW}\backslash \sL{G}/\IW \right)
\]
following the same recipe as in~\S\ref{sss:Whittaker}; we will
denote by $\Delta^{(J,\sW),\quar}_{\IW,e}$ the standard sheaf
%\footnote{We need to explain what is the cohomological shift here.} 
associated with the $(\I^-_{\varnothing\subset J,\mathrm{u}}\ltimes \sL{U_J})\times \Iu^-$-orbit through $e$, i.e.~the $!$-pushforward under the natural embedding of the dualizing sheaf on this orbit, shifted by $-\ell(w_0)$.
%(where no cohomological shift is involved in the definition). 

Convolution gives a functor
\[
\star^{\IW_J}: \Shv \left( \sL^+M_J\backslash \sL{M_J}/\IW_J \right) \times \Shv \left( \I^\quar_{J,\sW}\backslash \sL{G}/\IW \right) \rightarrow \Shv \left( \I^\semi_{J}\backslash \sL{G}/\IW \right). 
\] 
For $\mu\in \Xd_J^{++}$, let $\sA^{M_J}_{\sW,\mu}\in \Shv(\sL^+M_J\backslash \sL{M_J}/\IW_J)$ be the standard sheaf associated with the orbit $\sL^+M_J\cdot z^\mu\cdot \I^-_{\varnothing\subset J,\mathrm{u}}$, i.e.~the $!$-pushforward under the natural embedding of the dualizing sheaf on this orbit, shifted by $-\langle \mu, 2\rho_J \rangle$. 
For any $\nu\in \Xd$, multiplication induces an isomorphism 
\[
\sL^+M_J\cdot z^{\nu}\cdot \I^-_{\varnothing\subset J,\mathrm{u}} \times^{\I^-_{\varnothing\subset J,\mathrm{u}}} (\I^-_{\varnothing\subset J,\mathrm{u}}\ltimes \sL{U_J})\cdot \Iu^- \xs 
\I^\semi_J\cdot z^{\nu}\cdot \Iu^-.
\]
Hence we have an equivalence, for any $\mu\in \Xd_J^{++}$, 
\[
\sA^{M_J}_{\sW,\mu}\star^{\IW_J} \Delta^{(J,\sW),\quar}_{\IW,e} [-\langle \mu, 2\rho-2\rho_J \rangle]
\simeq \Delta^{J,\semi}_{\IW,\mu}. 
\]
Now the assertion follows from the fact that
\[
\sA_{\lambda}^{M_J}\star^{\sL^+M_J} \sA_{\sW,\zeta_J}^{M_J}\simeq \sA_{\sW,\lambda+\zeta_J}^{M_J},
%\qedhere 
\]
for any $\lambda \in \Xd_J^{+}$,
see~\cite[Proof of Theorem 3.9]{BGMRR}.
\end{proof}

\begin{proof}[Proof of Theorem~\ref{thm:texactness}]
By Lemma~\ref{lem:oblv-semi-quar-texact} the functor $\oblv^{\I_J^{\semi}|\I_J^{\quar}}$ is t-exact, and by Corollary~\ref{cor:Ras-right-texact} the functor $\mathrm{Ras}^{\I_J^{\quar}\mid \I^\semi}$ is left t-exact, as the inverse of a right t-exact equivalence. Hence $\mathsf{F}$ is left t-exact. 
By Lemma~\ref{lem:reduction-texactness}, to conclude we therefore only need to show the right t-exactness of $\mathsf{F}_\IW$. 
For this it suffices to show that for any $\lambda\in \Xd_J^+$ we have
\[
\mathsf{F}_\IW \left( \Delta^{J,\semi}_{\IW,\lambda+\zeta_J} \right) \in \Shv \left( \I^{\semi}\backslash \sL{G}/\IW \right)^{\leq 0},
\]
i.e.~(in view of Proposition~\ref{prop:tstr-convolution-Wak}) that for any $\mu \in \Xd$ we have
\[
\sW_{t_\mu} \star^\I \Ras^{\I^\semi \mid \I} \left( \mathsf{F}_\IW \left( \Delta^{J,\semi}_{\IW,\lambda+\zeta_J} \right) \right)
\in \Shv \left( \I \backslash \sL{G}/\IW \right)^{\leq 0}.
\]

By Lemma~\ref{lem:D-convolution}, we have 
\begin{multline*}
\mathsf{F}_\IW \left( \Delta^{J,\semi}_{\IW,\lambda+\zeta_J} \right) \\
\simeq \mathrm{Ras}^{\I_J^{\quar}\mid \I^\semi}\circ \oblv^{\I_J^\semi | \I_{J}^\quar } \left( \sA^{M_J}_{\lambda} \star^{\sL^+M_J} \Delta^{J,\semi}_{\IW,\zeta_J} \right) [-\langle \lambda, 2\rho-2\rho_J \rangle] \\ 
\simeq \mathrm{Ras}^{\I_J^{\quar}\mid \I^\semi} \left( \ZCen_{\varnothing \subset J}(\sA^{M_J}_{\lambda}) \star^{\I_{\varnothing \subset J}} \oblv^{\I_J^\semi | \I_J^\quar} (\Delta^{J,\semi}_{\IW,\zeta_J}) \right) [-\langle \lambda, 2\rho-2\rho_J \rangle]
\end{multline*}
where $\ZCen_{\varnothing \subset J}$ is as in~\eqref{eqn:ZCen-MJ}.
%\[
%\mathsf{Z}_J:\Shv(\sL^+M_J\backslash \sL{M_J}/\sL^+M_J)\rightarrow \Shv(\I_{\varnothing \subset J}\backslash \sL{M_J}/\I_{\varnothing \subset J})
%\]
%is the central functor constructed as in~\cite{ARbook}.
Now, recall the Wakimoto sheaves for $M_J$ considered in Lemma~\ref{lem:Ras-convolution-Wak}.
Since $\ZCen_{\varnothing \subset J}(\sA^{M_J}_{\lambda})$ admits a filtration by objects $\sW^{M_J}_{t_\mu}$ where $\mu + \Z\check{\Phi}_J = \lambda + \Z\check{\Phi}_J$ (so that $\langle \lambda, 2\rho-2\rho_J \rangle = \langle \mu, 2\rho-2\rho_J \rangle$), see e.g.~\cite[Theorem~4.4.5]{ARbook}, by Lemma~\ref{lem:Ras-convolution-Wak} it suffices to show that 
\begin{equation}
\label{eq:FIW-D}
\mathsf{F}_\IW \left( \Delta^{J,\semi}_{\IW,\zeta_J} \right) \in \Shv \left( \I^{\semi}\backslash \sL{G}/\IW \right)^{\leq 0}. 
\end{equation}

We set $w_{\zeta_J}:=w_Jt_{\zeta_J}w_0$. By 
%the result cited in the proof of 
Lemma~\ref{lem:JSJw} and its proof,
this element has the following properties.

\begin{claim}
\label{cl:properties-w-zetaJ}
\begin{enumerate}
\item 
\label{it:w-zetaJ-1}
$w_{\zeta_J}\in {}^JW_\ex^{S,\max} \cap {}^JW_\ex^{(S,J),\max}$; 
%\item 
%\label{it:w-zetaJ-2}
%$\ell(w_{\zeta_J})=\ell^J(w_{\zeta_J})$; 
\item 
\label{it:w-zetaJ-3}
$w_{\zeta_J}$ is a minimal element in ${}^JW_\ex^{S,\max}$ for the Bruhat order.
\end{enumerate}
\end{claim}

%\begin{proof}
%COMPLETE.
%\end{proof}

By Proposition~\ref{prop:Ras-D}\eqref{it:Ras-D-2} we have 
\[
\mathrm{Ras}^{\I_J^{\quar}\mid \I} \left( \Delta^{J,\quar}_{\chi_S,w_{\zeta_J}} \right)
\simeq \sW^{J}_{\chi_S,w_{\zeta_J}}.
%\simeq \Delta_{\chi_S,w_{\zeta_J}},
\]
%where the second equivalence follows from Claim~\ref{cl:properties-w-zetaJ}\eqref{it:w-zetaJ-2}.\footnote{SR: I don't understand this argument.}
Hence by Claim~\ref{cl:properties-w-zetaJ}\eqref{it:w-zetaJ-1}, Lemma~\ref{lem:av-oblv-whit}\eqref{it:av-oblv-whit-1}, and the diagram obtained from~\eqref{eq:comm-diag-ras-oblv} by passing to left adjoints,
%(and a similar statement for standard sheaves in $\Shv(\I\backslash \sL{G}/\IW)$)
we have 
\begin{multline*}
\mathrm{Ras}^{\I_J^{\semi}\mid \I_J} \left( \Delta^{J,\semi}_{\IW,\zeta_J} \right) \simeq
\mathrm{Ras}^{\I_J^{\semi}\mid \I_J} \circ \Av^{\I^\semi_J \mid \I^\quar_J}_! \left (\Delta^{J,\quar}_{\chi_S,w_{\zeta_J}} \right) [-\ell(w_J)] \\
\simeq \Av^{\I_J \mid \I}_! \circ \mathrm{Ras}^{\I_J^{\quar}\mid \I} \left (\Delta^{J,\quar}_{\chi_S,w_{\zeta_J}} \right) [-\ell(w_J)]
\simeq \Av^{\I_J \mid \I}_! \left( \sW^{J}_{\chi_S,w_{\zeta_J}} \right) [-\ell(w_J)].
%\Delta^J_{\chi_S,w_{\zeta_J}}.
%\simeq \nabla^J_{\chi_S,w_{\zeta_J}}, 
\end{multline*}
This complex is supported on the closure of the orbit parametrized by $w_{\zeta_J}$, and by Lemma~\ref{lem:av-oblv-whit}\eqref{it:av-oblv-whit-1} and the properties of Wakimoto objects explained in~\S\ref{sss:Wakimoto} its restriction to this orbit coincides with that of $\Delta^J_{\chi_S,w_{\zeta_J}}$.
Using Claim~\ref{cl:properties-w-zetaJ}\eqref{it:w-zetaJ-3}, we deduce equivalences
\begin{multline*}
\mathrm{Ras}^{\I_J^{\quar}\mid \I} \circ \oblv^{\I_J^{\semi} \mid \I_J^\quar} \left( \Delta^{J,\semi}_{\IW,\zeta_J} \right) \simeq
\oblv^{\I_J\mid \I} \circ \mathrm{Ras}^{\I_J^{\semi}\mid \I_J} \left( \Delta^{J,\semi}_{\IW,\zeta_J} \right) \\
\simeq \oblv^{\I_J\mid \I} (\Delta^J_{\chi_S,w_{\zeta_J}}) \simeq \oblv^{\I_J\mid \I} (\nabla^J_{\chi_S,w_{\zeta_J}}).
\end{multline*}
Hence, by Proposition~\ref{prop:tstr-convolution-Wak} (applied with ``$J$'' being the empty set), to conclude it suffices to prove that for any $\theta\in \Xd^+$ the complex
\[
\Delta_{t_\theta}\star^{\I}\oblv^{\I_J\mid \I}(\nabla^J_{\chi_S,w_{\zeta_J}})
\simeq ((p_J)_* \Delta_{t_\theta}) \star^{\I_J} \nabla^J_{\chi_S,w_{\zeta_J}}
\simeq \Delta_{J,t_\theta} \star^{\I_J} \nabla^J_{\chi_S,w_{\zeta_J}}
\]
belongs to $\Shv(\I\backslash \sL{G}/\IW)^{\leq 0}$. (Here, the second equivalence follows from the fact that $t_\theta \in W_\ex^J$.)
Now, interpreting the complexes as living on the ind-scheme $\I \backslash \sL G$, one can write the functor
\[
(-) \star^{\I_J} \nabla^J_{\chi_S,w_{\zeta_J}} : \Shv(\I \backslash \sL G / \I_J) \to \Shv(\I \backslash \sL G / \IW)
\]
as exterior product with a perversely shifted local system, followed by $*$-pushforward under an affine morphism, hence this functor is right t-exact by the same considerations as in the proof of Corollary~\ref{cor:right-t-exactness-convolution-nabla}. The desired claim follows, since $\Delta_{J,t_\theta}$ belongs to $\Shv(\I \backslash \sL G / \I_J)^{\leq 0}$. (In fact, this object is well known to be perverse.)
%
%Now, by the same argument as for~\eqref{eqn:computation-standard-perverse}, we have
%\[
%\Delta_{t_\theta}\star^{\I}\oblv^{\I_J\mid \I}(\nabla^J_{\chi_S,w_{\zeta_J}})
%\simeq \Delta^J_{t_\theta}\star^{\I} \nabla_{\chi_S,w_{\zeta_J}},
%\]
%and one concludes using Corollary~\ref{cor:right-t-exactness-convolution-nabla} (for ``$J$'' being $S$) and t-exactness of the functor $\Av_{\IW}$ (see~\cite[Lemma~3.2]{AR23} or Lemma~\ref{lem:pullback-Av-texact}).
%
%Finally, by a standard argument (as in the proof of Lemma~\ref{lem:standard-perverse}) the object 
%\[
%\Delta_{t_\theta}\star^{\I}\oblv^{\I_J\mid \I}(\nabla^J_{\chi_S,w_{\zeta_J}})
%\simeq \Delta_{J,t_\theta} \star^{\I_J} (\nabla^J_{\chi_S,w_{\zeta_J}})
%\]
%is perverse in $\Shv(\I\backslash \sL{G}/\IW)$, for any $\theta\in \Xd_J^0 \cap \Xd^+_{J^\comp}$. 
%Hence by Proposition~\ref{prop:tstr-convolution-Wak} we conclude that~\eqref{eq:FIW-D} holds, as desired. 
\end{proof}

%---------------------------------------------------------------
\subsection{Proof of Lemma~\ref{lem:bez-lin}} 
\label{ss:pf-lem}
%---------------------------------------------------------------

In this subsection, for the reader's convenience we explain the proof of Lemma~\ref{lem:bez-lin} following~\cite{BLin}.

\begin{proof}[Proof of Lemma~\ref{lem:bez-lin}]
The ``only if" part is easy, since the functor $(-) \star^{\I}\nabla_b$ is right t-exact (see Corollary~\ref{cor:right-t-exactness-convolution-nabla}) and $\Av_{\IW}$ is t-exact (see~\cite[Lemma~3.2]{AR23} or Lemma~\ref{lem:pullback-Av-texact}). 
It remains to show the ``if" part. 
Let $n$ be the maximal integer such that $\sG':={}^p\sH^n(\sG)\neq 0$ (which exists since $\sG\in \Shvfg(\I\backslash \Fl)$). 

\textit{Step 1.} 
We claim that  
\begin{equation}
\label{eq:AvIW-convolution-nonvanishing}
\Av_\IW(\sG'\star^{\I}\nabla_b) \neq 0 \quad \text{for some $b\in \mathfrak{B}_\af^+$.}
\end{equation}
The proof will use the following fact (see e.g.~\cite[Proposition~4.4.4]{ARbook}): there exists a finite subset $X \subset W_\ex$ such that for any $x\in W_\ex$ and $y\in W_\af$, 
\[
(\mathbf{s}^{\varnothing \subset S}_{\varnothing,x})^!(\sG'\star^{\I}\nabla_{y}) \neq 0 \quad \Rightarrow \quad x\in X \cdot y.
\]
%where $\widetilde{y}\in \mathfrak{B}^+_\af$ is the canonical lift of $y$. 

Now choose $\lambda\in \Z\check{\Phi}$ 
%be dominant enough 
such that $X \cdot t_\lambda$ consists of elements of the form $wt_\mu$, with $w\in W$ and $\mu\in \Xd^{++}$.
% (in particular, $A\cdot t_\lambda\subset W_\ex^{S,\mathrm{w}}$). 
%Then for each right $W$-coset in $W_\ex$, there is at most one element contained in $A\cdot t_\lambda$. 
If $x\in X \cdot t_\lambda$ satisfies $(\mathbf{s}^{\varnothing \subset S}_{\varnothing,x})^!(\sG'\star^\I \nabla_{t_\lambda})\neq 0$ and $x' \in W_\ex^{S,\max}$ is the unique element such that $x'W = xW$, then we have 
\begin{align*}
\Hom^\bullet(\Delta_{\chi_S,x'}, \Av_\IW(\sG'\star^{\I}\nabla_{{t_\lambda}}))
&\simeq \Hom^\bullet(\Av^S_!(\Delta_{\chi_S,x'}), q^!(\sG'\star^{\I}\nabla_{{t_\lambda}}))\\
&\simeq \Hom^\bullet(q^! \Delta_x, q^!(\sG'\star^{\I}\nabla_{{t_\lambda}}))\neq 0,
\end{align*}
where in the second equivalence we use the fact that $\Av^S_!(\Delta_{\chi_S,x})$ is an iterated extension of objects $q^! \Delta_{xw}$ for $w\in W$, each appearing once (see Lemma~\ref{lem:Av-chi-D}\eqref{it:Av-chi-D-2}), and that $x'W \cap (X \cdot t_\lambda)=\{x\}$. 
We deduce that $\Av_\IW(\sG'\star^{\I}\nabla_{{t_\lambda}}) \neq 0$, and hence \eqref{eq:AvIW-convolution-nonvanishing}. 

\textit{Step 2.} 
We let $b_0$ be an element of minimal length in $\{b\in \mathfrak{B}^+_\af \mid \Av_\IW(\sG'\star^{\I}\nabla_{b})\neq 0\}$. Then we claim that ${}^p\sH^0(\Av_\IW(\sG'\star^{\I}\nabla_{b_0}))\neq 0$. 

In fact, for any simple reflection $s\in S_\af$, we have a
%Consider the 
long exact sequence
\begin{equation*}
%\label{equ 1.5}
0\rightarrow \Delta_e \rightarrow \Delta_s \rightarrow \nabla_s \rightarrow \Delta_e \rightarrow 0
\end{equation*} 
in $\Shv(\I\backslash \Fl)^\heartsuit$.
%Hence using \eqref{equ 1.5} one can show 
Using these exact sequences one sees
that for any $b\in \mathfrak{B}_\af^+$ 
%(where $\Delta_{b^{-1}}$ is the inverse of $\nabla_{b}$) 
we have
\[
\sG'\star^{\I}\nabla_{b} \simeq \sG'\star^{\I}\Delta_{b} \quad \text{mod}\quad \langle \sG'\star\nabla_{b'} : b' \in \mathfrak{B}_\af^+ \text{ with } \ell(b')<\ell(b)\rangle
\]
where in the right-hand side we consider the full stable subcategory generated by the given objects, and the notation means an equivalence in the Verdier quotient.
By our assumption on $b_0$, we therefore have $\Av_{\IW}(\sG'\star^{\I}\nabla_{b_0}) \simeq \Av_{\IW}(\sG'\star^{\I}\Delta_{b_0})$. Since $(-) \star^{\I}\nabla_{b_0}$ is right t-exact, see Corollary~\ref{cor:right-t-exactness-convolution-nabla}, $(-)\star^{\I}\Delta_{b_0}$ is left t-exact by similar considerations, and $\Av_{\IW}$ is t-exact, this object belongs to both the ``$\leq 0$'' and the ``$\geq 0$'' parts of the t-structure, hence is perverse. 
%Since $\Av_{\IW}$ is t-exact, we have 
We deduce that ${}^p\sH^0(\Av_\IW(\sG'\star^{\I}\nabla_{b_0}))\simeq \Av_\IW(\sG'\star^{\I}\nabla_{b_0})\neq 0$, as desired.

\textit{Step 3.} 
We now conclude. Let $b_0$ be as in Step 2.
%We reduce the assertion to the following claim: 
%\begin{equation}
%{}^p\sH^0(\Av_\IW(\sG' \star^{\I}\nabla_b))\neq 0, \quad \text{for some } b\in \mathfrak{B}_\af^+. 
%\end{equation}
By right t-exactness of the $\mathfrak{B}_\af^+$-action and t-exactness of $\Av_\IW$, we have ${}^p\sH^{>n}(\Av_\IW(\sG\star^{\I}\nabla_{b_0}))=0$ and 
\[
{}^p\sH^n(\Av_\IW(\sG\star^{\I}\nabla_{b_0}))\simeq {}^p\sH^0(\Av_\IW(\sG'\star^{\I}\nabla_{b_0})) \neq 0. 
\]
Hence the fact that $\Av_\IW(\sG \star^{\I}\nabla_{b_0})\in \Shv(\I\backslash \sL{G}/\IW)^{\leq 0}$ forces $n\leq 0$, or in other words $\sG\in \Shv(\I\backslash \Fl)^{\leq 0}$. 
\end{proof}

%%%%%%%%%%%%%%%%%%%%%%%%%
\section{Gaitsgory sheaves} 
\label{sec:Gaitsgory}
%%%%%%%%%%%%%%%%%%%%%%%%%

%---------------------------------------------------------------
\subsection{Construction}
%---------------------------------------------------------------

In this subsection we explain the construction of parabolic counterparts of the ``Gaitsgory sheaf'' studied in~\cite{Gaits18} (when $\k$ has characteristic $0$) and~\cite{adr}. A similar construction is also considered (when $\k$ has characteristic $0$) in~\cite{DL23}. (The authors of~\cite{DL23} consider a more general construction; ours corresponds to the case $\mu=0$ in their notation.)

We fix $J \subset S$, and work in the $\infty$-category $\Shv\left( \I^\quar_J \backslash \Gr \right)$. According to our general constructions the $\I^\quar_J$-orbits in $\Gr$ are naturally parametrized by $W_\ex^{(S,J)}$, but this set is in bijection with $W_\ex / W$ (by Lemma~\ref{lem:bi-min}\eqref{it:bi-min-1}), hence also with $\Xd$. For $\nu \in \Xd$ we will therefore set
\[
\mathrm{S}_\nu^J = \I^\quar_J z^\nu \sL^+G/\sL^+G,
\]
and denote by $\sigma^J_\nu$ the embedding of $\mathrm{S}_\nu^J$ in $\Gr$. (So, we have $\mathrm{S}^J_\nu = \fS^{\varnothing \subset J}_{S,x^J_\nu}$ where $x^J_\nu$ is the unique element in $t_\nu W \cap W_{\ex}^{(S,J)}$. When $J=\varnothing$, $\mathrm{S}_\nu^\varnothing$ is denoted $\mathrm{S}_\nu$ in~\cite{adr}, and when $J=S$ we are considering the $\I$-orbits on $\Gr$, denoted $(\Gr_\nu : \nu \in \Xd)$ in~\cite{adr}.)

For $\lambda\in \Xd^+$, we denote by $\sI_*(\lambda)$ the 0-th perverse cohomology of the $*$-pushforward of $\omega_{\sL^+G z^\lambda \sL^+G/\sL^+G}[-\langle2\rho,\lambda\rangle]$ to $\Gr$. (Here $\sI_*(\lambda)$ corresponds to the induced $G^\vee_\k$-module $\nabla(\lambda)$ of highest weight $\lambda$ under the geometric Satake equivalence~\eqref{eqn:geom-Satake}.)
For $\lambda\in \Xd^0_J\cap \Xd^+_{J^\comp}$ 
we set 
\[
\sM^J_\lambda:= \Ras^{\I\mid \I^\quar_J}\big( \sW_{t_{-\lambda}} \star^\I \oblv^{\sL^+G\mid \I}(\sI_*(\lambda)) \big) 
\quad \in\ \Shv\left( \I^\quar_J \backslash \Gr \right). 
\]
If $\lambda,\mu \in \Xd^0_J\cap \Xd^+_{J^\comp}$ satisfy $\mu \unlhd^J \lambda$ the construction explained in~\cite[\S 8.1]{adr} provides a canonical morphism 
\[
\theta^J_{\mu,\lambda}: \sM^J_\mu \rightarrow \sM^J_\lambda,
\]
and this construction provides a functor from $(\Xd^0_J\cap \Xd^+_{J^\comp},\unlhd^J)$ to $\Shv\left( \I^\quar_J \backslash \Gr \right)$. (To justify this one can either use the corresponding fact when $J=\varnothing$ proved in~\cite[\S 8.1]{adr} and apply the equivalence $\Ras^{\I^\semi | \I_J^\quar}$, or repeat the same argument as in~\cite{adr} using Lemma~\ref{lem:MlambdaJ-perverse} below.)

We define the \textit{$J$-Gaitsgory sheaf}
\[
\Gait^{J,\semi}:= \colim_{\mu\in \Xd^0_J\cap \Xd^+_{J^\comp}} \sM^J_\mu. 
\]
Note that for any $\lambda\in \Xd^0_J\cap \Xd^+_{J^\comp}$ we have equivalences 
\begin{equation}
\label{equ:Mlambda-Jsemi}
\begin{aligned}
\sW_{t_{-\lambda}} \star^\I \oblv^{\sL^+G\mid \I}(\sI_*(\lambda))
&\simeq \nabla_{J,t_{-\lambda}} \star^{\I_J} \oblv^{\sL^+G\mid \I_J}(\sI_*(\lambda)) \\ 
&\simeq \oblv^{\I_J \mid \I}( \nabla^J_{J,t_{-\lambda}} \star^{\I_J} \oblv^{\sL^+G\mid \I_J}(\sI_*(\lambda))),
\end{aligned}
\end{equation}
where the first equivalence follows from the equality $(p_J)_* \nabla_{t_{-\lambda}} \simeq \nabla_{J,t_{-\lambda}}$, and the second one from the fact that $\I z^{-\lambda}\I_J/\I_J=\I_J z^{-\lambda}\I_J/\I_J$. 
These equivalences are compatible with the morphisms $\theta_{\mu,\lambda}^J$,
hence $\sM^J_\lambda$ and $\Gait^{J,\semi}$ are $\I^\semi_J$-equivariant, and by abuse of notation we view 
these objects as belonging to $\Shv\left( \I^\semi_J \backslash \Gr \right)$.

\begin{rmk}
\phantomsection
\label{rmk:Gait-functions-P/L}
\begin{enumerate}
\item
\label{it:OP/L-colim}
To understand the definition of $\Gait^{J,\semi}$, one should compare it with the fact that if $P^\vee_{J,\k} \subset G^\vee_\k$ is the (positive) standard parabolic subgroup associated with $J$ and $M^\vee_{J,\k}$ its Levi subgroup containing the canonical maximal torus $T^\vee_\k$ then we have
\[
\mathscr{O}(P^\vee_{J,\k} / M^\vee_{J,\k}) \simeq \colim_{\lambda \in \Xd^0_J\cap \Xd^+_{J^\comp}} \nabla(\lambda) \otimes \k_{P^\vee_{J,\k}}(-\lambda)
\]
as $P^\vee_{J,\k}$-modules, for the natural transition maps in the colimit. In fact one should consider $\Gait^{J,\semi}$ as a geometric incarnation of the $P^\vee_{J,\k}$-module $\mathscr{O}(P^\vee_{J,\k} / M^\vee_{J,\k})$. This idea will be made precise in~\cite{adrs}.
\item
The definition of $\Gait^{J,\semi}$ given above matches the description given (for cha\-racteristic-$0$ coefficients) in~\cite[\S 4.1.9]{DL23}. See~\S\ref{ss:Gaitsgory-alternative} for a description which matches the definition chosen in~\cite{DL23}.
\item
We expect that the object $\Gait^{J,\semi}$ is related to the intersection cohomology complex of Drinfeld's compactification of the stack $\mathrm{Bun}_{P_J}$ in the same way as in the characteristic-$0$ setting treated in~\cite[\S 4.2]{DL23} or in the case $J=\varnothing$ treated in~\cite{adr-Drinfeld}. This will be discussed in future work.
\end{enumerate} 
\end{rmk}

%---------------------------------------------------------------
\subsection{Main properties}
%---------------------------------------------------------------

\begin{lem}
\label{lem:MlambdaJ-perverse}
For any $\lambda\in \Xd^0_J\cap \Xd^+_{J^\comp}$,
the object $\sM^J_\lambda$ is perverse in $\Shv\left( \I^\quar_J \backslash \Gr \right)$. 
Therefore, we have
\[
\Gait^{J,\semi}\in \Shv\left( \I^\semi_J \backslash \Gr \right)^\heartsuit . 
\]
\end{lem}

\begin{proof}
Using Proposition~\ref{prop:Ras-D}\eqref{it:Ras-D-1} and then Lemma~\ref{lem:push-pull-pK-D-N-semi}\eqref{it:push-pull-pK-D-N-semi-1}
we have equivalences 
\begin{align*}
\Ras^{\I\mid \I^\quar_J}\big( \sW_{t_{-\lambda}} \star^\I \oblv^{\sL^+G\mid \I}(\sI_*(\lambda)) \big)
&\simeq \Delta^{J,\quar}_{t_{-\lambda}} \star^\I \oblv^{\sL^+G\mid \I}(\sI_*(\lambda)) \\ 
&\simeq \Delta^{J,\quar}_{S,t_{-\lambda}} \star^{\sL^+G} \sI_*(\lambda),
\end{align*}
and by Lemma~\ref{lem:standard-perverse} and Proposition~\ref{prop:right-conv-t-exact} this object belongs to $\Shv\left( \I^\quar_J \backslash \Gr \right)^\heartsuit$. The second assertion follows, using compatibility of our t-structures with colimits. (Note that by Lemma~\ref{lem:oblv-semi-quar-texact}, the questions of whether our objects are perverse when considered as $\I_J^\quar$-equivariant or as $\I_J^\semi$-complexes are equivalent.)
\end{proof}

As in~\cite[\S 6.7]{adr} we consider
the $q$-analogue of Kostant's partition function $\sP(\nu,q)\in \Z[q]$, for any $\nu\in \Xd$. 
%For any $\nu\in \Xd$, we denote by $\bs^J_\nu$ the embedding of the orbit $\I^\quar_J z^\nu \sL^+G/\sL^+G$ in $\Gr$. (So, we have $\bs^J_\nu = \bs^{\varnothing \subset J}_{S,x_\nu}$ where $x_\nu$ is the unique element in $t_\nu W \cap W_{\ex}^{(S,J)}$.)
%Recall that we have bijections 
%\[
%W_\ex^{(S,J)} \leftrightarrow W_\ex/W \leftrightarrow \Xd. 
%\]
%Note that on $\Gr$ we have $\rS^{J}_{S,e}=\rS^{J\subset J}_{S,e}$. 
For $\nu\in \Xd$, any $\I^{\quar}_J$-equivariant perverse sheaf $\sF$ on $\mathrm{S}^J_\nu$ is a direct sum of copies of the appropriate shift of the dualizing complex; in case the multiplicity is finite we will denote it $\rk \, \sF$. Finally, we will denote by ${\check \rho} \in \frac{1}{2}\Xd$ the halfsum of the positive coroots.

\begin{lem}
\phantomsection
\label{lem:costk-Mlambda}
%Let $\lambda\in \Xd^0_J\cap \Xd^+_{J^\comp}$. 
\begin{enumerate}
\item
\label{it:costk-Mlambda-2}
If $\lambda,\mu\in \Xd^0_J\cap \Xd^+_{J^\comp}$ satisfy $\mu \unlhd^J \lambda$, then the morphism $\theta^J_{\mu,\lambda}$ is injective. 
\item
\label{it:costk-Mlambda-1}
The objects $\sM^J_\lambda$ (for $\lambda\in \Xd^0_J\cap \Xd^+_{J^\comp}$) and $\Gait^{J,\semi}$ are supported on $\overline{\mathrm{S}^{J}_{0}}$. 
\item
\label{it:costk-Mlambda-3}
 If $\lambda\in \Xd^0_J\cap \Xd^+_{J^\comp}$ and if $\nu\in \Xd^+_J$ satisfies $\nu+\lambda\in \Xd^+$, then we have 
\[
\rk \, {}^{\mathrm{p}}\sH^0 \big( (\sigma^J_\nu)^!\sM^J_\lambda \big) = \begin{cases}
1 & \text{if $\nu=0$;} \\
0 & \text{otherwise},
\end{cases}
\quad
{}^{\mathrm{p}}\sH^1 \big( (\sigma^J_\nu)^!\sM^J_\lambda \big) = 0.
\]
\item
\label{it:costk-Mlambda-4} 
Assume that $\cha(\k)$ is good for $G$. 
If $\lambda\in \Xd^0_J\cap \Xd^+_{J^\comp}$ and if $\nu\in \Xd^+_J$ satisfies $\nu+\lambda\in \Xd^+$, then we have 
\[
\sum_{n \in \Z} \rk \, {}^{\mathrm{p}}\sH^n \big( (\sigma^J_\nu)^!\sM^J_\lambda \big) \cdot q^n 
= \sum_{w\in W} (-1)^{\ell(w)} \sP(w(\lambda+{\check \rho})- (\lambda+\nu+{\check \rho}), q^2). 
\]
\end{enumerate}
\end{lem}

\begin{proof}
\eqref{it:costk-Mlambda-2}
The image of $\theta^J_{\mu,\lambda}$ under the t-exact and conservative functor $\mathsf{F}_{\sL^+G}$ (see \textsection\ref{subsect:redtn-IW} and Theorem~\ref{thm:texactness}) is $\theta^\varnothing_{\mu,\lambda}$, which is injective by~\cite[Lemma~8.1]{adr}. Hence $\theta^J_{\mu,\lambda}$ is also injective.

\eqref{it:costk-Mlambda-1}
By~\eqref{it:costk-Mlambda-2} it suffices to prove that $\Gait^{J,\semi}$ is supported on $\overline{\mathrm{S}^{J}_{0}}$.
For $\lambda\in \Xd^0_J\cap \Xd^+_{J^\comp}$,
as in~\cite{adr} we will denote by $\overline{\nabla}_\lambda$ the costandard perverse sheaf on $\Gr$ associated with the $\I$-orbit of the image of $z^\lambda$, and
set $\sW^{\vee,\lambda}_S:=\sW_{t_{-\lambda}}\star^\I \overline{\nabla}_\lambda$. 
Since the restriction of $\sI_*(\lambda)$ to this $\I$-orbit is the perversely shifted dualizing sheaf, we have a natural morphism $\sI_*(\lambda)\rightarrow \overline{\nabla}_\lambda$, which induces a morphism 
\begin{equation}
\label{equ:Mlambda-pH0-dwaki-1}
\sM^J_\lambda \rightarrow {}^{\mathrm{p}}\sH^0\circ \Ras^{\I\mid \I^\quar_J}(\sW^{\vee,\lambda}_S) . 
\end{equation}
%By Corollary~\ref{cor:Ras-right-texact} the functor $\Ras^{\I\mid \I^\quar_J}$ is left t-exact, 
%\begin{claim}\label{claim:Mlambda-pH0-dwaki-inj}
%The morphism \eqref{equ:Mlambda-pH0-dwaki-1} is an injection. 
%\end{claim}
%\begin{proof}[Proof of Claim~\ref{claim:Mlambda-pH0-dwaki-inj}] 

We claim that the morphism~\eqref{equ:Mlambda-pH0-dwaki-1} is injective. In fact, $\overline{\nabla}_\lambda$ is $\I_J$-equivariant since 
%$\nabla_{S,t_\lambda}=\nabla^J_{S,t_\lambda}$ 
$\I z^{\lambda} \sL^+G/\sL^+G=\I_J z^{\lambda} \sL^+G/\sL^+G$. 
By a similar computation as in~\eqref{equ:Mlambda-Jsemi}, the right-hand side of~\eqref{equ:Mlambda-pH0-dwaki-1} is $\I^\semi_J$-equivariant, and~\eqref{equ:Mlambda-pH0-dwaki-1} is a morphism in $\Shv\left( \I^\semi_J \backslash \Gr \right)$. 
Applying the t-exact functor $\mathsf{F}_{\sL^+G}$ (see \textsection\ref{subsect:redtn-IW} and Theorem~\ref{thm:texactness}) to \eqref{equ:Mlambda-pH0-dwaki-1} and using Lemma~\ref{lem:Ras-shift-Wak}, we obtain a morphism
\begin{equation}
\label{equ:Mlambda-pH0-dwaki-2}
\sM^{\varnothing}_\lambda \rightarrow {}^{\mathrm{p}}\sH^0\circ \Ras^{\I\mid \I^\semi}(\sW^{\vee,\lambda}_S) 
\simeq {}^{\mathrm{p}}\sH^0\circ \Ras^{\I\mid \I^\semi}(\overline{\nabla}_\lambda)\langle -\lambda \rangle
\end{equation}
in $\Shv\left( \I^\semi \backslash \Gr \right)^\heartsuit$.
By \cite[Corollary 6.22]{adr}, the socle of $\sM^{\varnothing}_\lambda$ is the standard $\I^\semi$-equivariant perverse sheaf $\overline{\Delta}^\semi_{0}$ associated with the orbit of the base point. 
The precomposition of \eqref{equ:Mlambda-pH0-dwaki-2} with the embedding
\[
\Ras^{\I\mid \I^\semi}(\overline{\Delta}_{\lambda})\langle -\lambda \rangle \simeq \overline{\Delta}^\semi_{0} \hookrightarrow \sM^{\varnothing}_\lambda
\]
(where $\overline{\Delta}_{\lambda}$ is the standard perverse sheaf associated with the $\I$-orbit of the image of $z^\lambda$, and the equivalence follows from~\cite[Proposition~5.6]{adr})
coincides with the morphism induced by the natural map $\overline{\Delta}_{\lambda} \rightarrow \overline{\nabla}_{\lambda}$, and is therefore nonzero. (Here we use the fact that $\Ras^{\I | \I^\semi}$ is left t-exact by Corollary~\ref{cor:Ras-right-texact}, so that $\Ras^{\I\mid \I^\semi}(\overline{\nabla}_\lambda)$ belongs to $\Shv ( \I^\semi \backslash \Gr )^{\geq 0}$.)
Hence \eqref{equ:Mlambda-pH0-dwaki-2} is an injection, which implies that~\eqref{equ:Mlambda-pH0-dwaki-1} is injective as $\mathsf{F}_{\sL^+G}$ is t-exact and conservative. 
%\end{proof}
%\noindent 

By~\eqref{equ:costd-K-colim} we have 
\[
\nabla^{J,\quar}_{S,e}\simeq \colim_{\mu\in \Xd^0_J\cap \Xd^+_{J^\comp}} \Ras^{\I\mid \I^\quar_J}(\sW^{\vee,\mu}_S). 
\]
(Note that $e \in W_\ex^{(S,J)}$, so that the left-hand side makes sense.)
By our claim above, the colimit of the maps~\eqref{equ:Mlambda-pH0-dwaki-1} gives an injection 
\[
\Gait^{J,\semi} \hookrightarrow {}^{\mathrm{p}}\sH^0(\nabla^{J,\quar}_{S,e}), 
\]
from which~\eqref{it:costk-Mlambda-1} follows. 

\eqref{it:costk-Mlambda-3}--\eqref{it:costk-Mlambda-4}
The proofs are similar to those of the corresponding statements in~\cite[Lemma 6.18]{adr}; we only sketch the argument. The rank of ${}^{\mathrm{p}}\sH^n \big( (\sigma^J_\nu)^!\sM^J_\lambda \big)$ can be computed as the dimension of the morphism space in $\Shv(\I_{J,\mathrm{u}}^\quar \backslash \Gr)$ (where $\I_{J,\mathrm{u}}^\quar$ is the semidirect product of the pro-unipotent radical of $\I_{\varnothing \subset J}$ with $\sL U_J$) between the images of the standard $\I_J^\quar$-equivariant complex $\overline{\Delta}^{J,\quar}_{\nu}$ associated with $\mathrm{S}^J_\nu$ and $\sM^J_\lambda[n]$. Applying a pro-unipotent version of the Raskin equivalence and Lemma~\ref{lem:Ras-shift-Wak}, this space identifies with the space of morphisms in $\Shv(\Iu \backslash \Gr)$ between the images of $\Ras^{\I^\quar_J | \I}(\overline{\Delta}^{J,\quar}_{\lambda+\nu})$ and $\sI_*(\lambda)[n]$. (Here, $\Iu$ is the pro-unipotent radical of $\I$.) Since $\lambda+\nu$ is dominant, $t_{\lambda+\nu}$ belongs to $W_\ex^{(S,J)}$, and by Lemma~\ref{lem:push-pull-pK-D-N-semi}\eqref{it:push-pull-pK-D-N-semi-1} we have $\overline{\Delta}^{J,\quar}_{\lambda+\nu}=(p_S)_* \Delta^{J,\quar}_{t_{\lambda+\nu}}$. Using Proposition~\ref{prop:Ras-D}\eqref{it:Ras-D-1} we deduce that $\Ras^{\I^\quar_J | \I}(\overline{\Delta}^{J,\quar}_{\lambda+\nu})$ is the object denoted $\overline{\sW}_{\lambda+\nu}$ in~\cite{adr}. With this identification, the required dimensions are computed in~\cite{adr} and given as in the statement.
\end{proof}

\begin{rmk}
The restriction on the characteristic in Lemma~\ref{lem:costk-Mlambda}\eqref{it:costk-Mlambda-4} comes from the use of the Mirkovi{\'c}--Vilonen conjecture on costalks of the objects $\sI_*(\lambda)$, which is known to be true for all $\lambda \in \Xd^+$ if and only if $\mathrm{char}(\k)$ is good for $G$ (see~\cite[Proof of Lemma~6.18]{adr} for references). In fact, here we only need this property when $\lambda \in \Xd^0_J \cap \Xd_{J^\comp}^+$, so we might hope to relax this assumption. For instance, if $J=S$ it is easily seen that we do not need any restriction on $\mathrm{char}(\k)$. But we do not know how to implement this idea when $J \neq S$.
\end{rmk}

In the following proposition we consider the object $\nabla^{J,\semi}_{S,e}$. Note first that $e \in {}^J W_\ex^{(S,J)}$, so that this object makes sense. Note also that $\I_J^\semi \sL^+G = \I_J^\quar \sL^+ G$; so this object coincides with the object $\nabla^{J,\quar}_{S,e}$ considered in the proof of Lemma~\ref{lem:costk-Mlambda}.

\begin{prop}
\label{prop:costk-Ga}
\begin{enumerate}
\item
\label{it:costk-Ga-1} 
We have an isomorphism $\Gait^{J,\semi} \simeq {}^{\mathrm{p}}\sH^0(\nabla^{J,\semi}_{S,e})$. 
\item
\label{it:costk-Ga-2} 
Assume $\cha(\k)$ is good for $G$. 
For any $\nu\in \Xd^+_J$, we have 
\begin{equation*}
%\label{equ:costk-Ga}
\sum_{n} \rk\, {}^{\mathrm{p}}\sH^n\big( (\sigma^J_\nu)^!\Gait^{J,\semi} \big) \cdot q^n 
= \sum_{w\in W_J} (-1)^{\ell(w)} \sP(w({\check \rho})-{\check \rho} -\nu, q^2). 
\end{equation*}
\end{enumerate}
\end{prop}

\begin{proof}
%Note first that $e \in {}^J W_\ex^{(S,J)}$, so that the costandard object in~\eqref{it:costk-Ga-1} makes sense.
Here again the proof is similar to that of the counterpart of this statement in~\cite[\S 8.2]{adr}, so we will be brief. 
The claims follow from the computation of the rank of perverse cohomology objects of $!$-restrictions to strata of $\Gait^{J,\semi}$. In fact, since this object is $\I_J^\semi$-equivariant, and since $\mathrm{S}^J_\nu$ is open in the $\I_J^\semi$-orbit containing it when $\nu \in \Xd^+_J$, it suffices to consider the $!$-restrictions $(\sigma^J_\nu)^!\Gait^{J,\semi}$ when $\nu \in \Xd^+_J$. We compute the rank in this case 
as in the proof of Lemma~\ref{lem:costk-Mlambda}\eqref{it:costk-Mlambda-3}--\eqref{it:costk-Mlambda-4}, as the dimension of a certain morphism space in $\Shv(\I_{J,\mathrm{u}}^\quar \backslash \Gr)$. Since $\Gait^{J,\semi}$ is defined as a colimit and since the images of costandard objects in $\Shv(\I_{J,\mathrm{u}}^\quar \backslash \Gr)$ are compact, these spaces are colimits of their counterparts in $\sM_\lambda^J$, and to conclude we need to show that for fixed $\nu$, the latter spaces are independent of $\lambda$ for $\lambda$ ``large enough.'' 

This fact can be reduced to a statement which follows from the same considerations as in the proof of~\cite[Theorem~8.2]{adr}. Namely, using a ``pro-unipotent'' analogue of $\Ras^{\I_J^\quar | \I^\semi}$, it suffices to show that for any fixed $\sF$ in $\Shv_\omega(\I_{\mathrm{u}}^\semi \backslash \Gr)$ the space of morphisms from $\sF$ to the image of $\sM_\lambda^\varnothing$ is independent of $\lambda$ for $\lambda$ ``large enough.'' (See~\cite[\S 6.4]{adr} for the notation used here.) Any such $\sF$ is a finite extension of cohomological shifts of images of standard perverse sheaves, so one can assume that $\sF$ is such an object. Then one proceeds as in the proof of~\cite[Theorem~8.2]{adr}.
(Compared to that proof, in the present setting, all the elements $w \in W_J$ will contribute to the multiplicities, which also explains the sum over $W_J$ appearing here compared to the setting of~\cite{adr}.)
\end{proof}

% and 
%We remark that for any $\sF\in \Shv_\omega\left( \I^\quar_J \backslash \Gr \right)$, there exists $\lambda\in \Xd^0_J\cap \Xd^+_{J^\comp}$ such that for any $\mu,\mu' \in \Xd_J^0 \cap \Xd^+_{J^\comp}$ such that $\lambda \unlhd^J \mu \unlhd^J \mu'$, the morphism induced by $\theta_{\mu,\mu'}^J$ relating the vector space of morphisms from the image of $\sF$ to the image of $\sM^J_\mu$ and its counterpart for $\sM^J_{\mu'}$ is an isomorphism. In fact, applying a pro-unipotent version of $\Ras^{\I^\quar_J\mid \I^\semi}$ we reduce this claim to a similar claim for morphisms in $\Shv(\I^\semi_{\mathrm{u}} \backslash \Gr)$. Here $\Ras^{\I^\quar_J\mid \I^\semi}(\sF)$ is a finite extension of cohomological shifts of standard objects, so it suffices to prove a similar claim for each standard object, which follows from the same considerations as in the proof of~\cite[Theorem~8.2]{adr}. (Compared to that proof, in the present setting, all the elements $w \in W_J$ will contribute to the multiplicities.)
%
%\[
%\theta_{\mu,\mu'}^{J,\sF}: \Hom(\sF,\sM^J_\mu) \rightarrow \Hom(\sF,\sM^J_{\mu'}) 
%\]
%induced by $\theta_{\mu,\mu'}^J$ is an isomorphism.
%For any $\mu,\mu' \in \Xd_J^0 \cap \Xd^+_{J^\comp}$,
%applying the equivalence $\Ras^{\I^\quar_J\mid \I^\semi}$, $\theta_{\mu,\mu'}^{J,\sF}$ identifies with the morphism
%\[
%%\Ras^{\I^\quar_J\mid \I^\semi}(\theta_{\mu,\mu'}^{\sF}):\ 
%\Hom(\Ras^{\I^\quar_J\mid \I^\semi}(\sF),\sM^\emptyset_\mu) \rightarrow 
%\Hom(\Ras^{\I^\quar_J\mid \I^\semi}(\sF),\sM^\emptyset_{\mu'})
%\]
%induced by $\theta_{\mu,\mu'}^\varnothing$.
%Since $\Ras^{\I^\quar_J\mid \I^\semi}(\sF) \in \Shv_\omega\left( \I^\semi \backslash \Gr \right)$, this object is a finite extension of cohomological shifts of standard objects, 
%
%%it supports on finitely many $\I^\semi$-orbits by Lemma~\ref{lem:constr-complx}. 
%%We denote by $V\subset \Gr$ the union of such orbits. 
%The arguments in the proof of \cite[Theorem 8.2]{adr} show that,    there exists $\lambda\in \Xd^0_J\cap \Xd^+_{J^\comp}$, such that for any $\lambda \unlhd^J \mu \unlhd^J \mu'$ the $!$-restriction of the morphism $\sM^\emptyset_\mu\rightarrow \sM^\emptyset_{\mu'}$ on $V$ is an equivalence. 
%Hence $\Ras^{\I^\quar_J\mid \I^\semi}(\theta_{\mu,\mu'}^{\sF})$, and thus $\theta_{\mu,\mu'}^{\sF}$, is an isomorphism. 
%
%Now item \eqref{it:costk-Ga-2} follows from Lemma~\ref{lem:costk-Mlambda}\eqref{it:costk-Mlambda-3}. 
%The item \eqref{it:costk-Ga-1} can be proved similarly as in \cite[Proposition 8.4(1)]{adr}, provided Lemma~\ref{lem:costk-Mlambda}\eqref{it:costk-Mlambda-1} and \eqref{equ:costk-Ga}. 
%\end{proof}

\begin{rmk}
Comparing the arguments above with the formula in Remark~\ref{rmk:Gait-functions-P/L}\eqref{it:OP/L-colim}, one sees that for $\nu \in \Xd_J^+$ the integer $\sum_{w\in W_J} (-1)^{\ell(w)} \sP(w({\check \rho})-{\check \rho} -\nu, 1)$ is the dimension of the $T^\vee_\k$-weight space of weight $\nu$ in $\mathscr{O}(P^\vee_{J,\k} / M^\vee_{J,\k})$.
%Note that for any $\nu\in \Xd$ we have 
%\[
%\sum_{w\in W_J} (-1)^{\ell(w)} \sP(w(\rho)-\rho -\nu, 1)=\dim_\k\sO(U^{\vee,-}_J)_\nu. 
%\]
%The evaluation of the RHS of \eqref{equ:costk-Ga} at $q=1$ is equal to $\dim_\k\sO(U^{\vee,-}_J)_\nu$. 
\end{rmk}

%-----------------------------------------------------------
\subsection{Alternative description}
\label{ss:Gaitsgory-alternative}
%-----------------------------------------------------------

As in~\cite[\S 8.3]{adr} one can describe the object $\Gait^{J,\semi}$ as a different colimit. Namely, generalizing slightly the notation introduced in \S\ref{ss:convolution}, for $\mu \in \Xd$ we denote by $z^\mu \cdot (-)$ the pushforward functor under the automorphism of $\Gr$ given by the left action of $z^\mu \in \sL G(\F)$. It is explained in~\cite[\S 8.3]{adr} (following Gaitsgory) that the assignment $\lambda \mapsto z^{-\lambda} \cdot \sI_*(\lambda)[\langle \lambda,2\rho \rangle]$ defines a functor from $\Xd^+$ (endowed with the restriction of the preorder $\unlhd^\varnothing$) to $\Shv(T \backslash \Gr)$. Here we restrict this functor to $\Xd_J^0 \cap \Xd^+_{J^\comp}$. (Note that the restriction of $\unlhd^\varnothing$ to this subset coincides with the restriction of $\unlhd^J$.)

The following lemma follows from the same considerations as for~\cite[Proposition~8.6]{adr}.

\begin{lem}
\label{lem:Gait-alternative}
There exists a canonical equivalence
\[
\Gait^{J,\semi} \simeq \colim_{\lambda \in \Xd_J^0 \cap \Xd^+_{J^\comp}} z^{-\lambda} \cdot \sI_*(\lambda)[\langle \lambda,2\rho \rangle].
\]
\end{lem}

Consider the affine Grassmannians $\Gr_{M_J}$ and $\Gr_{P_J}$ of $M_J$ and $P_J$, and the diagram
\[
\begin{tikzcd}
\Gr_{M_J} \ar[r, "b"',shift right=0.5ex] & \Gr_{P_J} \ar[r, "c"] \ar[l, "a"', shift right=0.5ex] & \Gr
\end{tikzcd}
\]
where the morphisms are induced by the natural maps $P_J \to G$, $P_J \to M_J$ and $M_J \to P_J$. (This setting is similar to that considered in the proof of Proposition~\ref{prop:exactness-left-conv}, but for $\Gr$ instead of $\Fl$, so we use similar notation.) The connected components in $\Gr_{M_J}$ and $\Gr_{P_J}$ are naturally parametrized by $\Xd / \Z {\check \Phi}_J$; for $\zeta \in \Xd / \Z {\check \Phi}_J$ we will denote by $\Gr_{M_J, \zeta}$ and $\Gr_{P_J,\zeta}$ the corresponding components, and by $a_\zeta$, $b_\zeta$, $c_\zeta$ the restrictions of the maps above to these components. As in~\cite[Lemma~3.3.2]{DL23}, the functors $(a_\zeta)^!$ and $(b_\zeta)^!$ induce mutually inverse equivalences
\begin{equation}
\label{eqn:equiv-Gr-P-M}
\Shv \left( \I_J^\semi \backslash \Gr_{P_J, \zeta} \right) \simeq \Shv \left( \sL^+ M_J \backslash \Gr_{M_J, \zeta} \right).
\end{equation}

Let us now denote by $\Sat$ the equivalence~\eqref{eqn:geom-Satake}, and by $\Sat_{J}$ its counterpart of $M_J$. It is well known that for any $\zeta \in \Xd / \Z {\check \Phi}_J$ the functor
\[
(a_\zeta)_! (c_\zeta)^* [\langle \zeta, 2\rho-2\rho_J \rangle] : \Shvfg \left( \sL^+ G \backslash \Gr \right) \to \Shvfg \left( \sL^+ M_J \backslash \Gr_{M_J,\zeta} \right)
\]
is t-exact, and corresponds under the equivalences $\Sat$ and $\Sat_J$ to the composition of the functor $(-)_{M^\vee_{J,\k}}$ of restriction of the action of $G^\vee_\k$ to $M^\vee_{J,\k}$ with the functor $(-)_{\mathrm{Z}_J=\zeta}$ of projection to the direct summand where the center of $M_{J,\k}^\vee$ acts via the character defined by $\zeta$.

The following statement generalizes~\cite[Proposition~4.1.12]{DL23} and~\cite[Remark~8.3]{adr}. Here, for $\mu \in \Xd^0_J$ we denote by $\k_{M^\vee_{J,\k}}(\mu)$ the $1$-dimensional $M^\vee_{J,\k}$-module determined by $\mu$.

\begin{prop}
For any $\zeta \in \Xd / \Z {\check \Phi}_J$ we have an equivalence
\[
(c_\zeta)^* \Gait^{J,\semi} \simeq (a_\zeta)^! \left( \colim_{\lambda \in \Xd_J^0 \cap \Xd^+_{J^\comp}} \Sat_J \left( (\k_{M^\vee_{J,\k}}(-\lambda) \otimes \nabla(\lambda)_{| M^\vee_{J,\k}})_{\mathrm{Z}_J=\zeta} \right) \right)[- \langle \zeta, 2\rho-2\rho_J \rangle]
\]
where the maps in the colimit are induced by those considered in Remark~\ref{rmk:Gait-functions-P/L}\eqref{it:OP/L-colim}.
\end{prop}

\begin{proof}
The proof is identical to that of~\cite[Proposition~4.1.12]{DL23}. Namely, using the equivalence~\eqref{eqn:equiv-Gr-P-M}, it suffices to prove that
\[
(b_\zeta)^! (c_\zeta)^* \Gait^{J,\semi} \simeq \colim_{\lambda \in \Xd_J^0 \cap \Xd^+_{J^\comp}} \Sat_J \left( (\k_{M^\vee_{J,\k}}(-\lambda) \otimes \nabla(\lambda)_{| M^\vee_{J,\k}})_{\mathrm{Z}_J=\zeta} \right) [- \langle \zeta, 2\rho-2\rho_J \rangle].
\]
Now, using Lemma~\ref{lem:Gait-alternative} we obtain that
\[
(b_\zeta)^! (c_\zeta)^* \Gait^{J,\semi} \simeq \colim_{\lambda \in \Xd_J^0 \cap \Xd^+_{J^\comp}} z^{-\lambda} \cdot (b_{\overline{\lambda}+\zeta})_! (c_{\overline{\lambda}+\zeta})^* \sI_*(\lambda)[\langle \lambda,2\rho \rangle]
\]
where $\overline{\lambda}$ is the image of $\lambda$ in $\Xd / \Z {\check \Phi}_J$, and we still use the notation $z^{-\lambda} \cdot (-)$ for the counterpart of this functor for sheaves on $\Gr_{M_J}$. By the comments above we have
\[
(b_{\overline{\lambda}+\zeta})^! (c_{\overline{\lambda}+\zeta})^* \sI_*(\lambda) \simeq \Sat_J \left( (\nabla(\lambda)_{| M^\vee_{J,\k}})_{\mathrm{Z}_J=\overline{\lambda}+\zeta} \right) [-\langle \overline{\lambda}+\zeta, 2\rho-2\rho_J \rangle],
\]
and we deduce that
\begin{multline*}
z^{-\lambda} \cdot (b_{\overline{\lambda}+\zeta})^! (c_{\overline{\lambda}+\zeta})^* \sI_*(\lambda)[\langle \lambda,2\rho \rangle] \\
\simeq \Sat_J \left( (\k_{M^\vee_{J,\k}}(-\lambda) \otimes \nabla(\lambda)_{| M^\vee_{J,\k}})_{\mathrm{Z}_J=\zeta} \right) [- \langle \zeta, 2\rho-2\rho_J \rangle],
\end{multline*}
which finishes the proof.
\end{proof}

%%%%%%%%%%%%%%%%%%%%
\iffalse
%%%%%%%%%%%%%%%%%%%%%%%%%
\section{t-exactness for $L^+G$} 
%%%%%%%%%%%%%%%%%%%%%%%%%
Notations are mainly adopted from \cite{adr}. 
For subgroups $K,H\subset LG$, as long as the followings are well-defined, we denote by $\D_{H,K}=\mathsf{Shv}(H\backslash LG/K)$ the stable $\infty$-category of sheaves on $H\backslash LG/K$, and $\P_{H,K}=\mathsf{Shv}(H\backslash LG/K)^\heartsuit$ the heart of perverse t-structure. 
We denote by $\D^c_{H,K}$, $\P^c_{H,K}$ the full subcategories of constructible complexes. 

We also replace $K$ by $\IW$ in the notations above for the Iwahori--Whittaker sheaves (i.e. sheaves that are $I_u^-$-equivariant by a non-degenerate character $I_u^-\rightarrow \Ga$). 
We denote by $\Av: \D_{H,I_u}\rightarrow \D_{H,\IW}$ the averaging functor. 

\begin{prop}
The composition of functors 
\[
\mathsf{F}:\ \D_{L^+G,I_u}\xrightarrow{\For^{L^+G}_I} \D_{I,I_u} \xrightarrow{\mathrm{Ras}^{-1}} \D_{I^\semi,I_u}
\] 
is t-exact with respect to the perverse t-structure on both sides. 
\end{prop}
\begin{proof}
Note that $\mathsf{F}$ is left t-exact, since $\For^{L^+G}_I:\D_{L^+G,I_u}\rightarrow \D_{I,I_u}$ is t-exact and $\mathrm{Ras}^{-1}:\D_{I,I_u}\rightarrow \D_{I^\semi,I_u}$ is left t-exact. 
We need to show that $\mathsf{F}$ is right t-exact. 
It is enough to show that for any $\sF\in \P^c_{L^+G,I_u}$, we have $\mathsf{F}(\sF)\in \D^{\leq 0}_{I^\semi,I_u}$. 
By \cite[Lem 6.11]{adr}, it is equivalent to show that $\sW_{t_\lambda}\star^I \For^{L^+G}_I(\sF)\in \D^{\leq 0}_{I,I_u}$, for any $\lambda\in \Xd$. 

We consider the affine braid group $\mathfrak{B}_\af$, equipped with an action $\widehat{\star}$ on $\D_{I,I_u}$ (from the right), such that the simple reflection $s$ acts by convolution with the free monodromy costandard sheaf $\widehat{\nabla}_s\in \widehat{\D}_{I_u,I_u}$. 
We have the following lemma. 
\begin{lem}[{\cite[Prop 3]{BLin}}]\label{lem 1.2}
Let $\sG\in \D^c_{I,I_u}$. 
We have $\sG\in \D^{\leq 0}_{I,I_u}$ if and only if $\Av(\sG\ \widehat{\star}\ b)\in \D^{\leq 0}_{I,\IW}$ for any $b\in \mathfrak{B}^+_\af$ (positive element). 
\end{lem}

By Lemma~\ref{lem 1.2}, we need to show that $\sW_{t_\lambda}\star^I \Av( \For^{L^+G}_I(\sF\ \widehat{\star}\ b))\in \D^{\leq 0}_{I,\IW}$ for any $b\in \mathfrak{B}^+_\af$ and $\lambda\in \Xd$. 
Since $-\ \widehat{\star}\ b$ is right t-exact and $\Av$ is t-exact, it only needs to show that for any $\sF\in \P^c_{L^+G,I_u}$ and $\lambda \in \Xd$ we have 
\begin{equation}\label{equ 1.1}
\sW_{t_\lambda}\star^I (\Av \circ\For^{L^+G}_I(\sF))\in \D^{\leq 0}_{I,\IW}. 
\end{equation}
Consider $\Av(\sF)\in \P^c_{L^+G,\IW}$. 
By \cite[Thm 3.9]{BGMRR}, there is $\sF'\in \P^c_{L^+G,L^+G}$ such that $\Av(\sF)=\sF' \star^{L^+G} \overline{\nabla}^{\IW}_{\rho}$ where $\overline{\nabla}^{\IW}_{\rho}$ is the costandard sheaf in $\P^c_{L^+G,\IW}$ corresponding to the orbit $L^+G\backslash L^+G \cdot t^\rho I^-_u$. 
Hence we have 
\begin{align*}
\Av \circ\For^{L^+G}_I(\sF) &= \For^{L^+G}_I \circ \Av(\sF)\\
&= \For^{L^+G}_I( \sF' \star^{L^+G} \overline{\nabla}^{\IW}_{\rho}) \\
&= \mathsf{Z}(\sF') \star^{I} \For^{L^+G}_I(\overline{\nabla}^{\IW}_{\rho}), 
\end{align*}
where $\mathsf{Z}:\D_{L^+G,L^+G}\rightarrow \D_{I,I}$ is the central functor. 
Since $\mathsf{Z}(\sF')$ is composed by $\sW_{t_\mu}$ for some $\mu\in \Xd$ (see \cite[Thm 4.4.5]{ARbook}), \eqref{equ 1.1} is reduced to show that for any $\lambda \in \Xd$, 
\begin{equation}\label{equ 1.2}
\sW_{t_\lambda}\star^I \For^{L^+G}_I(\overline{\nabla}^{\IW}_{\rho})\in \D^{\leq 0}_{I,\IW}. 
\end{equation}
By a similar statement to \cite[Lem 4.4]{adr}, to show \eqref{equ 1.2} for any $\lambda\in \Xd$ is equivalent to show  \eqref{equ 1.2} for any dominant $\lambda$. 
Now $\sW_{t_\lambda}=\Delta_{t_\lambda}$ for $\lambda$ dominant, and standard argument shows that $\Delta_{t_\lambda}\star^I \For^{L^+G}_I(\overline{\nabla}^{\IW}_{\rho})$ is perverse. 
It completes the proof. 
\end{proof}

We include the proof of Lemma~\ref{lem 1.2} in the following. 
\begin{proof}[Proof of Lemma~\ref{lem 1.2}]
The ``only if" part is easy, since $-\ \widehat{\star}\ b$ is right t-exact and $\Av$ is t-exact. 
We show the ``if" part. 
Let $n$ be the maximal integer such that $\sG'={}^p\sH^n(\sG)\neq 0$, and let $\sL$ be a simple quotient of $\sG'$. 

\textit{Step 1.} 
We reduce the assertion to the following claim 
\begin{equation}
{}^p\sH^0(\Av(\sL\ \widehat{\star}\ b))\neq 0, \quad \text{for some } b\in \mathfrak{B}_\af^+. 
\end{equation}
Indeed, by the right t-exactness of $\mathfrak{B}_\af^+$-action and the t-exactness of $\Av$, we have ${}^p\sH^{>n}(\Av(\sG\ \widehat{\star}\ b))=0$ and 
\[
{}^p\sH^n(\Av(\sG\ \widehat{\star}\ b))={}^p\sH^0(\Av(\sG'\ \widehat{\star}\ b)) \twoheadrightarrow {}^p\sH^0(\Av(\sL\ \widehat{\star}\ b)). 
\]
%and ${}^p\sH^0(\Av(\sL\ \widehat{\star}\ b))$ is a quotient of ${}^p\sH^0(\Av(\sG'\ \widehat{\star}\ b))$. 
Hence the fact that $\Av(\sG\ \widehat{\star}\ b)\in \D^{\leq 0}_{I,\IW}$ forces $n\leq 0$, in other words, $\sG\in \D^{\leq 0}_{I,I_u}$. 

\textit{Step 2.} 
We claim that  
\begin{equation}\label{equ 1.4}
\Av(\sL\ \widehat{\star}\ b)\neq 0, \quad \text{for some } b\in \mathfrak{B}_\af^+. 
\end{equation}
We use the following fact (see e.g. \cite[Prop 4.4.4]{ARbook}): there exist a subset $A\subset W_\ex$ such that for any $x\in W_\ex$ and $y\in W_\af$, 
\[
i^!_x(\sL \ \widehat{\star}\ \widetilde{y}) \neq 0 \quad \Rightarrow \quad x\in A\cdot y,
\]
where $\widetilde{y}\in \mathfrak{B}^+_\af$ is the canonical lift of $y$ and $i_x:I\backslash I xI_u\rightarrow I\backslash LG$ is the embedding. 

Now let $\lambda$ be dominant enough (and in the coroot lattice) such that $A\cdot t_\lambda$ consists of elements of the form $wt_\mu$, with $w\in W$ and $\mu\in \Xd$ strictly dominant. 
Then for each right $W$-coset in $W_\ex$, there is at most one element contained in $A\cdot t_\lambda$. 
Let $x\in A\cdot t_\lambda$ such that $i^!_x(\sL \ \widehat{\star}\ \widetilde{t_\lambda}) \neq 0$, then we have 
\begin{align*}
\Hom^\bullet(\Delta^{\IW}_x, \Av(\sL \ \widehat{\star}\ \widetilde{t_\lambda}))
&= \Hom^\bullet(\Av^*(\Delta^{\IW}_x), \sL \ \widehat{\star}\ \widetilde{t_\lambda})\\
&= \Hom^\bullet(\Delta_x, \sL \ \widehat{\star}\ \widetilde{t_\lambda})\neq 0,
\end{align*}
where $\Delta^{\IW}_x\in \D_{I,\IW}$ is the standard sheaf corresponding to the orbit $I\backslash IxI_u^-$, $\Av^*$ is the right adjoint functor of $\Av$, and the second isomorphism is because that $\Av^*(\Delta^{\IW}_x)$ is composed by $\Delta_{xw}$ ($w\in W$, each of which appears once) and $i^!_{xw}(\sL \ \widehat{\star}\ \widetilde{t_\lambda}) = 0$ if $w\neq e$. 
We deduce that $\Av(\sL \ \widehat{\star}\ \widetilde{t_\lambda}) \neq 0$, and in particular \eqref{equ 1.4}. 

\textit{Step 3.} 
We let $b_0$ be an element of minimal length in $\{b\in \mathfrak{B}^+_\af \mid \Av(\sL\ \widehat{\star}\ b_0)\neq 0\}$, then we show that ${}^p\sH^0(\Av(\sL\ \widehat{\star}\ b_0))\neq 0$. 

Consider the long exact sequence for any simple reflection $s\in W_\ex$, 
\begin{equation}\label{equ 1.5}
0\rightarrow \Delta_e \rightarrow \Delta_s \rightarrow \nabla_s \rightarrow \Delta_e \rightarrow 0. 
\end{equation}
Since $\sL$ is $I$-equivariant (from the right), we have $\sL \ \widehat{\star}\ \widetilde{s}= \sL\ {\star^I} \nabla_s$ and $\sL \ \widehat{\star}\ \widetilde{s}^{-1}= \sL\ {\star^I} \Delta_s$. 
Hence using \eqref{equ 1.5} one can show that for any $b\in \mathfrak{B}_\af^+$, 
\[
\sL\ \widehat{\star}\ b \simeq \sL\ \widehat{\star}\ b^{-1}, \quad \text{mod}\quad \langle \sL\ \widehat{\star}\ b'\mid \ell(b')<\ell(b)\rangle.
\]
By our assumption on $b_0$, we have $\sL\ \widehat{\star}\ b_0 \simeq \sL\ \widehat{\star}\ b_0^{-1}$ and hence it is perverse (note that $-\ \widehat{\star}\ b_0$ is right t-exact and $-\ \widehat{\star}\ b_0^{-1}$ is left t-exact). 
Since $\Av$ is t-exact, we deduce that ${}^p\sH^0(\Av(\sL\ \widehat{\star}\ b_0))=\Av(\sL\ \widehat{\star}\ b_0)\neq 0$. 
\end{proof}
\fi
%%%%%%%%%%%%%%%%%%%%%

\bibliographystyle{plain} 
\bibliography{MyBibtex}

\end{document}

%%%%%%%%%%%%%%%%%%%%%%%%%
%\section{} 
%%%%%%%%%%%%%%%%%%%%%%%%%
%---------------------------------------------------------------
%\subsection{} 
%---------------------------------------------------------------
%===============================
%\subsubsection{} 
%===============================




